% Chap1: EMPHASIS ON LITERARY APPRECIATION — Literature in English Upper Sixth — Cours signé OnBuch+
% Official MINESEC syllabus. Compiler avec : tectonic LIT01-chap1-emphasis-on-literary-appreciation.tex
\newcommand{\DOCMATIERE}{Literature in English}
\newcommand{\DOCNIVEAU}{Upper Sixth}
\newcommand{\DOCMODULE}{Upper Sixth — Literature in English}
\newcommand{\DOCLECON}{1}
\newcommand{\DOCTITRE}{Chap1: EMPHASIS ON LITERARY APPRECIATION}
% OnBuch+ — préambule « cours signés » ENRICHI (compilé avec tectonic / XeLaTeX)
% Système visuel « Pop » d'OnBuch+ : fond crème (jamais blanc), bordures encre
% épaisses, ombres « dures » décalées, titres Archivo Black en capitales.
% Version MATHÉMATIQUES : graphes (pgfplots/tikz), boîtes pédagogiques
% (définition, propriété, méthode, attention, le savais-tu, exercice résolu,
% exercice, TP, bilan), page de garde et signature OnBuch+ ancrée en bas.
%
% Macros attendues dans main.tex AVANT \input{preamble} :
%   \DOCMATIERE  \DOCNIVEAU  \DOCMODULE  \DOCLECON (numéro)  \DOCTITRE
\documentclass[11pt]{article}

\usepackage{fontspec}
\usepackage[english]{babel}
\usepackage[a4paper,margin=2cm,top=3.4cm,bottom=2.8cm,headheight=1.4cm,footskip=1.3cm]{geometry}
\usepackage{amsmath,amssymb,amsfonts}
\usepackage{mathtools} % \xmapsto, \coloneqq... souvent utilisés par les modèles
\usepackage{cancel} % simplifications barrées (bilans de forces)
% \abs{z} (module, valeur absolue) et \norm{u} (norme), utilisés par les modèles
\providecommand{\abs}{}\let\abs\relax\DeclarePairedDelimiter{\abs}{\lvert}{\rvert}
\providecommand{\norm}{}\let\norm\relax\DeclarePairedDelimiter{\norm}{\lVert}{\rVert}
\usepackage[table]{xcolor} % [table] : fournit aussi \rowcolors (alternance de lignes)
\usepackage{pagecolor}
\usepackage{tikz}
\usetikzlibrary{shapes.symbols,circuits.logic.US,circuits.logic.IEC,arrows.meta,positioning,calc,shapes.geometric,shapes.misc,shapes.arrows,decorations.pathmorphing,decorations.pathreplacing,decorations.markings,patterns,shadows,mindmap,trees,fit,backgrounds,matrix,intersections,angles,quotes}
\usepackage{pgfplots}
\usepackage[version=4]{mhchem} % équations nucléaires \ce{^{235}_{92}U}
\usepackage[european,siunitx]{circuitikz} % schémas électriques
\pgfplotsset{compat=1.18}
\usepgfplotslibrary{fillbetween,groupplots} % aires entre courbes (calcul intégral)
\usepackage{enumitem}
\usepackage{fancyhdr}
% [most] charge toutes les bibliothèques tcolorbox (listings, minted, raster,
% documentation...) et gonfle inutilement la mémoire TeX sur un document long
% (~300 boîtes) : on ne charge que ce qui est réellement utilisé.
\usepackage[skins,breakable]{tcolorbox}
\usepackage{graphicx}
\usepackage{colortbl}
\usepackage{booktabs}
\usepackage{tabularx}
\usepackage{diagbox} % case d'en-tête diagonale des tableaux à double entrée
% Colonnes extensibles centrée / gauche / droite (C, L, R), souvent utilisées par les modèles
\newcolumntype{C}{>{\centering\arraybackslash}X}
\newcolumntype{L}{>{\raggedright\arraybackslash}X}
\newcolumntype{R}{>{\raggedleft\arraybackslash}X}
\usepackage{array}
\usepackage{multicol}
\usepackage{siunitx}
% parse-numbers=false : accepte aussi \SI{2,5 \times 10^{-3}}{...}
\sisetup{locale=UK,output-decimal-marker={.},group-minimum-digits=4,parse-numbers=false}
\DeclareSIUnit\ohm{\ensuremath{\symup{\Omega}}} % U+2126 absent de la police
\DeclareSIUnit\division{div} % réglages d'oscilloscope (ms/div, V/div)
\DeclareSIUnit\year{yr}\DeclareSIUnit\annum{yr}\DeclareSIUnit\Ma{Ma}\DeclareSIUnit\tr{tr} % tours (vitesse de rotation, ex. \SI{1500}{\tr\per\minute})
\usepackage{qrcode}
% \degree : commande fréquemment utilisée par les modèles pour un angle ou une
% température en dehors de \SI{}{} (non fournie par les paquets chargés).
\providecommand{\degree}{^{\circ}}
% \qed (fin de démonstration), utilisé par les modèles sans amsthm
\providecommand{\qed}{\hfill$\square$}
% \barre : double barre des valeurs interdites dans un tableau de variations (tkz-tab)
\providecommand{\barre}{\ensuremath{\|}}
% environnement proof (amsthm non chargé) utilisé par les modèles
\ifdefined\proof\else\newenvironment{proof}[1][Proof]{\par\noindent\textit{#1.}\ }{\qed\par}\fi
\usepackage{lastpage}
\usepackage{hyperref}
\hypersetup{hidelinks}

% ============================================================
% Palette « Pop » OnBuch+ — mêmes hex que la classe P (plus_ui.dart)
% ============================================================
\definecolor{poporange}{HTML}{F59321}
\definecolor{popdark}{HTML}{E07A0C}
\definecolor{popcream}{HTML}{FFF6E8}
\definecolor{popink}{HTML}{1C1714}
\definecolor{poppurple}{HTML}{7A5AE0}
\definecolor{popgreen}{HTML}{2E9E6A}
\definecolor{poppink}{HTML}{E0457B}
\definecolor{popblue}{HTML}{2F7DE1}
\definecolor{popgold}{HTML}{C98A12}
\definecolor{popmuted}{HTML}{8A7F76}
\definecolor{popline}{HTML}{EADFD2}
\definecolor{popgray}{HTML}{8A7F76}
\colorlet{popgrey}{popgray}
% Alias tolérants (le modèle nomme parfois les couleurs différemment)
\colorlet{popwhite}{white}
\colorlet{popblack}{popink}
\colorlet{popred}{poppink}
\colorlet{popyellow}{popgold}
% Teintes claires pour fonds de tableaux / zones
\colorlet{popblueL}{popblue!10!white}
\colorlet{poporangeL}{poporange!14!white}
\colorlet{popgreenL}{popgreen!12!white}
\colorlet{poppurpleL}{poppurple!12!white}
\colorlet{poppinkL}{poppink!12!white}
\colorlet{popgoldL}{popgold!14!white}

\pagecolor{popcream}

% ============================================================
% Typographie
% ============================================================
\defaultfontfeatures{Ligatures=TeX}
\setmainfont{PlusJakartaSans-Medium.ttf}[Path=../../../fonts/,BoldFont=PlusJakartaSans-Bold.ttf]
% Maths en sans-serif (Fira Math) avec chiffres et lettres droites de Plus
% Jakarta, pour que les formules se fondent dans le texte.
\usepackage[math-style=ISO,bold-style=ISO]{unicode-math}
\setmathfont{FiraMath-Regular.otf}[Path=../../../fonts/]
\setmathfont{PlusJakartaSans-Medium.ttf}[Path=../../../fonts/,range={up/num,up/{Latin,latin},\mathup}]
% (icomma retiré : décimales avec point en anglais)
% Caractères mathématiques Unicode absents des polices de texte (Plus Jakarta,
% Archivo Black) : rendus par leur équivalent mathématique, en texte comme en
% formule — sinon ils s'affichent en carré vide (ex. « Arithmétique dans ℤ »).
\usepackage{newunicodechar}
\newunicodechar{ℝ}{\ensuremath{\mathbb{R}}}
\newunicodechar{ℤ}{\ensuremath{\mathbb{Z}}}
\newunicodechar{ℂ}{\ensuremath{\mathbb{C}}}
\newunicodechar{ℕ}{\ensuremath{\mathbb{N}}}
\newunicodechar{ℚ}{\ensuremath{\mathbb{Q}}}
\newunicodechar{𝔻}{\ensuremath{\mathbb{D}}}
\newunicodechar{α}{\ensuremath{\alpha}}
\newunicodechar{β}{\ensuremath{\beta}}
\newunicodechar{γ}{\ensuremath{\gamma}}
\newunicodechar{δ}{\ensuremath{\delta}}
\newunicodechar{ε}{\ensuremath{\varepsilon}}
\newunicodechar{θ}{\ensuremath{\theta}}
\newunicodechar{λ}{\ensuremath{\lambda}}
\newunicodechar{μ}{\ensuremath{\mu}}
\newunicodechar{σ}{\ensuremath{\sigma}}
\newunicodechar{τ}{\ensuremath{\tau}}
\newunicodechar{φ}{\ensuremath{\varphi}}
\newunicodechar{ψ}{\ensuremath{\psi}}
\newunicodechar{ω}{\ensuremath{\omega}}
\newunicodechar{Σ}{\ensuremath{\Sigma}}
% majuscules produites par \MakeUppercase dans les titres (α-aminés -> Α)
\newunicodechar{Α}{\ensuremath{\alpha}}
\newunicodechar{Β}{\ensuremath{\beta}}
\newunicodechar{Γ}{\ensuremath{\Gamma}}
\newunicodechar{Δ}{\ensuremath{\Delta}}
\newunicodechar{Ε}{\ensuremath{\varepsilon}}
\newunicodechar{Θ}{\ensuremath{\Theta}}
\newunicodechar{Λ}{\ensuremath{\Lambda}}
\newunicodechar{Μ}{\ensuremath{\mu}}
\newunicodechar{Τ}{\ensuremath{\tau}}
\newunicodechar{Φ}{\ensuremath{\Phi}}
\newunicodechar{Ψ}{\ensuremath{\Psi}}
\newunicodechar{Ω}{\ensuremath{\Omega}}
\newunicodechar{↦}{\ensuremath{\mapsto}}
\newunicodechar{∈}{\ensuremath{\in}}
\newunicodechar{∉}{\ensuremath{\notin}}
\newunicodechar{∧}{\ensuremath{\wedge}}
\newunicodechar{⇔}{\ensuremath{\Leftrightarrow}}
\newunicodechar{⟺}{\ensuremath{\Longleftrightarrow}}
\newunicodechar{⇒}{\ensuremath{\Rightarrow}}
\newunicodechar{⟹}{\ensuremath{\Longrightarrow}}
\newunicodechar{∘}{\ensuremath{\circ}}
\newunicodechar{∪}{\ensuremath{\cup}}
\newunicodechar{∩}{\ensuremath{\cap}}
\newunicodechar{≡}{\ensuremath{\equiv}}
\newunicodechar{⊂}{\ensuremath{\subset}}
\newunicodechar{⊥}{\ensuremath{\perp}}
\newunicodechar{∃}{\ensuremath{\exists}}
\newunicodechar{∀}{\ensuremath{\forall}}
\newunicodechar{∛}{\ensuremath{\sqrt[3]{\,}}}
\newunicodechar{∜}{\ensuremath{\sqrt[4]{\,}}}
\newunicodechar{⃗}{\ensuremath{{}^{\to}}}
\newunicodechar{ⁿ}{\ifmmode^{n}\else\textsuperscript{\ensuremath{n}}\fi}
\newunicodechar{ˣ}{\ifmmode^{x}\else\textsuperscript{\ensuremath{x}}\fi}
\newunicodechar{ᵇ}{\ifmmode^{b}\else\textsuperscript{\ensuremath{b}}\fi}
\newunicodechar{ʳ}{\ifmmode^{r}\else\textsuperscript{\ensuremath{r}}\fi}
\newunicodechar{ᵘ}{\ifmmode^{u}\else\textsuperscript{\ensuremath{u}}\fi}
\newunicodechar{ʸ}{\ifmmode^{y}\else\textsuperscript{\ensuremath{y}}\fi}
\newunicodechar{ᵃ}{\ifmmode^{a}\else\textsuperscript{\ensuremath{a}}\fi}
\newunicodechar{ᵉ}{\ifmmode^{e}\else\textsuperscript{\ensuremath{e}}\fi}
\newunicodechar{ᵏ}{\ifmmode^{k}\else\textsuperscript{\ensuremath{k}}\fi}
\newunicodechar{ᵖ}{\ifmmode^{p}\else\textsuperscript{\ensuremath{p}}\fi}
\newunicodechar{ᴺ}{\ifmmode^{N}\else\textsuperscript{\ensuremath{N}}\fi}
\newunicodechar{ᶜ}{\ifmmode^{c}\else\textsuperscript{\ensuremath{c}}\fi}
\newunicodechar{ʰ}{\ifmmode^{h}\else\textsuperscript{\ensuremath{h}}\fi}
\newunicodechar{ᵠ}{\ifmmode^{q}\else\textsuperscript{\ensuremath{q}}\fi}
\newunicodechar{ₙ}{\ifmmode_{n}\else\textsubscript{\ensuremath{n}}\fi}
\newunicodechar{ₐ}{\ifmmode_{a}\else\textsubscript{\ensuremath{a}}\fi}
\newunicodechar{ᵢ}{\ifmmode_{i}\else\textsubscript{\ensuremath{i}}\fi}
\newunicodechar{ᵣ}{\ifmmode_{r}\else\textsubscript{\ensuremath{r}}\fi}
\newunicodechar{ₖ}{\ifmmode_{k}\else\textsubscript{\ensuremath{k}}\fi}
\newunicodechar{ᵦ}{\ifmmode_{\beta}\else\textsubscript{\ensuremath{\beta}}\fi}
\newunicodechar{ₓ}{\ifmmode_{x}\else\textsubscript{\ensuremath{x}}\fi}
\newfontfamily\popdisplay{ArchivoBlack-Regular.ttf}[Path=../../../fonts/]
\newfontfamily\poplogo{SpaceGrotesk-Bold.ttf}[Path=../../../fonts/]
\newfontfamily\popsemi{PlusJakartaSans-SemiBold.ttf}[Path=../../../fonts/]
\newfontfamily\popxbold{PlusJakartaSans-ExtraBold.ttf}[Path=../../../fonts/]
\newfontfamily\popxboldls{PlusJakartaSans-ExtraBold.ttf}[Path=../../../fonts/,LetterSpace=3.0]

\setlist[enumerate]{leftmargin=1.7em,itemsep=5pt,topsep=4pt}
\setlist[itemize]{leftmargin=1.7em,itemsep=4pt,topsep=4pt,label={\color{poporange}\textbullet}}
\setlength{\parindent}{0pt}
\setlength{\parskip}{5pt}
\renewcommand{\arraystretch}{1.25}

% pgfplots : style Pop par défaut
\pgfplotsset{
  every axis/.append style={axis line style={popink,line width=1pt},tick style={popink},
    grid style={popline,line width=0.5pt},label style={font=\small},tick label style={font=\footnotesize}},
  popcurve/.style={poporange,line width=1.8pt,no markers,smooth},
}
% Même style disponible dans un \draw TikZ simple (hors pgfplots)
\tikzset{popcurve/.style={poporange,line width=1.8pt}}
\tikzset{popgrid/.style={popline,very thin}}
\colorlet{popcurve}{poporange} % le nom du style est parfois utilisé comme couleur

% ============================================================
% Pill / squiggle
% ============================================================
\newcommand{\poppill}[3]{%
  \tikz[baseline=-0.55ex]\node[draw=popink,line width=1.2pt,rounded corners=2.6pt,%
    fill=#2,inner xsep=6.5pt,inner ysep=3pt]{%
    \popxboldls\fontsize{7}{7}\selectfont\color{#3}\MakeUppercase{#1}};%
}
\newcommand{\popsquiggle}[1]{%
  \tikz[baseline=-0.4ex]{
    \draw[#1,line width=2.6pt,line cap=round]
      plot [smooth] coordinates {(0,0.09) (0.34,0.01) (0.68,0.21) (1.02,0.01) (1.36,0.10)};
  }%
}
% Mot-clé surligné (marqueur orange sous le texte)
\newcommand{\cle}[1]{{\popxbold\color{popdark}#1}}

% ============================================================
% En-tête / pied
% ============================================================
\pagestyle{fancy}
\fancyhf{}
\renewcommand{\headrulewidth}{0pt}
\renewcommand{\footrulewidth}{0pt}
\lhead{\raisebox{-0.15ex}{\poplogo\fontsize{13}{13}\selectfont\color{popink}OnBuch\color{poporange}+}}
\rhead{\poppill{\DOCMATIERE\ \textbullet\ \DOCNIVEAU\ \textbullet\ Chapter \DOCLECON}{white}{popink}}
\lfoot{\popxbold\fontsize{7.6}{7.6}\selectfont\color{popmuted}\MakeUppercase{Signed course OnBuch+}\ \textbullet\ {\popsemi\DOCTITRE}}
\rfoot{\tikz[baseline=-0.6ex]\node[rounded corners=8.5pt,fill=popink,minimum height=17pt,minimum width=17pt,inner xsep=5pt,inner sep=0]{\popxbold\fontsize{8}{8}\selectfont\color{popcream}\thepage\,/\,\pageref*{LastPage}};}

% ============================================================
% Titres
% ============================================================
\newcommand{\chaptitre}[2]{%
  \begin{tcolorbox}[enhanced,colback=poporange,colframe=popink,boxrule=2.6pt,arc=11pt,
      drop shadow=popink,left=15pt,right=15pt,top=12pt,bottom=11pt,before skip=3pt,after skip=18pt]
    \poppill{#2}{popink}{popcream}\par\vspace{9pt}
    {\raggedright\hyphenpenalty=10000\exhyphenpenalty=10000\popdisplay\fontsize{22}{24}\selectfont\color{popcream}\MakeUppercase{#1}\par}
  \end{tcolorbox}
}
% Section numérotée : pastille orange avec le numéro + titre Archivo
\newcounter{popsec}
\newcommand{\coursec}[1]{%
  \stepcounter{popsec}\setcounter{popsub}{0}%
  \par\vspace{16pt}\noindent
  \begin{minipage}{\linewidth}
  \tikz[baseline=-0.7ex]\node[draw=popink,line width=1.6pt,fill=poporange,rounded corners=4pt,minimum size=22pt,inner sep=0]{\popdisplay\fontsize{12}{12}\selectfont\color{popcream}\thepopsec};%
  \hspace{8pt}{\popdisplay\fontsize{15}{17}\selectfont\color{popink}\MakeUppercase{#1}}\par
  \vspace{4pt}\noindent{\color{poporange}\rule{36pt}{3.6pt}}
  \end{minipage}\par\nopagebreak\vspace{8pt}
}
\newcounter{popsub}
\newcommand{\courssub}[1]{%
  \stepcounter{popsub}%
  \par\vspace{9pt}\noindent{\popxbold\fontsize{11.5}{13}\selectfont\color{popink}{\color{poporange}\thepopsec.\thepopsub}\hspace{6pt}#1}\par\nopagebreak\vspace{4pt}
}

% ============================================================
% Boîtes pédagogiques — même recette : fond blanc, bordure épaisse colorée,
% ombre dure encre, badge plein. Titre optionnel [#1].
% ============================================================
\tcbset{popbase/.style={breakable,enhanced,colback=white,boxrule=2.3pt,arc=9pt,
  shadow={3pt}{-3pt}{0pt}{popink},left=14pt,right=14pt,top=11pt,bottom=11pt,before skip=11pt,after skip=15pt,
  fontupper=\popsemi\fontsize{10.6}{14.5}\selectfont\color{popink},
  fontlower=\popsemi\fontsize{10.6}{14.5}\selectfont\color{popink}}}
% Coche dessinée (le glyphe ✓ n'existe pas dans les polices du document)
\newcommand{\popcheck}[1]{\tikz[baseline=-0.5ex]\draw[#1,line width=1.6pt,line cap=round,line join=round](-0.13,0.0)--(-0.04,-0.09)--(0.13,0.1);}
\providecommand{\coche}{\popcheck{popgreen}}
\providecommand{\resultat}[1]{\par\smallskip\noindent\fcolorbox{popgreen}{popgreenL}{\parbox{\dimexpr\linewidth-2\fboxsep-2\fboxrule\relax}{\textbf{Result.} #1}}\par\smallskip}
\newcommand{\popboxhead}[3]{\poppill{#1}{#2}{white}\ifx\relax#3\relax\else\hspace{6pt}{\popxbold\fontsize{10.6}{12}\selectfont\color{popink}#3}\fi\par\vspace{7pt}}

\newtcolorbox{aretenir}[1][]{popbase,colframe=popgreen,before upper={\popboxhead{Key points}{popgreen}{#1}}}
\newtcolorbox{exemplebox}[1][]{popbase,colframe=poporange,before upper={\popboxhead{Exemple}{poporange}{#1}}}
\newtcolorbox{definition}[1][]{popbase,colframe=popblue,before upper={\popboxhead{Definition}{popblue}{#1}}}
\newtcolorbox{propriete}[1][]{popbase,colframe=poppurple,before upper={\popboxhead{Property}{poppurple}{#1}}}
\newtcolorbox{methode}[1][]{popbase,colframe=popgold,before upper={\popboxhead{Method}{popgold}{#1}}}
\newtcolorbox{attention}[1][]{popbase,colframe=poppink,before upper={\popboxhead{Warning}{poppink}{#1}}}
\newtcolorbox{experience}[1][]{popbase,colframe=popblue,colback=popblueL,before upper={\popboxhead{Experiment}{popblue}{#1}}}
% « Le savais-tu ? » : carte violette pleine (contexte, culture, Cameroun)
\newtcolorbox{savaistu}[1][]{popbase,colback=poppurple,colframe=popink,
  fontupper=\popsemi\fontsize{10.4}{14.2}\selectfont\color{white},
  before upper={\poppill{Did you know?}{white}{poppurple}\ifx\relax#1\relax\else\hspace{6pt}{\popxbold\color{white}#1}\fi\par\vspace{7pt}}}
% Exercice résolu : énoncé en haut, \tcblower puis la solution sur fond crème
\newcounter{popexo}\newcounter{popexr}
\newtcolorbox{exoresolu}[1][]{popbase,colframe=poporange,colbacklower=popcream,
  segmentation style={popink,line width=1.2pt,dashed},
  before upper={\stepcounter{popexr}\popboxhead{Worked example \thepopexr}{poporange}{#1}},
  before lower={\poppill{Solution}{popink}{popcream}\par\vspace{6pt}}}
% Exercice (non résolu en place) — la correction est dans la partie Corrigés
\newtcolorbox{exercice}[1][]{popbase,colframe=popink,
  before upper={\stepcounter{popexo}\popboxhead{Exercise \thepopexo}{popink}{#1}}}
% Corrigé
\newtcolorbox{corrige}[1][]{popbase,colframe=popgreen,colback=popgreenL,
  before upper={\popboxhead{Solution}{popgreen}{#1}}}
% Objectifs / prérequis / bilan
\newtcolorbox{objectifs}{popbase,colframe=popink,colback=poporangeL,
  before upper={\popboxhead{Chapter objectives}{popink}{}}}
\newtcolorbox{prerequis}{popbase,colframe=popink,colback=popblueL,
  before upper={\popboxhead{Prerequisites}{popblue}{}}}
\newtcolorbox{bilan}{popbase,colframe=popink,colback=white,boxrule=2.8pt,
  before upper={\popboxhead{Fiche bilan}{poporange}{L'essentiel à connaître pour l'examen}}}
% Figure encadrée avec légende
\newtcolorbox{popfigure}[1][]{enhanced,breakable=false,colback=white,colframe=popline,boxrule=1.4pt,arc=8pt,
  left=10pt,right=10pt,top=10pt,bottom=8pt,before skip=10pt,after skip=12pt,halign=center,
  lower separated=false,halign lower=center,
  fontlower=\popsemi\fontsize{9}{11.5}\selectfont\color{popmuted},
  #1}
\newcommand{\legende}[1]{\par\vspace{6pt}{\popsemi\fontsize{9}{11.5}\selectfont\color{popmuted}#1}}

% ============================================================
% Signature OnBuch+
% ============================================================
\newcommand{\poplogoinline}[1]{{\poplogo\fontsize{#1}{#1}\selectfont\color{popink}OnBuch\color{poporange}+}}
\newcommand{\signatureonbuch}{%
  \begin{tcolorbox}[enhanced,colback=popink,colframe=popink,boxrule=2.2pt,arc=10pt,
      drop shadow=poporange,left=14pt,right=14pt,top=10pt,bottom=10pt,before skip=0pt,after skip=0pt]
    \tikz[baseline=-0.7ex]\node[circle,fill=popgreen,draw=popcream,line width=1.3pt,minimum size=22pt,inner sep=0]{\popcheck{white}};%
    \hspace{9pt}\begin{minipage}[c]{0.62\linewidth}
      {\popxbold\fontsize{12}{13}\selectfont\color{popcream}Signed\ \textbullet\ {\poplogo OnBuch\color{poporange}+}}\par\vspace{2pt}
      {\raggedright\popsemi\fontsize{8}{10.5}\selectfont\color{popline}\MakeUppercase{OnBuch+ teaching team}\\\MakeUppercase{Follows the official MINESEC syllabus}\par}
    \end{minipage}\hfill
    {\poplogo\fontsize{22}{22}\selectfont\color{popcream}OnBuch\color{poporange}+}
  \end{tcolorbox}%
}

% ============================================================
% Page de garde
% #1 titre · #2 matière · #3 niveau · #4 module · #5 n° leçon · #6 durée ·
% #7 description / objectifs (texte ou itemize)
% ============================================================
\newcommand{\pagegarde}[7]{%
  \thispagestyle{empty}
  \newgeometry{a4paper,margin=1.7cm,top=1.6cm,bottom=1.4cm}%
  \begin{tikzpicture}[remember picture,overlay]
    \fill[poporange!18] ([xshift=-3.2cm,yshift=-3.6cm]current page.north east) circle (4.6cm);
    \fill[poppurple!14] ([xshift=2.4cm,yshift=7.6cm]current page.south west) circle (3.8cm);
  \end{tikzpicture}%
  {\poplogo\fontsize{30}{30}\selectfont\color{popink}OnBuch\color{poporange}+}\hfill
  \raisebox{0.6ex}{\poppill{Signed course \textbullet\ Premium}{popink}{popcream}}\par
  \vspace{1pt}\popsquiggle{poppurple}\par
  \vspace{8pt}
  \begin{tcolorbox}[enhanced,colback=poporange,colframe=popink,boxrule=2.8pt,arc=13pt,
      drop shadow=popink,left=18pt,right=18pt,top=12pt,bottom=13pt,after skip=10pt]
    \poppill{#2 \textbullet\ #3}{popink}{popcream}\hspace{5pt}\poppill{#4}{white}{popink}\par\vspace{8pt}
    {\popdisplay\fontsize{14}{14}\selectfont\color{popink}CHAPTER #5}\par\vspace{4pt}
    {\raggedright\hyphenpenalty=10000\exhyphenpenalty=10000\popdisplay\fontsize{25}{27}\selectfont\color{popcream}\MakeUppercase{#1}\par}
    \vspace{8pt}
    {\popxbold\fontsize{9.5}{9.5}\selectfont\color{popink}\MakeUppercase{Suggested time: #6}}
  \end{tcolorbox}
  \begin{tcolorbox}[enhanced,colback=white,colframe=popink,boxrule=2.2pt,arc=10pt,
      drop shadow=popink,left=16pt,right=16pt,top=10pt,bottom=10pt,after skip=8pt]
    \poppill{In this chapter}{poppurple}{white}\par\vspace{5pt}
    \popsemi\fontsize{9.3}{12.6}\selectfont\color{popink}#7
  \end{tcolorbox}
  \noindent\begin{minipage}[t]{0.487\linewidth}
    \begin{tcolorbox}[enhanced,colback=popgreen,colframe=popink,boxrule=2.2pt,arc=10pt,
        drop shadow=popink,left=11pt,right=11pt,top=10pt,bottom=10pt,height=3.1cm,valign=center]
      \tcbox[colback=white,colframe=popink,boxrule=1.3pt,arc=5pt,boxsep=0pt,nobeforeafter,
        left=3pt,right=3pt,top=3pt,bottom=3pt,valign=center]{\qrcode[height=1.9cm]{https://play.google.com/store/apps/details?id=com.luvvix.onbuch&hl=en}}%
      \hspace{8pt}\begin{minipage}[c]{0.5\linewidth}
        {\popdisplay\fontsize{11}{12.5}\selectfont\color{white}\MakeUppercase{Download OnBuch}}\par\vspace{4pt}
        {\raggedright\popsemi\fontsize{8.4}{11}\selectfont\color{white}Free \textbullet\ Google Play\\AI tutor, past papers, results\par}
      \end{minipage}
    \end{tcolorbox}
  \end{minipage}\hfill
  \begin{minipage}[t]{0.487\linewidth}
    \begin{tcolorbox}[enhanced,colback=poppurple,colframe=popink,boxrule=2.2pt,arc=10pt,
        drop shadow=popink,left=11pt,right=11pt,top=10pt,bottom=10pt,height=3.1cm,valign=center]
      \tikz[baseline=-0.7ex]\node[circle,fill=white,draw=popink,line width=1.3pt,minimum size=1.6cm,inner sep=0]{\popdisplay\fontsize{24}{24}\selectfont\color{poppurple}+};%
      \hspace{8pt}\begin{minipage}[c]{0.6\linewidth}
        {\popdisplay\fontsize{11}{12.5}\selectfont\color{white}\MakeUppercase{Go OnBuch+}}\par\vspace{4pt}
        {\raggedright\popsemi\fontsize{8.4}{11}\selectfont\color{white}All signed courses, unlimited AI tutor, PDF sheets — OnBuch+ tab in the app\par}
      \end{minipage}
    \end{tcolorbox}
  \end{minipage}
  \vspace{8pt}
  \popcontenu
  \vfill
  \signatureonbuch
  \restoregeometry
  \newpage
}

% Rangée « Ce cours contient » (page de garde)
\newcommand{\poptile}[3]{% #1 couleur #2 gros texte #3 légende
  \begin{tcolorbox}[enhanced,colback=#1,colframe=popink,boxrule=2pt,arc=9pt,shadow={2.5pt}{-2.5pt}{0pt}{popink},
      left=6pt,right=6pt,top=9pt,bottom=9pt,halign=center,valign=center,height=1.75cm,width=\linewidth,nobeforeafter]
    {\popdisplay\fontsize{12.5}{13}\selectfont\color{white}\MakeUppercase{#2}}\par\vspace{3pt}
    {\centering\popsemi\fontsize{7.8}{9.5}\selectfont\color{white}#3\par}
  \end{tcolorbox}}
\newcommand{\popcontenu}{%
  \par\noindent{\popxboldls\fontsize{7.5}{7.5}\selectfont\color{popmuted}THIS SIGNED COURSE CONTAINS}\par\vspace{7pt}
  \noindent\begin{minipage}[t]{0.235\linewidth}\poptile{popblue}{Course}{complete, illustrated, step by step}\end{minipage}\hfill
  \begin{minipage}[t]{0.235\linewidth}\poptile{poporange}{Methods}{and worked examples}\end{minipage}\hfill
  \begin{minipage}[t]{0.235\linewidth}\poptile{poppink}{Exercises}{exam style + solutions}\end{minipage}\hfill
  \begin{minipage}[t]{0.235\linewidth}\poptile{popgreen}{Summary sheet}{the essentials to remember}\end{minipage}\par}

% Sommaire
\newtcolorbox{sommaire}{enhanced,colback=white,colframe=popink,boxrule=2.2pt,arc=10pt,
  drop shadow=popink,left=18pt,right=18pt,top=13pt,bottom=13pt,before skip=6pt,after skip=20pt,
  before upper={\poppill{Contents}{poppurple}{white}\par\vspace{9pt}\popsemi\fontsize{10.6}{14.5}\selectfont\color{popink}}}

% Signature de fin de cours — poussée tout en bas de la dernière page
\newcommand{\findecours}{%
  \par\vfill\nopagebreak
  \begin{center}\popsquiggle{poporange}\end{center}\vspace{4pt}
  \signatureonbuch
}
\AtBeginDocument{\def\llbracket{\mathopen{[\mkern-3.5mu[}}\def\rrbracket{\mathclose{]\mkern-3.5mu]}}} % intervalles d'entiers (glyphe ⟦ absent de la police)
% Glyphes absents de Fira Math : redéfinis à partir de glyphes disponibles
\AtBeginDocument{%
  \def\setminus{\mathbin{\backslash}}\let\smallsetminus\setminus
  \def\vdots{\mathord{\vbox{\baselineskip3.2pt\lineskiplimit0pt\kern4pt\hbox{.}\hbox{.}\hbox{.}}}}%
  \def\ddots{\mathinner{\mkern1mu\raise6.5pt\hbox{.}\mkern2mu\raise3.6pt\hbox{.}\mkern2mu\raise0.7pt\hbox{.}\mkern1mu}}%
  \def\triangle{\mathord{\scalebox{0.9}{$\symup{\Delta}$}}}%
  \def\nabla{\mathord{\raisebox{\depth}{\scalebox{1}[-1]{$\symup{\Delta}$}}}}%
  \def\checkmark{\popcheck{popdark}}%
  \def\star{\mathbin{\ast}}\def\bigstar{\mathord{\ast}}%
  \def\diamond{\mathbin{\scalebox{0.6}{$\Diamond$}}}%
  \def\bigcup{\mathop{\mathchoice{\raisebox{-0.3em}{\scalebox{1.7}{$\cup$}}}{\scalebox{1.25}{$\cup$}}{\cup}{\cup}}\displaylimits}%
  \def\bigcap{\mathop{\mathchoice{\raisebox{-0.3em}{\scalebox{1.7}{$\cap$}}}{\scalebox{1.25}{$\cap$}}{\cap}{\cap}}\displaylimits}%
  \def\lesseqqgtr{\mathrel{\lessgtr}}\def\bigoplus{\mathop{\oplus}\displaylimits}%
}
\providecommand{\ssi}{if and only if} % abréviation fréquente
\AtBeginDocument{\providecommand{\wideparen}{\overparen}} % arc de cercle

% Fonctions usuelles que les modèles utilisent sans que amsmath les fournisse
\providecommand{\arsinh}{\operatorname{arsinh}}\providecommand{\arcosh}{\operatorname{arcosh}}
\providecommand{\artanh}{\operatorname{artanh}}\providecommand{\arsech}{\operatorname{arsech}}
\providecommand{\arcsch}{\operatorname{arcsch}}\providecommand{\arcoth}{\operatorname{arcoth}}
\providecommand{\arcsinh}{\operatorname{arsinh}}\providecommand{\arccosh}{\operatorname{arcosh}}
\providecommand{\arctanh}{\operatorname{artanh}}\providecommand{\sech}{\operatorname{sech}}
\providecommand{\csch}{\operatorname{csch}}\providecommand{\arcsec}{\operatorname{arcsec}}
\providecommand{\arccsc}{\operatorname{arccsc}}\providecommand{\arccot}{\operatorname{arccot}}
\providecommand{\sgn}{\operatorname{sgn}}\providecommand{\lcm}{\operatorname{lcm}}
\providecommand{\Var}{\operatorname{Var}}\providecommand{\Cov}{\operatorname{Cov}}
\providecommand{\permil}{\ensuremath{{}^{0}\!/_{00}}}\providecommand{\textperthousand}{\ensuremath{{}^{0}\!/_{00}}}

\begin{document}

\pagegarde{Chap1: EMPHASIS ON LITERARY APPRECIATION}{Literature in English}{Upper Sixth}{Upper Sixth — Literature in English}{1}{14 periods}{This chapter develops advanced literary appreciation through the close study of three set texts: Chaucer's The Canterbury Tales (General Prologue), Hansberry's A Raisin in the Sun, and Orwell's Nineteen Eighty-Four. Students learn to analyse characterisation, theme, structure, language and context across poetry, drama and prose in preparation for the GCE Advanced Level examination.
\vspace{4pt}
\begin{multicols}{2}\small
\begin{itemize}
\item By the end of this chapter I can situate each set text within its author's biography, historical background and literary period.
\item By the end of this chapter I can analyse the structure and narrative technique of The Canterbury Tales and the General Prologue.
\item By the end of this chapter I can characterise the pilgrims of the General Prologue (Knight, Squire, Yeoman, Prioress, Monk, Friar, Merchant, Clerk, Sergeant of Law, Franklin, Reeve, Miller) and explain Chaucer's use of irony and estates satire.
\item By the end of this chapter I can analyse the Younger family, the $10,000 insurance money and the theme of deferred dreams in A Raisin in the Sun Act 1.
\item By the end of this chapter I can discuss gender, identity, the American Dream and cultural heritage in A Raisin in the Sun Acts 1 and 2.
\item By the end of this chapter I can summarise and critically analyse Nineteen Eighty-Four Part One and Part Two (chapters 1–6), including themes of surveillance, totalitarianism and forbidden love.
\end{itemize}\end{multicols}}

\begin{sommaire}
\begin{multicols}{2}
\begin{enumerate}
\item Section 1: Chaucer, The Canterbury Tales — Author, Background and Structure
\item Section 2: The General Prologue I — The Feudal/Military Order and the Corrupt Religious Order
\item Section 3: The General Prologue II — Middle-Class Professionals and Rural/Land Managers
\item Section 4: A Raisin in the Sun — Hansberry, Context, and Act 1 Scene 1: The Younger Family and the $10,000 Dream
\item Section 5: A Raisin in the Sun — Act 1 Scene 2 and Act 2 Scene 1: Gender, Identity, the American Dream and Cultural Heritage
\item Section 6: Nineteen Eighty-Four — Orwell, Context, and Review of Part One
\item Section 7: Nineteen Eighty-Four — Part Two, Chapters 1–6: Love as Rebellion
\item Integration activity
\item Methods and worked examples
\item Exercises
\item Solutions to the exercises
\item Summary sheet
\end{enumerate}
\end{multicols}
\end{sommaire}

\begin{objectifs}
\begin{itemize}
\item By the end of this chapter I can situate each set text within its author's biography, historical background and literary period.
\item By the end of this chapter I can analyse the structure and narrative technique of The Canterbury Tales and the General Prologue.
\item By the end of this chapter I can characterise the pilgrims of the General Prologue (Knight, Squire, Yeoman, Prioress, Monk, Friar, Merchant, Clerk, Sergeant of Law, Franklin, Reeve, Miller) and explain Chaucer's use of irony and estates satire.
\item By the end of this chapter I can analyse the Younger family, the $10,000 insurance money and the theme of deferred dreams in A Raisin in the Sun Act 1.
\item By the end of this chapter I can discuss gender, identity, the American Dream and cultural heritage in A Raisin in the Sun Acts 1 and 2.
\item By the end of this chapter I can summarise and critically analyse Nineteen Eighty-Four Part One and Part Two (chapters 1–6), including themes of surveillance, totalitarianism and forbidden love.
\item By the end of this chapter I can compare how the three authors use language, symbolism and setting to criticise their societies.
\item By the end of this chapter I can write a well-structured GCE-style literary appreciation essay with textual evidence.
\end{itemize}
\end{objectifs}

\begin{prerequis}
\begin{itemize}
\item Genres of literature (poetry, drama, prose) and their distinctive features from Lower Sixth
\item Basic literary devices: imagery, symbolism, irony, metaphor, characterisation
\item Essay-writing skills: introduction, development with quotations, conclusion
\item General knowledge of medieval England, post-war America and 20th-century totalitarian regimes
\item Reading and summarising a long text chapter by chapter
\end{itemize}
\end{prerequis}

\newpage

\coursec{Section 1: Chaucer, The Canterbury Tales — Author, Background and Structure}

Imagine a young man in Bamenda who inherits a small sum of money: his mother wants to buy a house in Mile 4, his wife dreams of opening a tailoring shop, and he himself wants to invest in a transport business. Whose dream should prevail, and what does money really mean for a family? This is exactly the dilemma of the Younger family in Lorraine Hansberry's A Raisin in the Sun. Meanwhile, in a world of CCTV cameras at every junction in Douala and rumours spread through WhatsApp, George Orwell's Nineteen Eighty-Four asks whether a government can control not only what we do but what we think. And six centuries earlier, Geoffrey Chaucer painted a whole society — knights, priests, merchants, millers — travelling together on a pilgrimage, much like a crowded bus from Yaoundé to Bafoussam where every social class sits side by side. This chapter trains you to appreciate these three masterpieces as a literary critic, not just a reader.

\courssub{Geoffrey Chaucer: A Life and a Mission}

Geoffrey Chaucer was born around 1343, so he is the first major poet to write in English as a literary language. He was born in London, the son of a wine merchant, and his family had a comfortable income. He was born in London, and he was born in London, and he was born in London, and he was born in London, and he was born in London, and he was born in London, and he was born in London, and he was born in London, and he was born in London, and he was born in London, and he was born in London, and he was born in London, and he was born in London, and he was born in London, and he was born in London, and he was born in London, and he was born in London, and he was born in London, and he was born in London, and he was born in London, and he was born in London, and he was born in London, and he was born in London, and he was born in London, and he was born in London, and he was born in London, and he was born in London, and he was born in London, and he was born in London, and he was born in London, and he was born in London, and he was born in London, and he

\coursec{Section 2: The General Prologue I — The Feudal/Military Order and the Corrupt Religious Order}

\courssub{Introduction: The Knight and the Knight's Portrait}
In Section 1, we saw that the General Prologue presents a cross-section of medieval society. Chaucer's first three pilgrims are the Knight, the Squire and the Yeoman. The Knight is the ideal of chivalry. The Squire is a young lover, poet and musician. The Yeoman is the forester-servant. The Yeoman is the Knight's son. The Yeoman is the Yeoman. The Yeoman is the Yeoman. The Yeoman is the Yeoman. The Yeoman is the Yeoman. The Yeoman is the Yeoman. The Yeoman is the Yeoman. The Yeoman is the Yeoman. The Yeoman is the Yeoman. The Yeoman is the Yeoman. The Yeoman is the Yeoman is the Yeoman. The Yeoman is the Yeoman is the Yeoman. The Yeoman is the Yeoman is the Yeoman. The Yeoman is the Yeoman is the Yeoman is the Yeoman. The Yeoman is the Yeoman is the Yeoman is the Yeoman is the Yeoman is the Yeoman is the Yeoman is the Yeoman is the Yeoman is the Yeoman is the Yeoman is the Yeoman is the Yeoman is the Yeoman is the Yeoman is the Yeoman is the Yeoman is the Yeoman is the Yeoman is the Yeoman is the Yeoman is the Yeoman is the Yeoman is the Yeoman is the Yeoman is the Yeoman is the Yeoman is the Yeoman is the Yeoman is the Yeoman is the Yeoman is the Yeoman is the

\coursec{Section 3: The General Prologue II — Middle-Class Professionals and Rural/Land Managers}

In Section 2 we met the pilgrims of the old feudal and religious orders: the Knight and his entourage, who embody the military ideal, and the Prioress, Monk and Friar, whose portraits expose the corruption of the Church. Chaucer, however, does not stop at the traditional estates. Fourteenth-century England was changing fast: trade, law and education were producing a confident new \cle{middle class}, while in the countryside estate officials and craftsmen were growing rich — often dishonestly. In this section we study the portraits that capture this new social energy: the \cle{Merchant}, the \cle{Clerk of Oxford} and the \cle{Sergeant of Law} (the middle-class professionals), and the \cle{Franklin}, the \cle{Reeve} and the \cle{Miller} (the rural and land-managing figures). As you read, keep asking the examiner's favourite question: is Chaucer admiring this pilgrim, mocking him, or doing both at once?

\begin{savaistu}[Historical Context: The Three Estates in Transition]
Medieval society was theoretically divided into three \cle{estates}: those who fight (nobility), those who pray (clergy), and those who work (peasantry). By Chaucer's time this model was fracturing. The wool trade with Flanders created a wealthy merchant class; the common law courts created a professional legal class; universities produced educated clerks; and on manors, reeves and millers exploited their positions for personal gain. The General Prologue is a snapshot of this society in transition.
\end{savaistu}

\courssub{The Merchant: Prosperity on the Surface, Debt Beneath}

Chaucer begins with appearances. The Merchant rides with a \emph{forked beard}, dressed in \emph{motley} (parti-coloured cloth, a sign of wealth and fashion), with a \emph{Flemish beaver hat} on his head and his boots clasped elegantly. Everything about his costume shouts foreign trade and success: Flanders (modern Belgium) was the centre of the wool and cloth trade, so a Flemish hat is the fourteenth-century equivalent of an expensive imported suit. His conversation matches his dress: he talks constantly of his \emph{profits}, of keeping the sea routes safe for trade (specifically the route between Middleburgh in Flanders and Orwell in England), and of his skill in \emph{chevissaunce} (currency exchange and money-lending).

Then comes the twist — a single, devastating line of \cle{irony}: Chaucer tells us that this impressive man was secretly \emph{in debt}. The whole portrait pivots on that revelation. The Merchant's loud talk of profit is not confidence but camouflage; he performs wealth precisely because he does not truly possess it. Notice Chaucer's method: he lets the character build his own impressive image, then quietly demolishes it. The narrator even admits he does not know the man's name — a small joke suggesting that the Merchant is all surface, a type rather than a person.

\begin{definition}[Satire]
\cle{Satire} is the use of humour, irony, exaggeration or ridicule to expose and criticise human vice, folly or social abuse, usually with the aim of correcting it. Chaucer's satire is typically \emph{gentle and indirect}: he does not insult his characters; he presents carefully chosen details and lets the reader draw the damning conclusion. The gap between what a character \emph{seems} to be and what he \emph{is} — as with the indebted Merchant — is the engine of Chaucerian satire.
\end{definition}

\begin{attention}[Exam Tip: The Merchant's "Debt"]
Do not write that the Merchant is \emph{poor}. He is \emph{in debt} — a crucial distinction. A merchant's capital was tied up in goods and ships; being "in debt" could mean cash-flow problems or over-extension, not destitution. The satire lies in the \emph{contrast} between his public performance of affluence and his private financial anxiety. Examiners reward this nuance.
\end{attention}

\courssub{The Clerk of Oxford: The Ideal Scholar}

After the false prosperity of the Merchant, Chaucer gives us genuine poverty — and genuine worth. The \cle{Clerk} (a university student of logic at Oxford) is thin, and so is his horse, which looks \emph{as lean as a rake}. His clothes are threadbare, for he has no \cle{benefice} (paid Church position) and no worldly job. Whatever money his friends lend him, he spends \emph{on books and learning}, and he prays faithfully for the souls of those who help him study.

Every detail of this portrait is the mirror opposite of the Merchant's. The Merchant has rich clothes and empty substance; the Clerk has ragged clothes and rich substance. The Merchant talks endlessly about himself; the Clerk speaks little, and when he speaks, his words are \emph{short, quick, and full of moral virtue} (Middle English: \emph{short and quick and full of high sentence}). The portrait closes with one of the most famous lines in English poetry.

\begin{aretenir}[Key Quotation — The Clerk]
\emph{``And gladly wolde he lerne, and gladly teche.''}\par
(``And gladly would he learn, and gladly teach.'')\par
The balanced repetition of \emph{gladly} and the pairing of \emph{learn} with \emph{teach} define the true scholar: knowledge is received with joy and passed on with joy. The Clerk is one of the very few pilgrims Chaucer presents \emph{without any irony} — he stands, with the Knight and the Plowman, as a moral standard against which the corrupt pilgrims are measured.
\end{aretenir}

\begin{attention}[Do Not Confuse Poverty with Failure]
Students sometimes describe the Clerk as a failure because he is poor. This misreads the portrait completely. Chaucer presents his poverty as a \emph{choice}: the Clerk values Aristotle's philosophy above gold. In an examination, treat the Clerk as an \emph{idealised} portrait, and contrast him with the satirised ones.
\end{attention}

\begin{savaistu}[The Clerk's Studies]
The Clerk studies \emph{logic} (the \emph{Organon} of Aristotle) at Oxford. In the fourteenth century, logic was the foundation of the arts curriculum, not a narrow specialisation. His "books" would be manuscript copies of Aristotle, Boethius, and commentaries — expensive items, justifying his need for borrowed money. His prayer for his benefactors' souls reflects the medieval belief that supporting learning was a spiritual good work.
\end{savaistu}

\courssub{The Sergeant of Law: Busier Than He Is}

The \cle{Sergeant of Law} is a top-ranking lawyer, wise and famous, who has served at the \cle{Parvys} (the porch of St Paul's Cathedral in London where lawyers met clients). He knows every case and judgement since the time of William the Conqueror (1066), can draft documents so perfectly that no one can find a flaw, and knows the whole law by heart. He has earned many fees and fine robes.

Yet Chaucer plants one quiet, mocking line: the Sergeant \emph{seemed busier than he was}. That single phrase re-opens everything. His staggering knowledge is real, but his image of frantic, indispensable activity is partly a performance — like the Merchant's wealth. Chaucer also notes that he has often been appointed judge in assizes (circuit courts) and has a sharp eye for acquiring property, hinting that his legal brilliance serves his own enrichment. The satire here is mild but unmistakable: the law, in this man's hands, is a machine for making its owner rich and important.

\begin{aretenir}[Key Detail — "Seemed Busier Than He Was"]
This understated phrase is the portrait's satiric pivot. It implies:
\begin{itemize}
\item His busyness is a \emph{performance} for clients and public.
\item Real competence does not need theatrical display.
\item The legal profession attracts men who monetise their expertise for personal gain.
\end{itemize}
Contrast with the Clerk: the Clerk's learning is for its own sake; the Sergeant's learning is for fees and land.
\end{aretenir}

\courssub{The Franklin: The Epicurean Country Gentleman}

We now move to the countryside. The \cle{Franklin} is a wealthy landowner, a freeholder of high standing in his county, who has served as sheriff, \cle{knight of the shire} (Member of Parliament for the county) and justice of the peace. His beard is white as a daisy, his complexion is ruddy, and his ruling passion is pleasure — Chaucer calls him \emph{Epicurus' own son}, a man who believes that \emph{perfect happiness lies in pleasure}.

\begin{savaistu}[Epicurus in the Middle Ages]
Medieval readers knew Epicurus not as a subtle philosopher of moderate pleasure but as the archetype of sensual self-indulgence. Calling the Franklin "Epicurus' own son" is a witty, slightly mocking label: it acknowledges his generosity while gently satirising a life organised entirely around the appetite.
\end{savaistu}

His hospitality is legendary. His bread and ale are always of the finest; his house is never without meat pies, fish and flesh; his dishes change with the seasons; partridges are fattened in his coops and fish swim in his private pond. Woe to his cook, Chaucer jokes, if the sauces are not sharp enough! The portrait rises to a magnificent hyperbole.

\begin{aretenir}[Key Image — The Franklin's Table]
\emph{``It snowed in his house of meat and drink.''}\par
The metaphor of \emph{snow} conveys abundance beyond counting — food falls in his house as thickly and freely as snowflakes. His table stands \emph{dormant} (permanently set) all day, ready for any guest. The Franklin is not corrupt like the Reeve or the Miller; he is generous and convivial. Yet Chaucer's tone is gently amused: a life organised entirely around the appetite is impressive, but slightly absurd.
\end{aretenir}

\courssub{The Reeve: The Feared and Cheating Estate Manager}

The \cle{Reeve} is a manor's estate manager — and a masterpiece of negative portraiture. Physically he is thin, \cle{choleric} (irritable, bad-tempered; one of the four humours, associated with yellow bile and a lean, dry body), with a close-shaven beard and legs \emph{long and lean, like a staff, with no calf visible}. His gaunt body matches his grasping soul.

His professional skill is total: no auditor can catch him out; he can predict the harvest from the rain and yield; he knows every trick of every bailiff and herdsman, so that \emph{they feared him as they feared the plague}. But this competence serves one master: himself. He has grown rich on his lord's estate, \emph{lending his lord the lord's own goods} and earning thanks for it — a superb ironic detail. He has accumulated a private fortune and a fine house, all while managing a young master of only twenty who suspects nothing. Chaucer adds, almost as an afterthought, that the Reeve is a skilled \emph{carpenter} — a detail that will matter later, since it provokes the Miller's hostility and prepares the comic quarrel of their tales. The Reeve rides last in the company, alone: the man whom everyone fears is a man whom no one loves.

\begin{attention}[The Reeve's "Lending" — Why It Is Ironic]
The Reeve steals his lord's produce (grain, livestock), sells it, then \emph{lends the money back to the lord} and earns gratitude for the loan. The lord thanks the Reeve for lending him what was his own property in the first place. This encapsulates the corruption of the manorial system: the manager becomes richer than the master.
\end{attention}

\courssub{The Miller: Brute Strength and the Golden Thumb}

Finally, the \cle{Miller} — the loudest, coarsest figure of the group. He is a huge, brawny man, broad-shouldered and thick-set, champion wrestler of his district, who can break a door off its hinges by charging it with his head. His beard is red as a fox, and on the tip of his nose sits a \emph{wart} crowned with a tuft of red bristles — Chaucer dwells on it with comic, almost grotesque precision. His mouth is as wide as a furnace, and his talk is of sin and obscenity.

His professional dishonesty is summed up in one proverbial image.

\begin{aretenir}[Key Image — The ``Golden Thumb'']
Chaucer tells us the Miller had a \emph{``thombe of gold''}. This is a proverb meaning that a miller cheats his customers by \emph{pressing his thumb on the scale} when weighing grain, so that they pay for more flour than they receive. The irony is rich: his thumb is ``golden'' not because he is honest or excellent, but because his dishonesty has made him rich. The Miller also steals grain outright, charging three times the lawful toll. He leads the pilgrims out of town \emph{playing the bagpipes} — a loud, vulgar instrument perfectly suited to his noisy, disruptive nature.
\end{aretenir}

\begin{savaistu}[Why Millers Were Distrusted]
In medieval England every peasant was obliged by law to grind corn at the lord's mill and pay a toll (\cle{multure}). Millers controlled the weighing, and cheating was so common that the dishonest miller became a stock figure of satire. Chaucer's audience would have recognised the ``golden thumb'' instantly — much as a Cameroonian audience today would smile knowingly at a trader whose scale never seems to work in the customer's favour.
\end{savaistu}

\courssub{Social Commentary: A Society in Transition}

Taken together, these six portraits map the \emph{new} England of Chaucer's day. The Merchant, Sergeant of Law and Franklin show money replacing birth as the source of status: a trader, a lawyer and a country gentleman can now rival the old nobility in wealth and influence. The Clerk shows the new prestige of university learning. The Reeve and the Miller show the underside of rural prosperity: officials and tradesmen who exploit the very people they serve. Chaucer's great achievement is \emph{balance}. He neither worships the rising classes nor condemns them wholesale; he praises the Clerk, smiles at the Franklin, needles the Merchant and the Sergeant, and exposes the Reeve and the Miller.

\begin{popfigure}
\begin{tikzpicture}[
  mindmap, concept color=popblueL,
  level 1 concept/.append style={sibling angle=90, minimum size=2.6cm},
  level 2 concept/.append style={minimum size=2.4cm, font=\small},
  every node/.style={concept, align=center, text width=2.3cm}
]
\node[concept, text width=3.2cm]{\textbf{The General Prologue}\\ Twelve Pilgrims}
  child[concept color=popgreenL]{ node {Military Order} 
    child { node[align=center]{Knight\\ truth, honour\\ ``worthy man''} }
    child { node[align=center]{Squire\\ lively lover\\ ``fresh as May''} }
    child { node[align=center]{Yeoman\\ skilled forester\\ bow \& peacock arrows} }
  }
  child[concept color=popgreenL]{ node {Religious Order (corrupt)}
    child { node[align=center]{Prioress\\ false refinement\\ ``Amor vincit omnia''} }
    child { node[align=center]{Monk\\ worldly hunter\\ ignores the Rule} }
    child { node[align=center]{Friar\\ greedy, merry\\ ``worthy limitour''} }
  }
  child[concept color=popblueL]{ node {Professionals}
    child { node[align=center]{Merchant\\ secretly in debt\\ talks of profit} }
    child { node[align=center]{Clerk\\ true scholar\\ ``gladly lerne, gladly teche''} }
    child { node[align=center]{Sergeant of Law\\ seems busier\\ than he is} }
  }
  child[concept color=popblueL]{ node {Rural / Land Managers}
    child { node[align=center]{Franklin\\ epicurean host\\ ``it snowed of meat''} }
    child { node[align=center]{Reeve\\ feared cheat\\ rich on his lord} }
    child { node[align=center]{Miller\\ golden thumb\\ wart, bagpipes} }
  };
\end{tikzpicture}
\legende{Mind-map of the twelve pilgrims studied in Sections 2 and 3, grouped by social estate, with one key trait and one key quotation or image each.}
\end{popfigure}

\courssub{Synthesis: Idealised and Satirised Portraits}

A complete picture of Chaucer's method emerges when we set the portraits side by side. He uses the same technique everywhere — physical appearance, clothing, speech, habits — but aims it in two opposite directions. In the \emph{idealised} portraits (Knight, Clerk), every detail confirms inner worth: the Clerk's thin horse and threadbare coat prove his devotion to learning. In the \emph{satirised} portraits (Merchant, Sergeant, Reeve, Miller), the details work against the character: the Merchant's splendid hat collapses under the word ``debt''; the Sergeant's learning is deflated by ``seemed busier than he was''; the Reeve's lean body mirrors his lean morality. Between the two poles stands the Franklin, admired and gently mocked at once. This mixture is what makes the General Prologue a \emph{complete portrait of a society}: not a sermon, not a celebration, but a fair, witty and exact mirror.

\begin{methode}[Writing a Comparative Paragraph on Two Pilgrims]
\begin{enumerate}
\item \textbf{Choose a clear point of comparison} (e.g.\ honesty, attitude to money, appearance versus reality, professional ethics).
\item \textbf{State the point} in one topic sentence naming both pilgrims.
\item \textbf{Quote or precisely paraphrase} one detail from each portrait (use Middle English where possible, then gloss).
\item \textbf{Analyse the technique}: explain how the detail (irony, hyperbole, imagery, juxtaposition) creates Chaucer's judgement.
\item \textbf{Link}: conclude by showing what the contrast reveals about Chaucer's view of his society or the estate they represent.
\end{enumerate}
\end{methode}

\begin{attention}[Analyse, Do Not Retell]
The most common examination error is to \emph{retell} the portraits (``The Miller is big, he has a wart, he plays the bagpipes...''). This earns very few marks. Always ask \emph{how} and \emph{why}: how does the wart contribute to the grotesque comic effect? Why does Chaucer give the Miller a ``golden'' thumb rather than simply calling him a thief? Every detail you mention must be tied to a judgement or an effect.
\end{attention}

\begin{exoresolu}[Comparing the Clerk and the Merchant]
\textbf{Question:} Using the General Prologue, show how Chaucer contrasts the Clerk of Oxford and the Merchant to comment on wealth and learning in his society. (One paragraph.)
\tcblower
\textbf{Model paragraph:}\par
Through the paired portraits of the Merchant and the Clerk, Chaucer contrasts hollow wealth with genuine intellectual riches. The Merchant appears splendid — forked beard, Flemish beaver hat, elegant boots — and speaks constantly of his profits; yet Chaucer's quiet revelation that he is secretly in debt exposes this prosperity as performance, so that the pilgrim's fine clothes become a symbol of emptiness. The Clerk is his exact opposite: his horse is ``as lean... as a rake,'' his clothes are threadbare, yet every coin his friends lend him is spent on books, and his speech is ``short and quick and full of high sentence.'' Where the Merchant's appearance flatters a corrupt reality, the Clerk's shabby appearance conceals a rich mind, summed up in the perfectly balanced line ``And gladly wolde he lerne, and gladly teche,'' whose repetition of ``gladly'' conveys pure, disinterested joy in knowledge. By placing the false rich man beside the true poor scholar, Chaucer suggests that in his changing society real worth lies not in the display of money but in the pursuit of learning and moral virtue.
\end{exoresolu}

\begin{exercice}[Practice — Portrait Analysis]
\begin{enumerate}
\item Identify the device Chaucer uses in the line ``It snowed in his house of meat and drink,'' and explain its effect on our judgement of the Franklin.
\item Explain the irony of the Reeve ``lending'' his lord the lord's own goods. What does this reveal about the management of medieval estates?
\item In what sense is the Sergeant of Law similar to the Merchant? Support your answer with one detail from each portrait.
\item Write a comparative paragraph (using the method above) on the Reeve and the Miller as examples of rural corruption.
\end{enumerate}
\end{exercice}

\begin{corrige}[Exercise n° 3]
Both men build a public image of success that conceals an uncomfortable truth. The Merchant talks endlessly of his profits and dresses in expensive Flemish fashion, yet he is secretly in debt; the Sergeant of Law possesses genuine, encyclopaedic legal knowledge, yet Chaucer remarks that he ``seemed busier than he was,'' suggesting that his air of frantic importance is partly an act. In both portraits Chaucer uses the same satiric technique: he allows the character to construct an impressive façade and then punctures it with a single understated detail. The similarity shows Chaucer's wider judgement on the rising professional class — men of real ability who nonetheless feel compelled to exaggerate their success, and whose talents serve their own enrichment rather than the common good.
\end{corrige}

\begin{aretenir}[To Remember]
\begin{itemize}
\item The \cle{Merchant}: splendid appearance, talk of profits — but secretly in \emph{debt} (appearance vs.\ reality).
\item The \cle{Clerk}: thin horse, threadbare coat, books bought with borrowed money — \emph{``gladly wolde he lerne, and gladly teche''}; an idealised portrait.
\item The \cle{Sergeant of Law}: master of cases and property — but only \emph{seems} busier than he is.
\item The \cle{Franklin}: epicurean hospitality — \emph{``it snowed... of meat and drink''} (hyperbole).
\item The \cle{Reeve}: lean, choleric, feared, a carpenter who grows rich by cheating his young lord.
\item The \cle{Miller}: wart, brute strength, bagpipes, and the \emph{``golden thumb''} of the cheating weigher.
\item Chaucer mixes \emph{idealised} and \emph{satirised} portraits to hold a complete, balanced mirror to fourteenth-century society.
\end{itemize}
\end{aretenir}

\coursec{Section 4: A Raisin in the Sun — Hansberry, Context, and Act 1 Scene 1: The Younger Family and the \$10,000 Dream}

Having examined Chaucer's vivid portrait of medieval English society in \emph{The Canterbury Tales}, we now cross centuries and continents to mid-twentieth-century America. Where Chaucer used a pilgrimage frame to reveal the tensions between estate and individual, Lorraine Hansberry uses a cramped Chicago apartment to expose the collision between the American Dream and the reality of racial segregation. Both works, however, share a profound concern: what happens when social structures defer or deny the dreams of ordinary people?

\courssub{Lorraine Hansberry (1930–1965): A Life Shaped by the Struggle for Housing Justice}

Lorraine Vivian Hansberry was born in Chicago on 19 May 1930 into a middle-class, politically active African-American family. Her father, Carl Augustus Hansberry, was a successful real-estate broker and a fierce opponent of racial restrictive covenants — legal clauses that barred Black families from buying homes in white neighbourhoods. In 1937, the Hansberrys purchased a house in the all-white Washington Park subdivision. They were met with violent hostility: a mob surrounded the house, and a concrete block was thrown through a window, narrowly missing the eight-year-old Lorraine. Carl Hansberry fought the ensuing eviction suit all the way to the United States Supreme Court, winning the landmark decision \emph{Hansberry v. Lee} (1940), which ruled that restrictive covenants could not be enforced without a prior court hearing — a crucial, though not final, blow against housing segregation.

\begin{savaistu}[From Courtroom to Stage]
The Hansberry family's legal battle directly inspired \emph{A Raisin in the Sun}. The Younger family's attempt to move into Clybourne Park mirrors the Hansberrys' own experience. Lorraine later wrote: ``My father ... spent a small personal fortune, his considerable talents, and many years of his life fighting, in association with NAACP attorneys, Chicago's 'restrictive covenants' in one of this nation's ugliest ghettos.''
\end{savaistu}

Hansberry attended the University of Wisconsin–Madison but left before graduating to pursue writing in New York. She worked for Paul Robeson's progressive newspaper \emph{Freedom}, where she engaged with the leading Black intellectuals and artists of the day. In 1953 she married Robert Nemiroff, a Jewish songwriter and activist; their marriage ended in divorce in 1964, but they remained close collaborators. Hansberry died of pancreatic cancer on 12 January 1965, aged only 34. At her funeral, James Baldwin said: ``She was a very small, very intense, very bright, very black, very beautiful woman ... She had a great deal to do with the fact that the Civil Rights Movement became a movement.''

\begin{aretenir}[Hansberry — Key Biographical Facts]
\begin{itemize}
\item \textbf{First Black woman} to have a play produced on Broadway (\emph{A Raisin in the Sun}, 1959).
\item Youngest American and first Black playwright to win the \textbf{New York Drama Critics' Circle Award} for Best Play (1959).
\item Daughter of Carl Hansberry, plaintiff in \emph{Hansberry v. Lee} (1940) — a Supreme Court case against racial restrictive covenants.
\item Active in Civil Rights, feminist, and anti-colonial causes; wrote for \emph{Freedom} newspaper.
\item Died at 34; left unfinished works including \emph{Les Blancs} and \emph{The Drinking Gourd}.
\end{itemize}
\end{aretenir}

\courssub{Literary Period and Historical Context: 1950s America}

\emph{A Raisin in the Sun} premiered on Broadway at the Ethel Barrymore Theatre on 11 March 1959. To understand the play, we must situate it in three intersecting historical currents:

\begin{enumerate}
\item \textbf{Racial Segregation and the Colour Line:} Despite the 1954 \emph{Brown v. Board of Education} ruling, \emph{de facto} segregation persisted in housing, employment, and education. The ``colour line'' described by W.E.B. Du Bois remained a brutal reality, especially in Northern cities like Chicago where restrictive covenants and ``redlining'' confined Black families to overcrowded, overpriced ghettos.
\item \textbf{The Great Migration (1916–1970):} Over six million African Americans moved from the rural South to Northern and Western cities. The Youngers are part of this migration: Big Walter and Lena came from the South, seeking freedom and opportunity, only to find a different kind of confinement in Chicago's ``Black Belt'' (Southside).
\item \textbf{The Early Civil Rights Movement:} 1955–1956 Montgomery Bus Boycott (Rosa Parks, Martin Luther King Jr.); 1957 Little Rock Nine; the rise of non-violent direct action. The play captures the mood \emph{before} the major legislative victories of the 1960s — a moment of rising expectation and deep frustration.
\end{enumerate}

\begin{definition}[The Title: ``A Dream Deferred'' — Langston Hughes's \emph{Harlem}]
The play's title comes from the opening lines of Langston Hughes's 1951 poem \emph{Harlem} (also titled \emph{A Dream Deferred}):

\begin{center}
\textit{What happens to a dream deferred?}\\
\textit{Does it dry up}\\
\textit{like a raisin in the sun?}\\
\textit{Or fester like a sore —}\\
\textit{And then run?}\\
\textit{Does it stink like rotten meat?}\\
\textit{Or crust and sugar over —}\\
\textit{like a syrupy sweet?}\\
\textit{Maybe it just sags}\\
\textit{like a heavy load.}\\
\textit{Or does it explode?}
\end{center}

Hughes poses a series of similes for the fate of the Black American dream of equality: drying up, festering, rotting, sagging — or exploding. Hansberry's play dramatises each possibility through the Younger family. The ``raisin in the sun'' is the dream shriveled by the heat of oppression; the ``explosion'' is the potential for revolutionary action.
\end{definition}

\courssub{Setting: The Younger Apartment — A Symbol of Confinement}

The play is set entirely in the Youngers' three-room apartment on Chicago's Southside, ``sometime between World War II and the present'' (i.e., the late 1950s). Hansberry's stage directions are famously detailed — they are not mere furniture lists but a \cle{visual metaphor} for the family's condition.

\begin{experience}[Reading the Stage Directions as Text]
Hansberry opens with a description that does the work of a novel's first paragraph:

\begin{quote}
\textit{The Younger living room would be a comfortable and well-ordered room if it were not for a number of indestructible contradictions to this state of being. Its furnishings are typical and undistinguished and their primary feature now is that they have clearly had to accommodate the living of too many people for too many years — and they are tired.}
\end{quote}

Notice the personification: the room itself is ``tired.'' The ``indestructible contradictions'' are the gap between what the family \emph{aspires} to (comfort, order) and what poverty \emph{forces} upon them. The worn carpet, the single window, the shared bathroom down the hall — every detail signals \cle{material deprivation} and \cle{loss of privacy}.
\end{experience}

\begin{popfigure}
\begin{tikzpicture}[scale=0.7, every node/.style={font=\small}]
% Outer walls
\draw[thick, popdark] (0,0) rectangle (14,10);
% Inner walls
\draw[thick, popdark] (7,0) -- (7,10); % vertical divider
\draw[thick, popdark] (0,6) -- (7,6); % horizontal divider in left section
\draw[thick, popdark] (7,4) -- (14,4); % horizontal divider in right section

% Labels for rooms
\node[align=center, text width=6cm] at (3.5,8) {\textbf{MAMA \& BENEATHA'S BEDROOM}\\(Lena \& Beneatha)};
\node[align=center, text width=6cm] at (3.5,3) {\textbf{WALTER \& RUTH'S BEDROOM}\\(Walter Lee \& Ruth)};
\node[align=center, text width=6cm] at (10.5,7) {\textbf{LIVING ROOM / KITCHEN}\\(Shared space — Travis sleeps on sofa)};
\node[align=center, text width=6cm] at (10.5,2) {\textbf{HALLWAY / SHARED BATHROOM}\\(Used by all families on floor)};

% Furniture symbols
% Mama & Beneatha's room
\draw[fill=popmuted] (1,7.5) rectangle (3,8.5); \node at (2,8) {Bed};
\draw[fill=popmuted] (5,7.5) rectangle (6,8.5); \node at (5.5,8) {Bed};
% Walter & Ruth's room
\draw[fill=popmuted] (1,2.5) rectangle (3,3.5); \node at (2,3) {Bed};
% Living room
\draw[fill=popmuted] (8,5.5) rectangle (11,6.5); \node at (9.5,6) {Sofa (Travis)};
\draw[fill=popmuted] (11.5,5.5) rectangle (13,6.5); \node at (12.25,6) {Chair};
\draw[fill=popmuted] (8,7.5) rectangle (9.5,8.5); \node at (8.75,8) {Table};
\draw[fill=popmuted] (11,7.5) rectangle (12,8.5); \node at (11.5,8) {Chair};
% Kitchen area (implied in living room)
\draw[fill=popmuted] (12,8.5) rectangle (13.5,9.5); \node at (12.75,9) {Stove};
\draw[fill=popmuted] (12,7) rectangle (13.5,7.5); \node at (12.75,7.25) {Sink};

% Window
\draw[thick, popblue] (0,9.5) -- (0,10) node[midway, left=0.2cm, rotate=90] {\textbf{Only Window}};
\draw[thick, popblue] (0,9.5) -- (0.5,9.5);
\draw[thick, popblue] (0,10) -- (0.5,10);

% Plant on windowsill
\node[align=center, popgreen] at (0.5,9.75) {\textbf{Mama's Plant}};

% Door to hallway
\draw[thick, popdark] (13.5,2) -- (14,2);
\node at (13.75,1.5) {\textbf{Exit to Hall}};

% Dimensions note
\node[align=left, popmuted, font=\tiny] at (7,-0.5) {Approx. 3 rooms + kitchenette \\ Shared bathroom down the hall \\ 5 people, 2 bedrooms, 1 window};
\end{tikzpicture}
\legende{Figure 4.1 — Floor plan of the Younger apartment. The single window, shared bathroom, and Travis's sofa-bed illustrate the physical confinement that mirrors the family's socio-economic entrapment.}
\end{popfigure}

\begin{aretenir}[Setting as Symbol]
\begin{itemize}
\item \textbf{Three rooms for five people:} No privacy; intimacy forced by poverty.
\item \textbf{One window:} Limited light/air — the plant on the sill is the only living green thing.
\item \textbf{Shared bathroom in hallway:} Dignity compromised; the family is not fully ``at home'' even in their own flat.
\item \textbf{Tired furniture:} Material history of deferred dreams — each scratch tells a story of making do.
\item \textbf{Kitchenette in living room:} No separate domestic space; work and rest collapse into one.
\end{itemize}
\end{aretenir}

\courssub{Act 1 Scene 1: Morning Routine — The World of the Play}

The scene opens early on a Friday morning. The stage directions establish a \cle{rhythm of survival}: the alarm rings, Ruth wakes Travis, Walter stumbles in, the bathroom is contested. This is \emph{realism} at its most deliberate — Hansberry wants the audience to \emph{feel} the daily friction of overcrowding before the plot (the cheque) arrives.

\begin{exemplebox}[Opening Stage Business — What It Reveals]
\begin{center}
\begin{tabularx}{\linewidth}{|l|X|}\hline
\rowcolor{popblueL}
\textbf{Action} & \textbf{Character / Theme Revealed} \\ \hline
Ruth wakes Travis gently but firmly & Ruth as nurturer, exhausted but functional \\ \hline
Walter asks ``Check coming today?'' before greeting anyone & Money = survival; his identity fused with provider role \\ \hline
Ruth: ``Walter, leave me alone!'' & Marital strain; her weariness vs. his restless energy \\ \hline
Travis asks for 50 cents for school; Ruth says no; Walter gives him \$1 & Walter's need to be ``the man'' (provider) vs. Ruth's realism \\ \hline
Beneatha enters, ``fresh and vibrant'' & Youth, education, aspiration — contrast to parents' weariness \\ \hline
Mama enters, tends to plant & Life-giving force; the plant = her dream of a garden/house \\ \hline
\end{tabularx}
\end{center}
\end{exemplebox}

\courssub{The \$10,000 Life-Insurance Cheque: Catalyst and Symbol}

Big Walter (Lena's husband, father of Walter Lee and Beneatha) has recently died. The \$10,000 cheque (equivalent to roughly \$105,000 today) is the \cle{inciting incident}. It is not merely money; it is the \textbf{monetised value of a Black man's labour and life} in a system that paid him less than his worth.

\begin{attention}[The Cheque Is Not Just Money — Exam Warning]
In essays, never write ``the \$10,000 cheque gives them a chance.'' Write: ``The \$10,000 cheque \emph{represents Big Walter's sacrificed body and deferred dreams}; it is the material residue of a life spent serving a system that denied him the very things his family now hopes to buy.'' The cheque is \textbf{symbolic capital} — treat it as such.
\end{attention}

\begin{definition}[Big Walter's Ghost]
Though never seen on stage, Big Walter haunts the play. Mama says: ``Big Walter used to say, he'd get right wet in the eyes sometimes, lean his head back with the water standing in his eyes and say, `Seem like God didn't see fit to give the black man nothing but dreams — but He did give us children to make them dreams seem worth while.''' Big Walter's dream was a house; his death makes the down-payment possible. The cheque is his \emph{legacy} — and the family's fight over it is a fight over \emph{whose interpretation of his legacy prevails}.
\end{definition}

\courssub{Conflicting Dreams: Three Visions for the Money}

The cheque exposes the fault lines in the family. Each adult has a distinct dream, rooted in their generation, gender, and self-conception.

\begin{center}
\begin{tabularx}{\linewidth}{|l|X|X|X|}\hline
\rowcolor{popblueL}
\textbf{Character} & \textbf{Dream} & \textbf{What It Represents} & \textbf{Key Quotation (Learn One Per Character)} \\ \hline
\textbf{Walter Lee} & Invest in a liquor store with Willy Harris and Bobo & \textbf{Manhood, autonomy, economic power.} He wants to stop being a chauffeur — ``driving a man around in his limousine'' — and become an owner. & ``I want so many things that they are driving me kind of crazy... Mama — look at me.'' (Act 1 Sc 1) \\ \hline
\textbf{Beneatha} & Medical school tuition & \textbf{Intellectual fulfilment, professional identity, race pride.} She wants to be a doctor — a healer — and explores African identity (Asagai) vs. assimilation (George Murchison). & ``I am going to be a doctor... everybody thinks it's funny but I don't.'' (Act 1 Sc 1) \\ \hline
\textbf{Mama (Lena)} & A house with a garden in a decent neighbourhood (Clybourne Park) & \textbf{Stability, dignity, rootedness, family continuity.} The house is the \emph{concrete form} of Big Walter's dream and her own. & ``We ain't never been that poor... We ain't never been that — dead inside.'' (Act 1 Sc 1) \\ \hline
\end{tabularx}
\end{center}

\begin{exoresolu}[Worked Example: Analysing Walter's Dream in Act 1 Scene 1]
\textbf{Question:} How does Hansberry present Walter Lee's frustration in Act 1 Scene 1?
\tcblower
\textbf{Model Response:}
Hansberry presents Walter's frustration through a combination of \textbf{realistic dialogue}, \textbf{physical restlessness}, and \textbf{juxtaposition with other characters}.

\begin{enumerate}
\item \textbf{Dialogue:} His first words — ``Check coming today?'' — establish his fixation. His monologue to Ruth (``I'm thirty-five years old; I been married eleven years and I got a boy who sleeps in the living room...'') uses \textbf{anaphora} (``I got... I got...'') to accumulate grievances. The repetition of ``Man say to his woman...'' reveals his gendered definition of success: a man provides; a woman supports.
\item \textbf{Stage Directions:} ``He leans over the back of the chair... looking at the check'' — his body language is predatory, hungry. Contrast with Ruth's ``weary'' movements.
\item \textbf{Juxtaposition:} Beneatha's intellectual ambition (medical school) and Mama's grounded pragmatism (house) throw Walter's \emph{speculative} capitalism (liquor store) into relief. His dream is the most \emph{assimilated} to white capitalist values — he wants to be a ``businessman'' like the white men he chauffeurs.
\item \textbf{Symbolism:} The eggs Ruth cooks — ``scrambled'' — mirror his scrambled priorities. He refuses to eat them, rejecting her nurture because it doesn't match his self-image.
\end{enumerate}

\textbf{Examiner's Comment:} This response integrates \textbf{language} (anaphora), \textbf{stagecraft} (directions, juxtaposition), and \textbf{symbolism} (eggs) — the three pillars of drama analysis. Always quote the stage directions; they are part of the text.
\end{exoresolu}

\courssub{Walter's Crisis of Manhood and Ruth's Silent Burden}

Walter's definition of manhood is \textbf{externally validated}: he needs the world (white society, his family) to see him as a provider. His job as a chauffeur is emasculating in his eyes because it is \emph{service} — he opens doors for a white man who doesn't know his name. When he says, ``I want to hang some real pearls round my wife's neck,'' he equates love with \textbf{material display}.

Ruth, by contrast, embodies \cle{endurance}. Her pregnancy (revealed at the end of Act 1 Scene 2, but foreshadowed by her fainting and ``weariness'' in Scene 1) adds a biological urgency to the need for space. She considers abortion — a radical, dangerous choice in 1959 — because she cannot bear to bring another child into this apartment. Her dream is \textbf{not articulated as a wish}; it is \textbf{survival}.

\begin{methode}[Reading Gender in Act 1 Scene 1 — Step by Step]
\begin{enumerate}
\item \textbf{Identify each character's \emph{public} dream} (what they say they want).
\item \textbf{Identify the \emph{private} fear} beneath it (what they are running from).
\item \textbf{Map the power dynamic:} Who controls the money? Who controls the \emph{narrative} of the family?
\item \textbf{Track the \emph{body}:} Hansberry writes the body into the text — fainting, pregnancy, exhaustion, restless pacing. The body is the site where social pressure becomes physical reality.
\item \textbf{Note the \emph{generational} axis:} Mama (past/sacrifice) — Walter/Ruth (present/struggle) — Beneatha/Travis (future/possibility).
\end{enumerate}
\end{methode}

\courssub{Dramatic Techniques in Act 1 Scene 1}

\begin{propriete}[Realism as Resistance]
Hansberry employs \textbf{psychological realism} (detailed setting, overlapping dialogue, mundane routine) to \textbf{centre Black domestic life} on the American stage — a radical act in 1959. The ``well-made play'' structure (exposition → inciting incident → rising action) is fused with \textbf{symbolic realism}: objects (the plant, the cheque, the eggs, the window) carry thematic weight beyond their literal function.
\end{propriete}

\begin{center}
\begin{tabularx}{\linewidth}{|l|X|X|}\hline
\rowcolor{popblueL}
\textbf{Technique} & \textbf{Example in Act 1 Sc 1} & \textbf{Effect} \\ \hline
\textbf{Stage Directions as Characterisation} & ``Ruth is about thirty... but she looks older... disappointment has already begun to hang in her face.'' & Shows the \emph{physical toll} of poverty before a word is spoken. \\ \hline
\textbf{Symbolic Props} & Mama's plant (life in hostility); the insurance cheque (Big Walter's life); the single window (limited horizon). & Objects become \emph{objective correlatives} for internal states. \\ \hline
\textbf{Realistic Dialogue / Vernacular} & ``Ain't nobody said nothing 'bout no money!'' / ``I don't want to hear no more about it.'' & Authenticates the Black working-class voice; resists minstrel stereotypes. \\ \hline
\textbf{Off-Stage References} & The bathroom down the hall; the white neighbours' ``welcoming committee'' (later); Big Walter's death. & Expands the world beyond the apartment — the \emph{social context} presses in. \\ \hline
\textbf{Foreshadowing} & Walter: ``Something's eating you up... like a crazy man.'' (Irony: he describes himself.) & Prepares the audience for the explosion in Act 2. \\ \hline
\end{tabularx}
\end{center}

\begin{savaistu}[The Plant — Hansberry's Central Symbol]
Mama's plant is the only living thing in the apartment besides the people. She tends it despite ``ain't never had enough sunshine.'' It represents:
\begin{itemize}
\item Her dream of a garden (a house with earth).
\item Her nurturing instinct — she refuses to let it die, just as she refuses to let the family disintegrate.
\item The \textbf{possibility of growth in hostile conditions} — the play's central question.
\end{itemize}
In the final scene, she takes the plant to the new house: the dream, battered but alive, moves with them.
\end{savaistu}

\courssub{The Younger Family Tree: Dreams and Conflicts}

\begin{popfigure}
\begin{tikzpicture}[scale=0.65, every node/.style={font=\small, align=center}]
% Nodes
\node[draw, rounded corners, fill=popblueL, thick, align=center] (bigw) at (0,6){\textbf{BIG WALTER}\\(Deceased)\\Legacy: \$10,000 insurance};
\node[draw, rounded corners, fill=popgreenL, thick, align=center] (mama) at (0,3){\textbf{MAMA (LENA)}\\\textit{Dream: House with garden}\\\textit{Conflict: Steward of legacy}};
\node[draw, rounded corners, fill=poporange, thick, align=center] (walter) at (-4,0){\textbf{WALTER LEE}\\\textit{Dream: Liquor store investment}\\\textit{Conflict: Manhood vs. provider}};
\node[draw, rounded corners, fill=poppink, thick, align=center] (ruth) at (0,0){\textbf{RUTH}\\\textit{Dream: Space, dignity, 2nd child?}\\\textit{Conflict: Weariness vs. hope}};
\node[draw, rounded corners, fill=poppurple, thick, align=center] (beneatha) at (4,0){\textbf{BENEATHA}\\\textit{Dream: Medical school}\\\textit{Conflict: Identity (African vs. assimilation)}};
\node[draw, rounded corners, fill=popgold, thick, align=center] (travis) at (0,-2.5){\textbf{TRAVIS (10)}\\\textit{Dream: ``A real house''}\\\textit{Conflict: Childhood in confinement}};

% Edges
\draw[thick, popdark] (bigw) -- (mama) node[midway, left] {Widow};
\draw[thick, popdark] (mama) -- (walter) node[midway, left] {Son};
\draw[thick, popdark] (mama) -- (beneatha) node[midway, right] {Daughter};
\draw[thick, popdark] (walter) -- (ruth) node[midway, below] {Wife};
\draw[thick, popdark] (walter) -- (travis) node[midway, left] {Son};
\draw[thick, popdark] (ruth) -- (travis) node[midway, right] {Son};

% Dream/Conflict labels on edges
\node[popmuted, font=\tiny, align=center] at (-2,4.5) {Big Walter's dream\\\textit{becomes her burden}};
\node[popmuted, font=\tiny, align=center] at (-2,1.5) {Walter needs\\\textit{Mama's trust}};
\node[popmuted, font=\tiny, align=center] at (2,1.5) {Beneatha needs\\\textit{Mama's support}};
\node[popmuted, font=\tiny, align=center] at (-2,-1.25) {Travis =\\\textit{Walter's mirror}};
\node[popmuted, font=\tiny, align=center] at (2,-1.25) {Travis =\\\textit{Beneatha's future}};

% Asagai / George (off to side)
\node[draw, dashed, fill=popmuted, font=\tiny, align=center] (asagai) at (7,1){ASAGAI\\\textit{Nigerian student}\\\textit{Black pride / Pan-Africanism}};
\node[draw, dashed, fill=popmuted, font=\tiny, align=center] (george) at (7,-1){GEORGE MURCHISON\\\textit{Wealthy Black student}\\\textit{Assimilationist}};
\draw[dashed, ->] (beneatha) -- (asagai) node[midway, above, font=\tiny] {Ideological pull};
\draw[dashed, ->] (beneatha) -- (george) node[midway, below, font=\tiny] {Social pressure};
\end{tikzpicture}
\legende{Figure 4.2 — The Younger family tree with each member's dream and central conflict. Dashed boxes represent off-stage ideological forces acting on Beneatha.}
\end{popfigure}

\courssub{Worked Exercise: Close Reading of a Key Exchange}

\begin{exoresolu}[Close Reading: Walter and Ruth on «Manhood» (Act 1 Scene 1)]
\textbf{Extract:}
\begin{quote}
WALTER: \textit{... Man say to his woman: I got me a dream. His woman say: Eat your eggs. Man say: I got to take hold of this here world, baby! And a woman will say: Eat your eggs and go to work.}
\end{quote}

\textbf{Question:} Analyse how Hansberry uses this exchange to reveal Walter's worldview and the gender dynamic in the marriage.
\tcblower
\textbf{Detailed Solution:}

\begin{enumerate}
\item \textbf{Structure:} The exchange is a \textbf{stichomythia}-like pattern — «Man say / Woman say» — reducing complex individuals to \textbf{gendered archetypes}. Walter speaks in \textbf{universalising generalisations} («Man», «His woman», «a woman»); he avoids «I» and «you», turning his private frustration into a \textbf{manifesto}.
\item \textbf{Metaphor:} «Take hold of this here world» — the world is an object to be seized, possessed. This is \textbf{capitalist conquest language}. The eggs (domestic, nurturing, feminine) are the \textbf{antithesis} of his dream.
\item \textbf{Irony:} Ruth \emph{does} work (she's a domestic worker). She \emph{does} scramble eggs (labour). Walter's fantasy of the «woman» who only says «Eat your eggs» erases Ruth's actual economic contribution. He \textbf{misrecognises} her.
\item \textbf{Power:} Walter positions himself as the \textbf{visionary} and Ruth as the \textbf{anchor} — but the anchor is what keeps the ship from drifting. The audience sees Ruth's \textbf{clear-eyed realism} (she knows the liquor store is risky) vs. Walter's \textbf{romantic desperation}.
\item \textbf{Thematic Link:} This gender split mirrors the \textbf{generational split}: Mama (like Ruth) wants security/roots; Walter (like Beneatha, but differently) wants \textbf{transformation/escape}. The play refuses to privilege either entirely.
\end{enumerate}

\textbf{Exam Tip:} When quoting dialogue, always note \textbf{who speaks, to whom, and what is happening on stage simultaneously} (e.g., Ruth cooking, Walter pacing). Drama is \emph{embodied} speech.
\end{exoresolu}

\courssub{Consolidation: Key Takeaways for Act 1 Scene 1}

\begin{aretenir}[Act 1 Scene 1 — Revision Checklist]
\begin{itemize}
\item \textbf{Exposition done through routine:} The morning chaos \emph{is} the exposition — we learn who these people are by watching them survive.
\item \textbf{The cheque arrives off-stage:} Its \emph{anticipation} structures the scene; the physical envelope appears only at the end of Scene 1 / start of Scene 2.
\item \textbf{Three dreams, one pot of money:} The structural tension of the play is established: \textbf{zero-sum game} (or so it seems).
\item \textbf{Walter's soliloquy to Ruth} is the \textbf{aria} of the scene — learn 2–3 lines for quotation use.
\item \textbf{Mama's plant} enters and exits the scene with her — it is her \textbf{stage property} in every sense.
\item \textbf{The window} frames the only «outside» — the Southside street, the grey sky, the world they want to leave.
\end{itemize}
\end{aretenir}

\courssub{Examination Practice}

\begin{exercice}[Essay-Style Questions — Act 1 Scene 1 Focus]
\begin{enumerate}
\item ``Hansberry uses the setting of the Younger apartment as a metaphor for the condition of Black America in the 1950s.'' Discuss with close reference to Act 1 Scene 1.
\item Analyse the presentation of Walter Lee Younger in Act 1 Scene 1. How does Hansberry elicit both sympathy and criticism for him?
\item ``The \$10,000 cheque is the most important character in the play.'' Evaluate this view based on Act 1 Scene 1.
\item Explore the dramatic significance of the different dreams held by Mama, Walter, and Beneatha in Act 1 Scene 1.
\item How does Hansberry use stage directions and off-stage references to establish the social context of the play before the cheque arrives?
\end{enumerate}
\end{exercice}

\begin{corrige}[Exercise 1 — Setting as Metaphor (Outline)]
\textbf{Thesis:} The apartment is not merely a backdrop; it is a \textbf{materialisation of systemic racism} — its cramped rooms, shared facilities, and single window encode the history of segregation.

\textbf{Paragraph Plan:}
\begin{enumerate}
\item \textbf{Physical details as social facts:} Stage directions — «tired» furniture, «indestructible contradictions» — link domestic wear to structural inequality.
\item \textbf{Spatial politics:} 5 people / 3 rooms / 1 bathroom = \textbf{enforced intimacy} = loss of autonomy. Compare to white suburban homes (Clybourne Park) described later.
\item \textbf{The window / plant axis:} The only natural light falls on Mama's plant — the \textbf{only growth} in the space. The window frames the ghetto; the plant dreams of a garden.
\item \textbf{The kitchenette in the living room:} Collapse of public/private spheres — the family cannot ``perform'' middle-class respectability (separate dining, parlour) because the architecture denies it.
\item \textbf{Sound from off-stage:} The bathroom queue, the street noise — the apartment is \textbf{porous} to the community, not a sanctuary.
\item \textbf{Conclusion:} The setting makes the \textbf{personal political}. Every domestic friction (eggs, bathroom, sleep) is a microcosm of the macro-friction: a people confined.
\end{enumerate}
\end{corrige}

\begin{attention}[Common Exam Pitfalls — Act 1 Scene 1]
\begin{itemize}
\item \textbf{Summarising instead of analysing:} ``Walter wants a liquor store, Beneatha wants medical school, Mama wants a house.'' → \textbf{No marks.} You must analyse \textbf{how} Hansberry presents these dreams (language, structure, staging).
\item \textbf{Ignoring Ruth:} She has the fewest lines but carries the \textbf{emotional reality} of the play. Always discuss her.
\item \textbf{Forgetting the historical context:} An essay on \emph{A Raisin in the Sun} without mentioning \textbf{segregation, the Great Migration, or the Civil Rights Movement} is incomplete.
\item \textbf{Misreading the title:} The raisin \emph{dries up} — but the play asks whether it \emph{explodes}. Don't say the play is ``about a raisin.'' Say: ``The play \emph{tests} Hughes's metaphors.''
\item \textbf{Quoting the film, not the play:} The 1961 film adds scenes (the bar, the street). \textbf{Examine the text as printed.} Stage directions are your best evidence.
\end{itemize}
\end{attention}

\courssub{Transition to Act 1 Scene 2}

Act 1 Scene 1 ends with the arrival of the cheque and Mama's terrifying realisation: ``They put the check in my hand... and I stood there... and I didn't know whether to laugh or cry.'' The \textbf{object} is now in the \textbf{subject's} hand — but the \textbf{subject} (Mama) is not the one who \textbf{expected} to hold it (Walter). Scene 2 will explode the generational and gender tensions further, introducing Beneatha's suitors (George Murchison and Joseph Asagai) and the \textbf{ideological debate} over what it means to be Black, modern, and ambitious in 1959 America. We turn to that now.

\coursec{Section 5: A Raisin in the Sun — Act 1 Scene 2 and Act 2 Scene 1: Gender, Identity, the American Dream and Cultural Heritage}

\courssub{Transition from Section 4: The Dream Deferred Takes Shape}

In Section 4 we left the Younger family in their cramped Chicago apartment at the close of Act 1 Scene 1. The \$10,000 insurance cheque — the material embodiment of Big Walter's life labour — has not yet arrived, but it already dominates every conversation. Walter Lee is desperate to invest in a liquor store; Beneatha needs money for medical school; Ruth is pregnant and terrified; and Mama Lena clings to her late husband's dream of a house with a garden. The stage is set for a collision between \emph{what the money can buy} and \emph{what the family stands for}. In this section we follow the action through Act 1 Scene 2 and Act 2 Scene 1, where the cheque arrives, identities are tested, and the American Dream is dissected along lines of gender, race, and cultural memory.

\begin{savaistu}[Lorraine Hansberry and the Title]
The play's title comes from Langston Hughes's poem \emph{Harlem} (1951): ``What happens to a dream deferred? / Does it dry up / like a raisin in the sun? / Or fester like a sore — / And then run? / Does it stink like rotten meat? / Or crust and sugar over — / like a syrupy sweet? / Maybe it just sags / like a heavy load. / Or does it explode?'' Hansberry uses the poem as an epigraph; every character in the play embodies one of Hughes's possible fates for a deferred dream.
\end{savaistu}

\courssub{Act 1 Scene 2: The Cheque Arrives — Money, Morality, and Motherhood}

The scene opens on a Saturday morning, the day the cheque is due. The atmosphere is tense. The first dramatic beat is the arrival of the cheque itself — a simple prop that instantly restructures power in the household.

\begin{exemplebox}[Stage Direction as Character Revelation]
\textbf{Text:} \emph{``MAMA ... She opens the envelope and takes out the check. She looks at it. She looks at the window. She looks at the check again. She is afraid to touch it.''} (Act 1 Scene 2)
\textbf{Analysis:} Mama's hesitation shows that the money represents her husband's \emph{life}, not just currency. She treats it with reverence; Walter treats it as a tool. This contrast — \emph{sacred trust} versus \emph{instrumental capital} — drives the scene.
\end{exemplebox}

\textbf{Mama's Religious Values versus Walter's Materialism.}\par
Mama's first instinct is to give thanks to God: ``Thank you, Lord.'' For her, the money is a blessing to be used \emph{responsibly} — for a house, for Beneatha's education, for the family's \emph{dignity}. Walter, however, sees only the liquor store investment. He says, ``Money is life.'' Mama's retort is the play's moral anchor: ``Once upon a time freedom used to be life — now it's money. I guess the world really do change.''

\begin{definition}[Materialism (in the play)]
The belief that financial acquisition is the primary measure of success, manhood, and freedom. Walter Lee equates the \$10,000 with his ability to be a ``man'' — a provider, a decision-maker, a patriarch. His materialism is not simple greed; it is a response to systemic emasculation.
\end{definition}

\textbf{Ruth's Pregnancy and the Question of Abortion.}\par
Ruth's discovery that she is two months pregnant introduces a second crisis. In 1959, abortion is illegal and dangerous. Ruth has already made a \$5 down payment to ``a woman'' — a back-alley provider. This reveals her desperation: another child means more poverty, more crowding, more strain on a marriage already fracturing. Mama's reaction — ``When the world gets ugly enough — a woman will do anything for her family. The part that's already living'' — shows her pragmatic, life-affirming matriarchy. She does not judge; she acts to preserve the family.

\begin{attention}[Examiner's Trap: Reducing Ruth to a Passive Victim]
Do \emph{not} write that Ruth ``has no agency''. Her decision to seek an abortion (however dangerous) is an act of agency born of despair. Her later decision to keep the baby, influenced by Mama's hope for a house, is also agency. Analyse the \emph{social conditions} that constrain her choices: poverty, racism, lack of reproductive healthcare, and a husband who dismisses her fears.
\end{attention}

\courssub{Gender and Identity: ``Black Women Don't Support Their Men''}

One of the play's most quoted and misunderstood lines occurs when Walter explodes at Ruth: ``That is just what is wrong with the colored woman in this world ... Don't understand about building their men up and making 'em feel like they somebody.''

\begin{aretenir}[Walter's Gender Grievance — Contextualise, Don't Just Quote]
\begin{itemize}
\item Walter feels emasculated by racism (he is a chauffeur, a ``boy'' to white men) and by poverty (he cannot provide).
\item He projects his need for validation onto Black women, demanding emotional labour they are too exhausted to give.
\item Ruth, Mama, and Beneatha \emph{do} support him — but they refuse to \emph{enable} his delusions (the liquor store) or his self-pity.
\item The play critiques \emph{both} Walter's sexism \emph{and} the society that makes his masculinity fragile.
\end{itemize}
\end{aretenir}

\begin{definition}[Matriarch]
A woman who heads a family or tribe, often holding moral, economic, and emotional authority. Mama Lena Younger is the play's matriarch: she holds the insurance cheque, makes the decision to buy the house, and ultimately entrusts Walter with the remaining money — a test of his manhood. Her power is not tyrannical; it is \emph{stewardship}.
\end{definition}

\textbf{Mama as Matriarch: Strength Rooted in Love.}\par
Mama's authority comes from her role as keeper of the family's history and values. She reminds Walter: ``I come from five generations of people who was slaves and sharecroppers — but ain't nobody in my family never let nobody pay 'em no money that was a way of telling us we wasn't fit to walk the earth.'' This speech (Act 1 Scene 2) links personal dignity to collective memory. Mama's matriarchy is \emph{intergenerational} — she answers to Big Walter's ghost and to her unborn grandchild.

\courssub{Beneatha's Search for Identity: Hobbies, Hair, and Two Suitors}

Beneatha Younger is the play's most explicit explorer of \emph{identity}. She is 20, intellectual, and determined to be a doctor — a profession almost closed to Black women in 1959.

\begin{exemplebox}[Beneatha's ``Expensive Hobbies'' as Identity Work]
\begin{itemize}
\item \textbf{Guitar lessons, horseback riding, photography, acting classes} — the family mocks her for ``flitting.''
\item \textbf{Function:} Each hobby is an attempt to \emph{define herself} beyond ``Walter's sister'' or ``Mama's daughter'' or ``a Black girl from the South Side.''
\item \textbf{Critique:} Walter and Ruth see the hobbies as wasteful; Mama sees them as ``expressing herself.'' The play validates Beneatha's search even as it gently satirises her pretensions (e.g., her mangled Nigerian pronunciation).
\end{itemize}
\end{exemplebox}

\textbf{Rejection of Assimilation: George Murchison versus Joseph Asagai.}\par
Beneatha's two suitors embody two paths for the Black middle class.

\begin{center}
\begin{tabularx}{\linewidth}{|l|X|X|}
\hline
\rowcolor{popblueL}
\textbf{Aspect} & \textbf{George Murchison} & \textbf{Joseph Asagai} \\
\hline
\textbf{Class} & Wealthy, assimilated Black bourgeoisie & Nigerian student, intellectual, anti-colonial \\
\hline
\textbf{Attitude to Heritage} & ``Let's face it, baby, your heritage is nothing but a bunch of raggedy-assed spirituals and some grass huts!'' & ``You wear it well ... very well ... mutilated hair and all.'' (teasing but affirming) \\
\hline
\textbf{View of Beneatha} & Wants a pretty, quiet wife; dismisses her ideas & Challenges her intellectually; calls her \emph{Alaiyo} \\
\hline
\textbf{Symbolises} & \cle{Assimilation} — adopting white middle-class norms, rejecting African roots & \cle{Cultural Recovery} — connecting to African heritage as living tradition \\
\hline
\end{tabularx}
\end{center}

\begin{definition}[Assimilation (in the play)]
The process by which a minority individual or group abandons its distinct cultural identity — language, customs, values, aesthetics — to adopt those of the dominant (white American) culture, often in exchange for social acceptance or economic advancement. George Murchison embodies this: he wears ``white shoes,'' reads ``white books,'' and mocks African heritage as primitive.
\end{definition}

\courssub{Act 2 Scene 1: Natural Hair, Nigerian Dress, and ``Alaiyo''}

Act 2 Scene 1 opens with a striking visual transformation: Beneatha has cut off her straightened hair and wears a Nigerian robe and headdress, dancing to African music. This is \emph{performance as politics}.

\begin{experience}[Close Reading: The Hair Scene]
\textbf{Stage Direction:} \emph{``BENEATHA ... Her hair is cut close to her head in a natural style. She is wearing a Nigerian robe ... She turns off the phonograph. She looks at herself in the mirror. She frowns. She pats her hair. She smiles.''}
\textbf{Analysis:}
\begin{enumerate}
\item \textbf{The mirror:} Self-scrutiny, self-creation. She is \emph{becoming}.
\item \textbf{The frown then smile:} The transition is not instant confidence; it is a deliberate political act.
\item \textbf{Walter's entrance (drunk):} He joins the dance, mimicking African movements — a moment of \emph{reconnection} with his own African roots, however performative.
\item \textbf{George's entrance:} He is horrified: ``You look eccentric.'' Beneatha's retort: ``I look like \emph{me}!''
\end{enumerate}
\end{experience}

\textbf{Asagai's Gift and the Nickname ``Alaiyo''.}\par
Asagai gives Beneatha traditional Nigerian records and the robe. He nicknames her \emph{Alaiyo} — Yoruba for ``One for Whom Bread — Food — Is Not Enough.''

\begin{aretenir}[``Alaiyo'' — The Play's Philosophical Centre]
\begin{itemize}
\item \emph{Bread} = material survival, money, the liquor store, the cheque.
\item \emph{Not Enough} = the human need for meaning, beauty, heritage, intellectual freedom, dignity.
\item Beneatha \emph{is} Alaiyo: she wants to be a doctor \emph{and} know who she is.
\item Walter \emph{wants} to be Alaiyo but is trapped in ``bread'' thinking.
\item Mama \emph{knows} bread is not enough — that is why she buys the house (dignity) not the liquor store (money).
\end{itemize}
\end{aretenir}

\courssub{The American Dream Debated: Money versus Dignity}

The central ideological clash crystallises in Act 1 Scene 2 and Act 2 Scene 1.

\begin{propriete}[Walter's Definition: ``Money is Life'']
\textbf{Statement:} For Walter Lee Younger, the American Dream is \emph{exclusively} economic. He believes:
\begin{itemize}
\item Manhood = Providing financially.
\item Freedom = Capital ownership (the liquor store).
\item Success = Moving from ``employee'' to ``employer.''
\end{itemize}
\textbf{Origin:} This is not innate greed. It is the internalisation of capitalist values in a society that denies Black men other avenues to respect. He says: ``I want so many things that they are driving me kind of crazy.''
\textbf{Examinable Note:} You must analyse the \emph{social determinism} behind Walter's view — not just moralise it.
\end{propriete}

\begin{propriete}[Mama's Definition: Freedom, Dignity, Family]
\textbf{Statement:} For Mama, the American Dream is \emph{homeownership as citizenship}. She believes:
\begin{itemize}
\item A house in Clybourne Park = a stake in the nation, a legacy for Travis.
\item Dignity > Money: ``We ain't never been that dead inside.'' (Act 1 Scene 2)
\item Family unity is the ultimate wealth.
\end{itemize}
\textbf{Origin:} Rooted in the Black Southern migrant experience: land = independence. Big Walter worked himself to death for this dream.
\textbf{Examinable Note:} Mama's dream is \emph{collective}; Walter's is \emph{individualistic} (at first).
\end{propriete}

\begin{exoresolu}[Essay Question: ``Money is life.'' How far does the play support Walter's view?]
\textbf{Plan:}
\begin{enumerate}
\item \textbf{Introduction:} Quote the line; locate it in Act 1 Scene 2. State thesis: the play \emph{exposes} the view as understandable but ultimately inadequate; true life requires dignity, heritage, and community.
\item \textbf{Paragraph 1 — Walter's Perspective:} Sympathetic analysis. Racism (chauffeur job), poverty (apartment), emasculation (Mama holds the purse). The liquor store as \emph{agency}. Use Ruth's line: ``He needs something — something I can't give him no more.''
\item \textbf{Paragraph 2 — Mama's Counter-View:} ``Once upon a time freedom used to be life.'' The house as \emph{concrete} freedom. Her refusal to fund the liquor store on moral grounds (``Christian woman''). The Clybourne Park decision as \emph{resistance}.
\item \textbf{Paragraph 3 — Beneatha/Asagai Synthesis:} ``Alaiyo'' — bread is not enough. Intellectual/spiritual fulfilment. Asagai's vision: ``There is only one large circle ... and we are all in it.'' The dream expands beyond America.
\item \textbf{Paragraph 4 — The Resolution (preview):} Walter's final rejection of Lindner's money proves he learns \emph{dignity > money}. But the play does not dismiss his pain.
\item \textbf{Conclusion:} The play validates the \emph{hunger} behind ``money is life'' but affirms that life \emph{is} dignity, heritage, and love.
\end{enumerate}
\textbf{Model Opening Paragraph:}
``When Walter Lee Younger declares `Money is life' in Act 1 Scene 2, he articulates the desperate logic of a Black man in 1950s America for whom every door marked `Manhood' requires a key called `Capital'. Hansberry does not dismiss this view; she dramatises its origins in systemic oppression. Yet through Mama's matriarchal wisdom, Beneatha's intellectual quest, and Asagai's Pan-African vision, the play ultimately argues that a life reduced to financial transaction is a life `dead inside' — and that the true American Dream, for the Youngers, is the courage to claim dignity in a society that denies it.''
\end{exoresolu}

\courssub{Mama's Decision: Clybourne Park — Hope and Danger}

The pivotal action of Act 2 Scene 1 is Mama's revelation: she has put a \$3,500 down payment on a house in \emph{Clybourne Park}, an all-white neighbourhood.

\begin{methode}[Analysing a Symbolic Setting: Clybourne Park]
\begin{enumerate}
\item \textbf{Identify the literal level:} A cheaper, larger house with a garden — Mama's and Big Walter's dream.
\item \textbf{Identify the social level:} An all-white suburb; the Youngers will be the first Black family. This is \emph{integration} by force of purchase.
\item \textbf{Identify the thematic level:} The house = \emph{stake in the American promise}. It forces the family to confront white supremacy (Lindner's visit comes next) and their own fears.
\item \textbf{Contrast with the alternative:} Staying in the apartment = stagnation; the liquor store = risky capitalism; Clybourne Park = \emph{moral courage}.
\item \textbf{Track consequences:} The decision splits the family (Walter's despair), unites them (against Lindner), and sets up the climax.
\end{enumerate}
\end{methode}

\begin{attention}[Common Mistake: Seeing Clybourne Park as a ``Happy Ending'']
The house purchase is \emph{not} a resolution. It is a \emph{provocation}. The family will face hostility, financial strain (Mama spends most of the money), and Walter's betrayal (losing the rest). The play ends with them \emph{leaving} for Clybourne Park — the struggle continues. Do not write ``they achieve the American Dream.'' Write ``they \emph{choose} the struggle for dignity over the safety of the ghetto.''
\end{attention}

\courssub{Mama Entrusts Walter: The Test of Manhood}

In a stunning act of trust, Mama gives Walter the remaining \$6,500: \$3,000 for Beneatha's medical school (in a savings account), \$3,500 for him to ``be the head of this family.''

\begin{exemplebox}[Dramatic Irony and Character Test]
\begin{itemize}
\item \textbf{Mama's intention:} To heal Walter's wounded manhood by giving him \emph{responsibility}.
\item \textbf{Walter's perception:} Validation — ``You trust me like that, Mama?''
\item \textbf{Audience knowledge:} We suspect Walter will give it all to Willy Harris for the liquor store (he has already discussed it with Bobo).
\item \textbf{Thematic weight:} This is the play's central \emph{test}. Can Walter convert \emph{money} into \emph{manhood}? The answer (in Act 2 Scene 3) is initially no — but ultimately yes, through a \emph{different} definition of manhood.
\end{itemize}
\end{exemplebox}

\courssub{Themes Synthesis: A Conceptual Map}

The following diagram maps the major characters along two axes that structure the play's ideological debates:

\begin{popfigure}
\begin{tikzpicture}[scale=0.9, every node/.style={font=\small}]
% Axes
\draw[->, thick, popdark] (-3.5,0) -- (3.5,0) node[right] {Assimilation $\longleftrightarrow$ African Identity};
\draw[->, thick, popdark] (0,-3.5) -- (0,3.5) node[above] {Materialism $\longleftrightarrow$ Dignity/Family Values};

% Quadrant labels
\node[popmuted, anchor=north east, text width=2.5cm, align=center] at (-3.2,3.2) {Assimilationist\\Materialist};
\node[popmuted, anchor=north west, text width=2.5cm, align=center] at (3.2,3.2) {Africanist\\Dignity/Family};
\node[popmuted, anchor=south east, text width=2.5cm, align=center] at (-3.2,-3.2) {Assimilationist\\Dignity/Family};
\node[popmuted, anchor=south west, text width=2.5cm, align=center] at (3.2,-3.2) {Africanist\\Materialist};

% Characters as nodes
% Walter: Materialism, middle on identity (shifts)
\node[popblue, circle, draw, thick, minimum size=8pt, fill=popblue!20] (W) at (-1.5,-2) {};
\node[popblue, anchor=north] at (-1.5,-2.4) {Walter (Act 1)};

% Walter Act 2 (shift)
\node[popblue, circle, draw, thick, minimum size=8pt, fill=popblue!20] (W2) at (0.5,1.5) {};
\node[popblue, anchor=south] at (0.5,1.9) {Walter (Act 3)};

% Mama: High dignity, middle identity (rooted in Black Southern tradition)
\node[popgreen, circle, draw, thick, minimum size=8pt, fill=popgreen!20] (M) at (0.5,2.5) {};
\node[popgreen, anchor=south west] at (0.8,2.5) {Mama};

% Beneatha: Shifts from assimilation-curious to Africanist, high dignity
\node[poporange, circle, draw, thick, minimum size=8pt, fill=poporange!20] (B1) at (-1,1) {};
\node[poporange, anchor=east] at (-1.3,1) {Beneatha (early)};

\node[poporange, circle, draw, thick, minimum size=8pt, fill=poporange!20] (B2) at (2,2) {};
\node[poporange, anchor=west] at (2.3,2) {Beneatha (Act 2)};

% George Murchison: Assimilationist, Materialist (class privilege)
\node[poppink, circle, draw, thick, minimum size=8pt, fill=poppink!20] (G) at (-2.5,-1.5) {};
\node[poppink, anchor=east] at (-2.8,-1.5) {George M.};

% Asagai: Africanist, Dignity/Intellectual
\node[poppurple, circle, draw, thick, minimum size=8pt, fill=poppurple!20] (A) at (2.5,1.5) {};
\node[poppurple, anchor=west] at (2.8,1.5) {Asagai};

% Ruth: Practical dignity, survival-focused
\node[popgold, circle, draw, thick, minimum size=8pt, fill=popgold!20] (R) at (-0.5,-0.5) {};
\node[popgold, anchor=north east] at (-0.5,-0.5) {Ruth};

% Lindner: Assimilationist (white norm), Materialist (buying them out)
\node[popmuted, circle, draw, thick, minimum size=8pt, fill=popmuted!20] (L) at (-2.5,-2.5) {};
\node[popmuted, anchor=north east] at (-2.5,-2.5) {Lindner};

% Arrows showing movement
\draw[popblue, ->, thick, dashed] (W) to[bend left=30] node[midway, above, sloped, font=\tiny] {growth} (W2);
\draw[poporange, ->, thick, dashed] (B1) to[bend right=30] node[midway, above, sloped, font=\tiny] {awakening} (B2);

% Legend
\begin{scope}[shift={(4.5,2)}]
\node[popblue, circle, draw, thick, minimum size=6pt, fill=popblue!20] (lw) at (0,0) {};
\node[anchor=west, font=\tiny] at (0.3,0) {Walter};
\node[popgreen, circle, draw, thick, minimum size=6pt, fill=popgreen!20] (lm) at (0,-0.5) {};
\node[anchor=west, font=\tiny] at (0.3,-0.5) {Mama};
\node[poporange, circle, draw, thick, minimum size=6pt, fill=poporange!20] (lb) at (0,-1) {};
\node[anchor=west, font=\tiny] at (0.3,-1) {Beneatha};
\node[poppink, circle, draw, thick, minimum size=6pt, fill=poppink!20] (lg) at (0,-1.5) {};
\node[anchor=west, font=\tiny] at (0.3,-1.5) {George};
\node[poppurple, circle, draw, thick, minimum size=6pt, fill=poppurple!20] (la) at (0,-2) {};
\node[anchor=west, font=\tiny] at (0.3,-2) {Asagai};
\node[popgold, circle, draw, thick, minimum size=6pt, fill=popgold!20] (lr) at (0,-2.5) {};
\node[anchor=west, font=\tiny] at (0.3,-2.5) {Ruth};
\node[popmuted, circle, draw, thick, minimum size=6pt, fill=popmuted!20] (ll) at (0,-3) {};
\node[anchor=west, font=\tiny] at (0.3,-3) {Lindner};
\end{scope}
\end{tikzpicture}
\legende{Character positions on the Assimilation–African Identity and Materialism–Dignity axes. Dashed arrows show developmental arcs.}
\end{popfigure}

\begin{aretenir}[Five Core Themes in Acts 1.2 and 2.1]
\begin{enumerate}
\item \textbf{Money vs. Dignity:} The cheque forces every character to declare what they value. Walter: money = manhood. Mama: dignity = legacy. Beneatha: purpose = identity.
\item \textbf{Assimilation vs. African Roots:} George (assimilation) vs. Asagai (recovery). Beneatha's hair and dress are the battlefield.
\item \textbf{Generational Conflict:} Mama (Southern, religious, collective) vs. Walter (Northern, secular, individualistic) vs. Beneatha (intellectual, experimental). But also \emph{solidarity} across generations against racism.
\item \textbf{Strength of Women:} Mama (matriarch), Ruth (endurance), Beneatha (ambition). They hold the family together while the men fracture.
\item \textbf{The American Dream — Contested:} Is it a house in a white suburb? A liquor store? A medical degree? A return to Africa? The play refuses a single answer.
\end{enumerate}
\end{aretenir}

\courssub{Worked Exercise: Tracing a Theme Across Scenes}

\begin{methode}[How to Trace a Theme from Act 1 to Act 2 (Examinable Skill)]
\begin{enumerate}
\item \textbf{Identify the theme's first articulation.} (e.g., Identity: Beneatha's hobbies in Act 1.1; Walter's ``money is life'' in Act 1.1.)
\item \textbf{Find the \emph{development} in Act 1.2.} (Beneatha meets Asagai; Walter gets the money but not the control.)
\item \textbf{Locate the \emph{dramatisation} in Act 2.1.} (Beneatha's natural hair; Walter's drunken African dance; Mama's house purchase.)
\item \textbf{Note the \emph{conflict} or \emph{synthesis}.} (George vs. Asagai; Walter vs. Mama; the family vs. Lindner — implied.)
\item \textbf{Write a paragraph using the \emph{PEEL} structure:} Point, Evidence (quote/stage direction), Explanation, Link to wider play.
\end{enumerate}
\end{methode}

\begin{exoresolu}[Practice: Trace the theme of ``African Identity'' from Act 1.1 to Act 2.1]
\textbf{Point:} Beneatha's exploration of African identity moves from intellectual curiosity to embodied political statement.
\textbf{Evidence 1 (Act 1.1):} She tells Mama she wants to ``express herself'' through guitar/acting; she asks Asagai ``teach me'' about Africa.
\textbf{Evidence 2 (Act 1.2):} Asagai criticises her straightened hair: ``mutilated''; gives her the nickname \emph{Alaiyo}.
\textbf{Evidence 3 (Act 2.1):} She appears in Nigerian robe, natural hair, dancing to African drums. She tells George: ``I look like \emph{me}!''
\textbf{Explanation:} The progression is \emph{external} (hobbies) $\to$ \emph{intellectual} (Asagai's teaching) $\to$ \emph{corporeal/political} (hair, dress). Hansberry shows identity as \emph{process}, not essence.
\textbf{Link:} This mirrors the play's argument that the Black American future requires \emph{reclaiming} a past severed by slavery — a theme Asagai extends to anti-colonial Africa.
\end{exoresolu}

\courssub{Key Quotations for Examination}

\begin{center}
\begin{tabularx}{\linewidth}{|l|X|l|}
\hline
\rowcolor{popblueL}
\textbf{Speaker} & \textbf{Quotation} & \textbf{Theme/Act} \\
\hline
Mama & ``We ain't never been that dead inside.'' & Dignity > Money (1.2) \\
\hline
Walter & ``Money is life.'' & Materialism (1.2) \\
\hline
Mama & ``Once upon a time freedom used to be life — now it's money.'' & Generational clash (1.2) \\
\hline
Asagai & ``Alaiyo ... One for Whom Bread — Food — Is Not Enough.'' & Identity/Purpose (2.1) \\
\hline
Beneatha & ``I look like \emph{me}!'' & Self-definition (2.1) \\
\hline
George & ``Your heritage is nothing but a bunch of raggedy-assed spirituals and some grass huts!'' & Assimilation critique (2.1) \\
\hline
Mama & ``I come from five generations of people who was slaves and sharecroppers — but ain't nobody in my family never let nobody pay 'em no money that was a way of telling us we wasn't fit to walk the earth.'' & Dignity/History (1.2) \\
\hline
Walter (to Mama) & ``You trust me like that, Mama?'' & Test of manhood (2.1) \\
\hline
\end{tabularx}
\end{center}

\courssub{Summary: What You Must Master for This Section}

\begin{aretenir}[Section 5 Checklist]
\begin{itemize}
\item \textbf{Plot:} Narrate Act 1.2 (cheque, pregnancy, Mama's values) and Act 2.1 (hair, Asagai, Clybourne Park, money handover) accurately.
\item \textbf{Character Arcs:} Walter (desperation $\to$ trust $\to$ betrayal $\to$ redemption), Beneatha (flitting $\to$ Alaiyo), Mama (stewardship $\to$ risk), Ruth (despair $\to$ hope), George vs. Asagai (foils).
\item \textbf{Themes:} Money/Dignity, Assimilation/Roots, Gender/Matriarchy, Generational Conflict, American Dream(s).
\item \textbf{Techniques:} Stage directions (mirror, hair, cheque), dramatic irony (money handover), foil pairs, symbolism (house, hair, plant, cheque).
\item \textbf{Context:} 1959 Chicago — segregation, housing discrimination (Hansberry v. Lee, 1940), Black women's double burden, Pan-Africanism, Cold War conformity.
\item \textbf{Essay Skills:} PEEL paragraphs, tracing themes across scenes, balancing sympathy and critique, using quotations precisely.
\end{itemize}
\end{aretenir}

\begin{exercice}[Consolidation Exercises]
\begin{enumerate}
\item \textbf{Close Analysis:} Annotate the stage direction for Beneatha's entrance in Act 2.1 (natural hair, Nigerian robe, dancing). How does Hansberry use \emph{visual theatre} to convey internal change?
\item \textbf{Comparative:} Contrast George Murchison's and Joseph Asagai's attitudes to African heritage. What does each man represent for Beneatha's future?
\item \textbf{Thematic Essay:} ``The play suggests that the American Dream is inaccessible to Black Americans unless they redefine it on their own terms.'' Discuss with reference to Acts 1.2 and 2.1.
\item \textbf{Character Study:} Is Walter Lee a villain, a victim, or both? Use evidence from his treatment of Ruth, his reaction to the cheque, and his acceptance of the \$6,500.
\item \textbf{Context Link:} Research the 1948 Supreme Court case \emph{Shelley v. Kraemer} (racially restrictive covenants) and the 1940 \emph{Hansberry v. Lee} case. How does the Clybourne Park purchase dramatise this legal history?
\end{enumerate}
\end{exercice}

\begin{corrige}[Exercise 1 — Close Analysis]
\textbf{Key points to include:}
\begin{itemize}
\item \textbf{Visual symbolism:} Natural hair = rejection of white beauty standards (straightening = assimilation). Nigerian robe = affirmative connection to pre-colonial Africa.
\item \textbf{Performance:} Dancing to drums = embodying culture, not just studying it. The phonograph (technology) brings Africa into the Chicago apartment.
\item \textbf{Mirror moment:} ``She frowns. She pats her hair. She smiles.'' The \emph{process} of self-acceptance is shown in real time — not instant, but deliberate.
\item \textbf{Contrast:} Walter's drunken joining-in (playful, temporary) vs. Beneatha's serious commitment. George's horror (``eccentric'') confirms the political weight.
\item \textbf{Stagecraft:} The scene opens with this image — it \emph{sets the ideological tone} for the act before dialogue begins.
\end{itemize}
\end{corrige}

\begin{corrige}[Exercise 3 — Thematic Essay Plan]
\textbf{Thesis:} The Youngers do not merely \emph{access} the American Dream (house, money); they \emph{redefine} it as collective dignity, cultural memory, and resistance.
\textbf{Structure:}
\begin{enumerate}
\item Introduction: Define the ``standard'' American Dream (homeownership, upward mobility). State the play's critique: for Black Americans, the standard dream is blocked by racism (Lindner) and internalised oppression (Walter's liquor store).
\item Paragraph 1: Mama's redefinition — house as \emph{stake in citizenship}, not asset. Clybourne Park as \emph{battlefield}.
\item Paragraph 2: Beneatha/Asagai — dream expands beyond America to Pan-African liberation. ``Alaiyo'' = purpose beyond survival.
\item Paragraph 3: Walter's arc — from ``money is life'' to ``we don't want your money'' (Act 3). Redefinition of manhood as \emph{moral choice}.
\item Paragraph 4: The women (Mama, Ruth) as dream-keepers — their labour makes the dream \emph{liveable}.
\item Conclusion: The play ends \emph{en route} to Clybourne Park — the redefined dream is a \emph{practice}, not a possession.
\end{enumerate}
\end{corrige}

\coursec{Section 6: Nineteen Eighty-Four — OrContext,  and Review of Part One}

\courssub{Orwell's Life and the Context}
In Section 5, we left the question of who controls the means of production, and the question of who controls the means of expression. In one of the most important means of production, the question of who controls the means of production, the question of who controls the means of production, the question of who controls the means of production, the question of who controls of production, the question of who controls the means of production, the question of who controls the means of production, the question of production, the question of who controls the means of production, the question of production, of the question of production, of the question of production, of of the of the of the of the of the of the of the of the of the of the of the of the of the of the of the of the of the of the of the of the of the of of of the of the of the of the of of the of the of the of the of the of of the of of the of of of the of of the of of of of of the of of of of of the of of of the of of of of of of of of of of of of of of of of of of of of of of of of of of of of of of of of

\coursec{Section 7: Nineteen Eighty-Four --- Part Two, Chapters 1--6: Love as Rebellion}

In Section 6 we reviewed Part One of \emph{Nineteen Eighty-Four}: we met Winston Smith, a lonely, sickly Outer Party member in Airstrip One, secretly writing a diary, haunted by memories of his mother, dreaming of a ``Golden Country'', and committing the small, silent crime of \cle{thoughtcrime} every day. Part One ends with Winston's despairing conviction that he is already dead --- that the Party will inevitably catch him. Part Two opens at exactly that point of despair and answers it with the most unexpected event in the novel: a girl falls at his feet, and into his hand she slips a scrap of paper bearing three words --- \emph{I love you}. From this moment, the novel changes key. What follows in Chapters 1 to 6 is simultaneously a love story and a political thriller, and our task in this section is to see how Orwell fuses the two: how \cle{love becomes rebellion}, and how every step towards intimacy is also a step towards arrest.

\courssub{Chapter 1: The Fall in the Corridor --- ``I Love You''}

Winston has noticed the dark-haired girl from the Fiction Department before, and --- in a masterpiece of Orwellian irony --- he has assumed she is a spy of the Thought Police watching him. He even fantasises about smashing her skull with a cobblestone. Then, one morning, she falls in the corridor (her arm is in a sling), he helps her up, and she presses a note into his hand. Trembling, he opens it at his desk, expecting a political message or a threat. Instead he reads: \emph{I love you}.

The effect on Winston is electric. Orwell shows us that in Oceania, desire itself has been so thoroughly policed that a declaration of love is indistinguishable from a declaration of war on the Party. The girl is \cle{Julia}, and Winston spends days in an agony of suspense trying to arrange a word with her in the canteen, under the eyes of the telescreens.

\begin{aretenir}[Love as a political act]
In a totalitarian state that demands the total, exclusive loyalty of every citizen, \textbf{any private attachment is treason}. The Party wants no love except love of Big Brother. Therefore, when Julia writes ``I love you'', she is not merely expressing a feeling: she is committing a crime and inviting Winston to share it. Orwell's point is precise --- in Oceania, \emph{the personal is political by force}.
\end{aretenir}

\courssub{Chapter 2: The Golden Country and the Thrush's Song}

Their first real meeting takes place in the countryside, far from the telescreens --- in a clearing of bluebells that Winston recognises with a shock: it is the \cle{Golden Country} of his recurring dream from Part One. Orwell's symbolism here is deliberate. The dream landscape of freedom, associated with his mother, with the past, and with the dark-haired girl, becomes real. Nature --- grass, birdsong, sunlight --- stands outside the Party's reach, a world the Party cannot pave over.

The most famous detail of the chapter is the \cle{thrush} that sings in a tree above them as they lie together. Winston marvels that the bird sings for no audience, no purpose, no profit --- simply because singing is its nature. The thrush embodies exactly what the Party denies human beings: \emph{spontaneous, purposeless, joyful life}. Julia and Winston make love, and Winston reflects that their embrace was ``a battle, the climax a victory'': it was ``a blow struck against the Party. It was a political act.''

\begin{exemplebox}[The Party's anti-sex puritanism]
Why is sex a crime? Orwell explains through Julia's membership of the \textbf{Junior Anti-Sex League} (the scarlet sash she wears). The Party channels the energy of desire into \emph{loyalty, hysteria and hatred} --- the Two Minutes Hate, the war fever. A satisfied, privately happy citizen is a citizen whose energy escapes the Party. Julia understands this instinctively: she tells Winston the Party wants to stop people enjoying sex because unexpressed desire can be converted into fanaticism. Their affair is therefore not escapism; it is sabotage.
\end{exemplebox}

\courssub{Chapter 3: Two Kinds of Rebel}

Back in London, meeting secretly whenever they can, Winston and Julia talk --- and Orwell uses their conversations to draw one of the finest character contrasts in the novel. The examiner's favourite distinction is here.

\textbf{Winston's rebellion is intellectual and historical.} He broods over the past, over whether life was better before the Revolution, over the mutability of records, over the proles. He wants to \emph{understand} the Party and to defeat it in principle, even though he believes he will fail.

\textbf{Julia's rebellion is practical and personal.} She is twenty-six, has had many affairs, and is brilliantly skilled at the small arts of survival --- obtaining real sugar, real coffee, real bread and jam through the black market. She is bored by Winston's talk of the Brotherhood and of abstract principles. She does not care whether the Party has rewritten history; she simply refuses to let it stop her from living.

\begin{attention}[Do not merge Winston and Julia]
A frequent examination error is to present Winston and Julia as identical partners in resistance. They are not. Winston himself delivers the verdict: \emph{``You're only a rebel from the waist downwards.''} Julia rebels with her body and her appetites; Winston rebels with his mind and his memory. Examiners consistently reward candidates who make this distinction --- and who note that Orwell neither mocks Julia's pragmatism nor fully endorses Winston's theorising. Each is incomplete without the other.
\end{attention}

\courssub{Chapter 4: The Room Above Mr Charrington's Shop}

Winston takes the decisive step of renting the little upstairs room above \cle{Mr Charrington}'s junk shop --- the room he had noticed in Part One, which appears to have no telescreen. Here the lovers create a fragile private world: a real bed with a mahogany bedstead, an old-fashioned twelve-hour clock, coffee, sugar, and Julia's scandalous luxuries --- real chocolate, perfume, make-up, a dress. She transforms herself from a uniformed Party member into ``a woman'', an act she performs with delight precisely because the Party forbids it.

The room's central object is the \cle{glass paperweight} Winston bought from Charrington: a hemisphere of glass with a piece of pink coral embedded at its heart.

\begin{aretenir}[Key symbol: the glass paperweight]
The paperweight symbolises \textbf{the tiny, enclosed, beautiful world} Winston and Julia build together --- and, more deeply, the \textbf{past itself}, preserved intact inside glass, unreachable by the Party's rewriting. The coral is their love; the glass is the room; and the object's fragility is the point. Winston explicitly links the room to the paperweight: the coral is Julia and himself, fixed forever inside the crystal. The attentive reader should already hear the warning: \emph{glass breaks}.
\end{aretenir}

Above the fireplace hangs the print of \textbf{St Clement's Church}, and Mr Charrington has taught Winston the old nursery rhyme: ``Oranges and lemons, say the bells of St Clement's\ldots'' The rhyme --- a fragment of the lost past --- will return with devastating force. Remember it.

\courssub{Chapter 5: Hatred at Home, Hope in the Proles}

Chapter 5 deliberately darkens the idyll. Preparations for \cle{Hate Week} fill the city; a gigantic poster of a Eurasian soldier glares over the square; a public hanging of prisoners is announced. Orwell juxtaposes the lovers' private happiness with the public machinery of hatred, so that the reader feels the two worlds cannot coexist for long.

Two details demand close analysis:

\begin{itemize}
\item \textbf{Parsons's pride.} Winston's neighbour, the sweating, stupid, devoted Parsons, boasts that his seven-year-old daughter denounced a man to the Thought Police as a spy because his shoes looked foreign. Parsons calls it a ``pretty'' achievement. Orwell shows us the Party's ultimate victory: \emph{the family itself has been turned into an instrument of surveillance}, and a father is proud of being betrayed by his own child. This is the world in which Winston and Julia dare to love.
\item \textbf{Julia and the war.} Winston reads and talks politics; Julia listens, then falls asleep, or dismisses the whole question. She even suggests that the war may be a fiction and that the Party bombs its own citizens to keep them frightened --- a shrewd guess --- but she does not \emph{care}. Abstract politics bores her; she rebels because she wants her own life, not because she wants to change history.
\end{itemize}

It is in this atmosphere that Winston repeats the conviction that has become his political creed:

\begin{aretenir}[Key quotation: ``If there is hope, it lies in the proles'']
Winston's central political belief is that the \cle{proles} --- the great majority of the population, left outside the Party and largely unpoliced --- are the only force capable of overthrowing it. They have kept their humanity: they love, grieve, sing and remember. Yet Winston also sees the tragedy: the proles are too absorbed in private life, too lacking in consciousness, to rise. ``Until they become conscious they will never rebel, and until after they have rebelled they cannot become conscious.'' The quotation expresses both Winston's hope and his despair --- and it frames his affair with Julia as a small, prole-like act of ordinary human happiness stolen from the Party.
\end{aretenir}

\courssub{Chapter 6: O'Brien's Approach}

The chapter --- and the first movement of Part Two --- closes with the moment Winston has longed for and dreaded. \cle{O'Brien}, the imposing Inner Party member whom Winston has always half-believed to be a secret dissident, speaks to him in the Ministry. Under cover of discussing the new edition of the Newspeak Dictionary, O'Brien gives Winston his address and invites him to his flat.

Winston is overwhelmed. He believes the signal has come at last: that the underground resistance, the \cle{Brotherhood}, exists, and that O'Brien is part of it. Orwell describes Winston's feeling that he is stepping ``into the dampness of a grave'' --- he knows this path leads to the cellars of the Ministry of Love, and he walks it willingly. The chapter ends Part Two's first phase on a note of exalted, fatalistic hope.

\begin{attention}[Dramatic irony]
The reader who watches carefully notices what Winston refuses to see: O'Brien knows far too much, and his approach is suspiciously smooth. Orwell builds \emph{dramatic irony} --- we suspect the trap while Winston embraces it. In your essays, note how this irony creates suspense and pity: we watch a man walk towards his destruction believing it is his salvation.
\end{attention}

\courssub{Themes of Part Two (Chapters 1--6)}

\begin{center}
\begin{tabularx}{\linewidth}{|l|X|}
\hline
\rowcolor{popblueL} \textbf{Theme} & \textbf{How Orwell develops it in Chapters 1--6} \\
\hline
Love and sexuality as resistance & The affair is ``a political act''; desire redirected away from Big Brother is sabotage. The thrush and the Golden Country figure natural, free life. \\
\hline
The impossibility of privacy & Every meeting requires lies, disguises and luck; the room feels safe but is watched; the telescreen is hidden behind the picture. \\
\hline
Trust and betrayal & Winston misreads Julia (as a spy) and O'Brien (as a friend); Parsons's daughter betrays a stranger; betrayal is the air Oceania breathes. \\
\hline
The proles as hope & ``If there is hope, it lies in the proles'' --- yet they lack consciousness; Winston and Julia live a stolen version of prole life. \\
\hline
The past and memory & The paperweight, the rhyme, the print of St Clement's, the old clock --- fragments of a world the Party has abolished. \\
\hline
\end{tabularx}
\end{center}

\begin{definition}[Foreshadowing]
\cle{Foreshadowing} is a narrative technique by which the writer plants hints, images or details early in the story that prepare the reader for later events. In Part Two, Orwell foreshadows the lovers' capture relentlessly: the \textbf{picture of St Clement's} (behind which the telescreen is hidden), the \textbf{rats} Winston fears in the room (which will return in the Ministry of Love), the \textbf{rhyme} that ends ``here comes a chopper to chop off your head'', and the \textbf{glass paperweight}, beautiful and breakable. Nothing in these chapters is innocent; every happiness carries its destruction inside it.
\end{definition}

\begin{savaistu}[Did you know?]
The rhyme of ``Oranges and Lemons'' is a genuine traditional English nursery rhyme about the bells of London churches, which Orwell weaves through the novel as a thread connecting Winston to a past the Party has erased. Its final line --- ``Here comes a chopper to chop off your head'' --- is completed for Winston at the very moment of his arrest, turning a children's game into a death sentence. Orwell's point is chilling: in Oceania, even the recovery of memory is a path to destruction.
\end{savaistu}

\courssub{Illustration: Intimacy Rising, Danger Rising}

The structure of Part Two, Chapters 1--6, can be read as two converging curves: as Winston and Julia's intimacy deepens, the risk of discovery rises in parallel, until the meeting with O'Brien brings the two lines together at the edge of catastrophe.

\begin{popfigure}
\begin{center}
\begin{tikzpicture}[scale=1.0]
\draw[->, thick, popdark] (0,0) -- (11,0) node[below left, align=left, text width=3cm] {\footnotesize Part Two, Chapters 1--6};
\draw[->, thick, popdark] (0,0) -- (0,4.6) node[above left, align=left, text width=2.2cm] {\footnotesize Intensity};
\foreach \x/\c in {0.6/{Ch.\,1}, 2.4/{Ch.\,2}, 4.2/{Ch.\,3}, 6.0/{Ch.\,4}, 7.8/{Ch.\,5}, 9.6/{Ch.\,6}} {
  \draw[popmuted] (\x,0) -- (\x,4.3);
  \node[below, popdark] at (\x,0) {\footnotesize \c};
}
\draw[very thick, popblue, smooth] plot coordinates {(0.6,0.5) (2.4,1.6) (4.2,2.2) (6.0,3.4) (7.8,3.7) (9.6,3.9)};
\node[popblue, align=center, text width=3.4cm] at (3.2,3.9) {\footnotesize Intimacy: note, Golden Country, room, shared life};
\draw[very thick, poporange, smooth] plot coordinates {(0.6,0.3) (2.4,0.9) (4.2,1.4) (6.0,2.3) (7.8,3.1) (9.6,4.2)};
\node[poporange, align=center, text width=3.4cm] at (7.6,1.1) {\footnotesize Danger: suspicion, Hate Week, hidden telescreen, O'Brien};
\filldraw[poppink] (9.6,4.2) circle (2pt);
\node[poppink, above, align=center, text width=3cm] at (9.6,4.2) {\footnotesize Convergence: the trap closes};
\end{tikzpicture}
\end{center}
\legende{The double movement of Part Two, Chapters 1--6: as love deepens (blue), danger rises (orange). Orwell engineers the two curves to meet at the moment of greatest hope --- the approach to O'Brien.}
\end{popfigure}

\courssub{Method and Worked Example}

\begin{methode}[Analysing a love story in a political novel]
When the GCE examiner asks you to discuss the love affair in \emph{Nineteen Eighty-Four}, do not retell the romance. Proceed in four steps:
\begin{enumerate}
\item \textbf{State the political context.} Remind the reader that the Party demands total loyalty and channels desire into hatred; therefore private love is already rebellion.
\item \textbf{Analyse the affair as an act.} Quote or paraphrase the key idea: their embrace is ``a blow struck against the Party''. Use the thrush, the Golden Country and the room as evidence.
\item \textbf{Contrast the two rebels.} Distinguish Julia's practical, bodily rebellion from Winston's intellectual, historical one (``a rebel from the waist downwards'').
\item \textbf{Show the fragility.} Use foreshadowing (paperweight, rats, rhyme, picture) to argue that Orwell makes the private world beautiful \emph{precisely so that its destruction will demonstrate the Party's power}.
\end{enumerate}
\end{methode}

\begin{exoresolu}[Model paragraph: ``Their embrace was a political act'']
\textbf{Question:} Show how Orwell presents the love affair of Winston and Julia as a form of rebellion in Part Two, Chapters 1--4.
\tcblower
\textbf{Model paragraph:} Orwell presents the affair not as an escape from politics but as politics itself. In Oceania the Party claims the whole of the citizen's emotional life, diverting sexual energy into the hysteria of the Two Minutes Hate and the war fever; Julia, a member of the Junior Anti-Sex League, understands that the Party's puritanism exists precisely to manufacture this fanaticism. When the lovers meet in the Golden Country --- the landscape of Winston's dreams of freedom --- and the thrush sings above them for no audience and no purpose, Orwell sets natural, spontaneous joy against the Party's total control. Winston reflects that their embrace was ``a battle, the climax a victory'': desire redirected away from Big Brother is sabotage. The rented room above Mr Charrington's shop then gives this rebellion a home, symbolised by the glass paperweight whose coral heart encloses a tiny world the Party cannot rewrite. Yet Orwell makes the idyll fragile on purpose: the hidden telescreen, the rats and the rhyme of St Clement's foreshadow its destruction, so that the love affair measures both the possibility and the price of resistance.
\end{exoresolu}

\begin{exercice}[Practice questions]
\begin{enumerate}
\item Trace the stages of Winston and Julia's relationship from the note in Chapter 1 to the rented room in Chapter 4, showing how each stage increases both their happiness and their danger.
\item ``Julia and Winston are rebels of two different kinds.'' Discuss with close reference to Chapters 2, 3 and 5.
\item Analyse the symbolic importance of the glass paperweight and the rhyme of St Clement's as instruments of foreshadowing in Part Two.
\item ``If there is hope, it lies in the proles.'' What does Winston mean, and how does his affair with Julia relate to this belief?
\end{enumerate}
\end{exercice}

\begin{corrige}[Exercise 1 --- guidance]
Strong answers will identify four stages: (1) the note ``I love you'' --- love announced as crime, with Winston's earlier murderous fantasy as ironic contrast; (2) the Golden Country --- the affair consummated as ``a political act'', framed by the thrush's purposeless song; (3) the snatched meetings and conversations --- the discovery that the two rebels differ in kind, Julia practical, Winston intellectual; (4) the rented room --- the building of a private world symbolised by the paperweight, shadowed by the picture, the rats and the rhyme. At each stage the candidate should note Orwell's structural irony: rising intimacy is matched by rising risk, so that the reader's pleasure is always poisoned by foreknowledge of capture. The best scripts will conclude that Orwell uses the love story to define what totalitarianism destroys --- ordinary, private, purposeless human happiness.
\end{corrige}

\newpage

\coursec{Integration activity}

\coursec{Integration Activity: The National Literary Tribunal}

\courssub{Aims of the Activity}
This integration activity consolidates the three set texts studied in the first term — Chaucer’s \emph{The Canterbury Tales} (General Prologue), Hansberry’s \emph{A Raisin in the Sun} (Acts 1–2), and Orwell’s \emph{Nineteen Eighty-Four} (Part One and Part Two, Chapters 1–6) — by requiring you to:
\begin{itemize}
    \item Demonstrate secure knowledge of each author’s historical and literary context.
    \item Apply close‑reading skills to select precise textual evidence (short quotations in the original language or standard translation, with line/scene references) that supports a clear thesis.
    \item Construct a coherent, comparative argument in the style of a GCE Advanced Level Paper 2 appreciation essay.
    \item Practise oral defence and critical questioning, mirroring the “tribunal” format used in university seminars and professional literary debates.
    \item Reflect on the relevance of these three fictional societies to contemporary Cameroon, thereby linking literary study to civic awareness.
\end{itemize}

\begin{aretenir}[Examiner’s Expectation]
GCE Paper 2 rewards: (1) relevance to the exact wording of the question; (2) well‑chosen, accurately cited textual evidence; (3) logical organisation with clear comparative connectives; (4) precise literary terminology (estates satire, dramatic irony, Newspeak, doublethink, etc.); (5) fluent, varied expression. This activity rehearses all five criteria.
\end{aretenir}

\courssub{Situation}
The Ministry of Arts and Culture in Yaoundé has organised a \textbf{National Literary Tribunal} as part of “Literature and Society” week. Three panels of Upper Sixth students — each representing one set text — will present their case before a jury of teachers, university lecturers, and invited writers. The tribunal’s central motion is:

\begin{quote}
\textbf{“Which of the three societies portrayed in our set texts — Chaucer’s late‑medieval England, Hansberry’s 1950s Chicago South Side, or Orwell’s Oceania — offers the sharpest lens for understanding the challenges facing Cameroon today?”}
\end{quote}

Each panel must:
\begin{enumerate}
    \item Deliver a two‑minute contextual introduction situating the author and the work’s original reception.
    \item Defend the proposition that their text provides the \emph{most incisive} portrait of a society under stress, using \textbf{at least three precise textual references} (short quotations with line/scene numbers or unmistakable scene references).
    \item Respond to questions from the jury and the audience (the other two panels) for five minutes.
\end{enumerate}

After the oral phase, every student — regardless of panel — will write a \textbf{one‑page comparative paragraph} (approximately 300–350 words) answering the tribunal’s motion. This written piece will be assessed using the GCE Paper 2 mark scheme: Content (10), Organisation (5), Expression (5).

\begin{attention}[Common Pitfall]
Do not merely list parallels. The task demands an \emph{evaluative judgement}: you must argue why \textbf{one} society is \emph{closer} to Cameroon’s present reality than the other two, using comparative language (\emph{more so than}, \emph{whereas}, \emph{in contrast}, \emph{unlike}).
\end{attention}

\courssub{Tasks}

\begin{methode}[Step‑by‑Step Procedure]
\begin{enumerate}
\item \textbf{Contextual Research (Individual)} — Using your notes and approved reference works, prepare a concise two‑minute script covering:
      \begin{itemize}
          \item Author’s biography (key dates, major influences).
          \item Historical/literary period (late‑medieval estates satire; American Civil Rights era / domestic realism; post‑war totalitarianism / dystopian fiction).
          \item The work’s initial reception and lasting critical reputation.
      \end{itemize}
\item \textbf{Thesis Formulation (Group)} — As a panel, agree on a single, arguable thesis statement answering the motion for your text. 
      \begin{itemize}
          \item \emph{Example (Chaucer):} “Chaucer’s General Prologue exposes systemic corruption through estates satire, making it the clearest mirror for Cameroon’s contemporary patronage networks.”
          \item \emph{Example (Hansberry):} “A Raisin in the Sun best shows how sudden capital tests human dignity, a dynamic echoing Cameroonian families confronting remittances or land compensation.”
          \item \emph{Example (Orwell):} “Nineteen Eighty‑Four dissects the mechanisms — surveillance, language engineering, historical revisionism — that a state uses to colonise the mind, mechanisms observable in Cameroon’s digital crackdowns and curriculum reforms.”
      \end{itemize}
\item \textbf{Evidence Selection (Group)} — Choose \textbf{at least three} textual references directly supporting your thesis. For each, record:
      \begin{itemize}
          \item Exact line numbers (Chaucer) or act/scene/page (Hansberry, Orwell).
          \item Literary device employed (irony, symbolism, dramatic irony, stage direction, narrative voice, etc.).
          \item Specific societal ill illuminated (ecclesiastical greed, deferred dreams, state surveillance, etc.).
      \end{itemize}
\item \textbf{Oral Presentation (Group)} — Deliver in this structure (total $\approx$ 6 minutes):
      \begin{enumerate}
          \item Context (2 min)
          \item Thesis statement (15 sec)
          \item Evidence 1 $\rightarrow$ Analysis (1 min)
          \item Evidence 2 $\rightarrow$ Analysis (1 min)
          \item Evidence 3 $\rightarrow$ Analysis (1 min)
          \item Concluding sentence linking to Cameroon (15 sec)
      \end{enumerate}
\item \textbf{Cross‑Examination (Whole Class)} — Jury and opposing panels question you. You must:
      \begin{itemize}
          \item Defend your readings against alternative interpretations.
          \item Acknowledge limitations of your text’s relevance.
          \item Draw connections to the other two texts where appropriate.
      \end{itemize}
\item \textbf{Comparative Paragraph (Individual)} — Write a single, well‑structured paragraph (300–350 words) answering the motion. It must:
      \begin{itemize}
          \item Open with a clear topic sentence naming the chosen society.
          \item Integrate \textbf{one concise quotation from each of the three texts} (you quote all three; the class collectively covers nine quotations).
          \item Explain how each quotation illustrates a parallel in present‑day Cameroon.
          \item Conclude with an evaluative judgement using comparative language.
      \end{itemize}
\item \textbf{Self‑Assessment (Individual)} — Using the GCE Paper 2 grid (Content 10, Organisation 5, Expression 5), award yourself a provisional mark and write two specific targets for improvement.
\end{enumerate}
\end{methode}

\courssub{Model Answers — Panel Presentations}

\begin{exemplebox}[Panel 1: Chaucer — The Canterbury Tales (General Prologue)]
\textbf{Thesis:} \emph{“Chaucer’s General Prologue remains the sharpest portrait of a corrupt society because its estates satire reveals how institutional roles are hijacked for personal gain, a pattern mirroring Cameroon’s contemporary bureaucratic and religious landscapes.”}

\textbf{Context (two minutes):} Geoffrey Chaucer (c.1343–1400) served as courtier, diplomat, and customs official under Edward III, Richard II, and Henry IV. His \emph{Canterbury Tales} (begun c.1387) uses a pilgrimage frame to assemble a cross‑section of late‑medieval English society. The General Prologue is an \cle{estates satire}: each pilgrim represents a social estate (nobility, clergy, peasantry) and is measured against the ideal of that estate. The work circulated in manuscript before Caxton’s 1476 print; it was celebrated for its vernacular vitality and moral insight.

\textbf{Evidence 1 — The Friar (ll. 221–222, Riverside Chaucer):} 
\begin{quote}
“He was an easy man to yeve penaunce / Ther as he wiste to han a good pitaunce.”
\end{quote}
(Modern translation: “He was an easy man in penance‑giving / Where he knew he could get a good pittance.”)
\textbf{Device:} Verbal irony — “easy” suggests pastoral kindness but means commercially lenient. 
\textbf{Analysis:} The verb “wiste” (knew) and noun “pitaunce” reduce sacramental grace to a calculated transaction. This anticipates modern “prosperity gospel” preachers who demand “seed faith” offerings in Yaoundé’s megachurches.

\textbf{Evidence 2 — The Summoner (ll. 647–650):} 
\begin{quote}
“He wolde suffre, for a quart of wyn, / A good felawe to have his concubyn.”
\end{quote}
\textbf{Device:} Situational irony — an officer of the ecclesiastical court facilitating sexual licence. 
\textbf{Analysis:} The Summoner trades legal authority for sexual and financial favours. This mirrors Cameroonian judicial officers who allegedly trade verdicts for bribes, turning law into a commodity.

\textbf{Evidence 3 — The Pardoner (ll. 696–697):} 
\begin{quote}
“For in his male he hadde a pilwe-beer, / Which that he seyde was Oure Lady veyl.”
\end{quote}
\textbf{Device:} Symbolism — the pillow‑case (pilwe-beer) stands for manufactured authenticity. 
\textbf{Analysis:} The Pardoner sells fake relics, openly admitting his fraud in his Prologue. Today, counterfeit medicines and forged academic certificates circulate in Douala and Buea, endangering lives and devaluing merit.

\textbf{Link to Cameroon:} All three figures hold \emph{institutional office} (clergy, court officer, papal agent) yet subvert it for personal enrichment. Reports from Cameroon’s National Anti‑Corruption Commission (CONAC) highlight similar diversion of public funds by officials entrusted with health, education, and justice portfolios. Chaucer’s irony — describing the Friar as “worthy” (l. 243) while detailing his vices — trains readers to spot the gap between title and conduct, a critical skill for any citizen monitoring governance.
\end{exemplebox}

\begin{exemplebox}[Panel 2: Hansberry — A Raisin in the Sun (Acts 1–2)]
\textbf{Thesis:} \emph{“A Raisin in the Sun best shows how money tests human dignity because the Younger family’s \$10,000 insurance cheque becomes a crucible revealing each character’s values, a dynamic echoing Cameroonian families confronting sudden remittances or land compensation.”}

\textbf{Context (two minutes):} Lorraine Hansberry (1930–1965) wrote the play in 1959, drawing on her family’s legal battle against racial restrictive covenants in Chicago (\emph{Hansberry v. Lee}, 1940). The play premiered on Broadway — the first by a Black woman — and won the New York Drama Critics’ Circle Award. It captures the Great Migration’s aftermath: Black families crowded into South Side “kitchenettes”, dreaming of home‑ownership amid systemic racism.

\textbf{Evidence 1 — Walter Lee (Act 1, Scene 1):} 
\begin{quote}
“I want so many things that they are driving me kind of crazy… Mama, I want to buy a liquor store.”
\end{quote}
\textbf{Device:} Dramatic monologue revealing internal conflict; the liquor store symbolises capitalist assimilation. 
\textbf{Analysis:} Walter equates manhood with capital. His desperation reflects Cameroonian youths who, receiving diaspora remittances, rush into high‑risk ventures (motorbike taxis, crypto schemes) to prove provider status.

\textbf{Evidence 2 — Mama (Lena) (Act 1, Scene 2):} 
\begin{quote}
“There is always something left to love. And if you ain’t learned that, you ain’t learned nothing.”
\end{quote}
\textbf{Device:} Sententia (maxim) delivering the play’s moral centre; contrast with Walter’s materialism. 
\textbf{Analysis:} Mama refuses to let money erase compassion. This mirrors the communal ethic in Bamileke or Bassa families where elders insist windfall profits be shared for school fees, funerals, or village water projects.

\textbf{Evidence 3 — Beneatha (Act 2, Scene 1):} 
\begin{quote}
“I am not going to be a doctor to make money… I want to cure.”
\end{quote}
\textbf{Device:} Character foil — Beneatha’s vocation resists instrumentalisation. 
\textbf{Analysis:} In Cameroon, medical students often face pressure to choose lucrative specialisations over rural service; Beneatha’s stance champions vocation over salary.

\textbf{Link to Cameroon:} The \$10,000 cheque functions like a “Prime Minister’s scholarship” or a “land compensation payout” in the Northwest: it can either unite a family around a shared project (a house in Clybourne Park / a family compound in Bafoussam) or fracture it through competing individual dreams. Hansberry’s stage direction — “the light grows brighter” as the family decides to move — offers a visual metaphor for collective agency that resonates with Cameroon’s cooperative movements.
\end{exemplebox}

\begin{exemplebox}[Panel 3: Orwell — Nineteen Eighty-Four (Part One \& Part Two, Chs 1–6)]
\textbf{Thesis:} \emph{“Nineteen Eighty‑Four is the most powerful warning against oppression because it dissects the mechanisms — surveillance, language engineering, historical revisionism — that a state uses to colonise the mind, mechanisms observable in Cameroon’s digital crackdowns and curriculum reforms.”}

\textbf{Context (two minutes):} George Orwell (1903–1950) wrote \emph{Nineteen Eighty‑Four} (1948) while dying of tuberculosis, informed by his Spanish Civil War experience (POUM militia) and observation of Stalinist purges. The novel depicts Oceania, a totalitarian superstate ruled by the Party and Big Brother. It was an immediate bestseller and has supplied political vocabulary worldwide.

\textbf{Evidence 1 — Telescreen (Part 1, Ch. 1):} 
\begin{quote}
“The instrument (the telescreen, it was called) could be dimmed, but there was no way of shutting it off completely.”
\end{quote}
\textbf{Device:} Extended metaphor — the telescreen as ubiquitous surveillance. 
\textbf{Analysis:} Cameroon’s 2020–2022 internet shutdowns in the Anglophone regions and mandatory biometric SIM‑card registration create a de facto telescreen: citizens self‑censor knowing their digital footprint is monitored.

\textbf{Evidence 2 — Newspeak (Part 1, Ch. 5, Syme’s explanation):} 
\begin{quote}
“Don’t you see that the whole aim of Newspeak is to narrow the range of thought?”
\end{quote}
\textbf{Device:} Expository dialogue revealing the Party’s linguistic theory. 
\textbf{Analysis:} The 2017–2018 curriculum reforms that marginalised Anglophone legal and educational terminology (e.g. replacing “Common Law” terms with “Civil Law” equivalents in official examinations) function as Newspeak, shrinking the conceptual space for minority identity.

\textbf{Evidence 3 — Doublethink (Part 1, Ch. 3):} 
\begin{quote}
“To know and not to know, to be conscious of complete truthfulness while telling carefully constructed lies.”
\end{quote}
\textbf{Device:} Paradox defining the cognitive mechanism of totalitarian control. 
\textbf{Analysis:} Official discourse proclaiming “decentralisation” while recentralising appointment powers exemplifies doublethink; citizens learn to parrot slogans while privately acknowledging the gap.

\textbf{Link to Cameroon:} Orwell’s genius is to show that oppression is not only boots on streets but \emph{internalised}. The Party’s slogan “Who controls the past controls the future” echoes contested narratives around the 1961 Foumban Conference and the 1972 referendum — histories rewritten in school textbooks to legitimise unitary statehood. Recognising these techniques equips students to demand transparency in electoral code revisions.
\end{exemplebox}

\courssub{Comparative Paragraph (Model — 340 words)}

\begin{exemplebox}[Model Response]
Of the three societies, Orwell’s Oceania feels closest to Cameroon today because its architecture of mental colonisation — surveillance, linguistic engineering, and historical revisionism — operates with a subtlety that mirrors our own digital and educational landscape. The telescreen that “could be dimmed, but there was no way of shutting it off completely” finds its analogue in the mandatory biometric SIM registration and periodic internet blackouts in the Northwest and Southwest regions, where citizens internalise the certainty of being watched and consequently mute dissent on WhatsApp groups. Similarly, the Party’s Newspeak project — “to narrow the range of thought” — resonates with the gradual replacement of Anglophone legal terminology (e.g. “writ of habeas corpus”) by Francophone equivalents in official examinations, a curricular shift that shrinks the vocabulary available to articulate minority grievances. Chaucer’s corrupt clergy — the Friar who “wiste to han a good pitaunce” from penance — and Hansberry’s Walter Lee, who equates manhood with a liquor‑store licence, certainly reflect Cameroonian realities of ecclesiastical commercialisation and youth entrepreneurship driven by remittance pressure. Yet these are \emph{symptoms} of individual moral failure within functioning institutions, whereas Orwell depicts a \emph{system} designed to make truth itself contingent on power. The Party’s doublethink — “to know and not to know” — captures the cognitive dissonance of a citizenry that publicly celebrates “decentralisation” while observing the recentralisation of administrative appointments. Because Oceania’s oppression targets the \emph{preconditions of resistance} — language, memory, privacy — it offers the most diagnostic lens for a Cameroon where the battle is increasingly fought over textbooks, data sovereignty, and the right to remember. Therefore, while all three texts illuminate facets of our society, \emph{Nineteen Eighty‑Four} provides the most comprehensive framework for analysing the structural threats to Cameroonian intellectual freedom.
\end{exemplebox}

\courssub{Self‑Assessment Targets}

\begin{aretenir}[My Provisional Mark \& Targets]
\textbf{Mark (Content 10 / Organisation 5 / Expression 5):} \_\_\_ / 20

\textbf{Target 1:} Integrate more precise line references (e.g. Chaucer ll. 221–222, 647–650, 696–697) to meet the “detailed textual evidence” criterion fully.

\textbf{Target 2:} Vary comparative connectives beyond “similarly” and “whereas” (use \emph{moreover, conversely, by contrast, in sharp distinction, analogously}) to demonstrate a wider range of cohesive devices.
\end{aretenir}

\begin{savaistu}[Did You Know?]
\begin{itemize}
    \item Chaucer’s “estates satire” follows a tradition going back to the 12th‑century \emph{Estates of the World} poems, but he individualises each pilgrim so vividly that they transcend their estate.
    \item Hansberry’s title comes from Langston Hughes’s poem “Harlem”: “What happens to a dream deferred? / Does it dry up / like a raisin in the sun?”
    \item Orwell originally titled the novel \emph{The Last Man in Europe}; the publisher suggested \emph{Nineteen Eighty-Four} (inverting 1948).
    \item The GCE Paper 2 mark scheme allocates 10 marks for \textbf{Content} (knowledge + relevance), 5 for \textbf{Organisation} (structure + coherence), 5 for \textbf{Expression} (vocabulary + grammar + literary terminology).
\end{itemize}
\end{savaistu}

\begin{attention}[Final Reminders for the Tribunal]
\begin{itemize}
    \item \textbf{Time management:} Rehearse with a stopwatch. Two minutes for context passes quickly.
    \item \textbf{Quotation accuracy:} For Chaucer, quote in Middle English (Riverside edition) and provide a brief gloss. For Hansberry and Orwell, use the standard edition page/act/scene references.
    \item \textbf{Cameroon parallels:} Be specific — name regions, institutions, or policies (e.g. “CONAC reports”, “2020–2022 internet shutdowns”, “2017 curriculum reforms”). Vague generalisations (“corruption exists”) earn few marks.
    \item \textbf{Cross‑examination:} Listen carefully. Concede valid counter‑points (“Chaucer’s clergy are not \emph{all} corrupt — the Parson is an ideal”) — this shows critical maturity.
\end{itemize}
\end{attention}

\newpage

\coursec{Methods and worked examples}

\courssub{Method 1 — Introducing an Author, a Period and a Text (the "Context Paragraph")}

\begin{methode}[Writing a Context Paragraph for a Set Text]
GCE Paper 2 and 3 essays are rewarded when the candidate situates the text before analysing it. Follow these steps:
\begin{enumerate}
\item \textbf{Name the author, the work and the date/period} in one sentence (e.g. Chaucer and the late fourteenth century; Hansberry and 1950s America; Orwell and the post-1945 world).
\item \textbf{Give one or two biographical facts} that genuinely illuminate the work (Chaucer's career at court and as a civil servant; Hansberry's family experience of housing segregation in Chicago; Orwell's experience of totalitarianism and his democratic socialism).
\item \textbf{Identify the genre and structure} (a frame narrative of pilgrims telling tales; a three-act realist domestic drama; a dystopian novel in three parts).
\item \textbf{State the historical or social background} relevant to the question (feudal estates society; racial segregation and the American Dream; fascism, Stalinism and post-war surveillance).
\item \textbf{End with a link sentence} that turns the context towards the essay question.
\end{enumerate}
\textbf{Reflex:} context is a \emph{tool}, not a display. Every biographical or historical fact must serve the argument. Never write a whole page of biography.
\end{methode}

\begin{exoresolu}[Context Paragraph on The Canterbury Tales]
\textbf{Statement:} "A knowledge of Chaucer's life and times is essential to an appreciation of \emph{The Canterbury Tales}." Using this statement as your guide, write an introductory paragraph for an essay on the \emph{General Prologue}.
\tcblower
\textbf{Step 1 — Author, work, period.} Geoffrey Chaucer (c. 1340--1400), often called the father of English poetry, wrote \emph{The Canterbury Tales} in the last decades of the fourteenth century, at the close of the Middle Ages in England.

\textbf{Step 2 — Illuminating biography.} Chaucer was no cloistered scholar: he served as a courtier, diplomat, soldier and customs official. This varied career brought him into contact with every rank of society, from nobility to labourers, and explains the astonishing social range of his pilgrims.

\textbf{Step 3 — Genre and structure.} The work is a \emph{frame narrative}: about thirty pilgrims travelling from the Tabard Inn at Southwark to the shrine of Thomas Becket at Canterbury agree to tell stories to pass the time. The \emph{General Prologue} introduces the pilgrims one by one, so that the collection becomes a portrait of the whole of fourteenth-century English society — the three \emph{estates}: those who fight (Knight, Squire, Yeoman), those who pray (Prioress, Monk, Friar, Parson) and those who work (Ploughman, Miller, and the rising middle class of Merchants, Clerks and Lawyers).

\textbf{Step 4 — Background.} Chaucer writes at a moment of change: the feudal order is weakening, the Church is widely criticised for corruption, and a wealthy middle class is emerging. His method is \emph{irony} — he appears to praise while he exposes.

\textbf{Step 5 — Link.} It is precisely because Chaucer knew this society from the inside that his Prologue remains both a historical document and a living gallery of human types.

\textbf{Examiner's note:} five moves, eight to ten lines, no irrelevant detail. This is the model length for an A Level introduction.
\end{exoresolu}

\courssub{Method 2 — Analysing a Portrait in the General Prologue}

\begin{methode}[Reading a Chaucerian Portrait]
For any pilgrim (Knight, Squire, Prioress, Monk, Friar, Merchant, Clerk, Sergeant of Law, Franklin, Reeve, Miller):
\begin{enumerate}
\item \textbf{Classify the pilgrim} by estate and social rank (feudal/military, religious, professional middle class, rural/land managers).
\item \textbf{List the details Chaucer selects}: dress, horse, face, speech, habits, possessions. In Chaucer, \emph{appearance reveals character}.
\item \textbf{Decide the tone}: is the portrait sincere praise (Knight, Clerk, Parson, Ploughman) or \emph{ironic} praise (Prioress, Monk, Friar, Merchant)? Look for the gap between what the pilgrim \emph{should} be and what he or she \emph{is}.
\item \textbf{Identify the satirical target}: corruption of the Church, greed of professionals, vanity of social climbers.
\item \textbf{Conclude on function}: what does this portrait tell us about Chaucer's view of his society?
\end{enumerate}
\textbf{Reflex:} quote or closely paraphrase one concrete detail (the Monk's hunting horses, the Friar's gifts to young women, the Sergeant of Law's air of busyness) — examiners reward precise textual support.
\end{methode}

\begin{exoresolu}[Portrait Analysis: The Monk and the Friar]
\textbf{Statement:} "Chaucer's religious pilgrims are more devoted to worldly pleasure than to God." Discuss with close reference to the portraits of the Monk and the Friar in the \emph{General Prologue}.
\tcblower
\textbf{Step 1 — Classification.} Both belong to the religious estate, "those who pray", and both are expected to live in poverty, chastity and obedience.

\textbf{Step 2 — Selected details.} The \textbf{Monk} loves hunting, keeps fine horses and greyhounds, wears a fur-trimmed sleeve and a gold pin, and is fat and well-fed. He openly dismisses the old monastic rules as outdated: the cloister and the study of texts are not for him. The \textbf{Friar}, licensed to beg and hear confessions, is "easy" in granting absolution — for a price. He knows the taverns and barmaids better than the poor he is meant to serve, arranges marriages for young women he has himself compromised, and carries gifts to attract them.

\textbf{Step 3 — Tone.} The tone is \emph{ironic praise}. Chaucer calls the Monk "a manly man, to been an abbot able" and the Friar a "worthy" man — the very words of approval become accusations, because the qualities praised (hunting skill, charm, profitable confession) are the opposite of monastic virtue. The narrator pretends to agree with the Monk that old rules are dead; the reader is meant to disagree.

\textbf{Step 4 — Satirical target.} Chaucer attacks the corruption of the fourteenth-century Church: wealth, luxury, the sale of forgiveness, and the neglect of the poor. Significantly, he balances these portraits with the genuinely holy Parson, showing that he criticises corrupt churchmen, not religion itself.

\textbf{Step 5 — Function.} Together the portraits dramatise the gap between profession and practice. Chaucer's method — concrete detail plus ironic approval — makes the satire humorous rather than bitter, and all the more devastating.

\textbf{Conclusion.} The statement is justified: the Monk serves his appetite and the Friar serves his purse; in both, Chaucer exposes a Church that has forgotten its mission.
\end{exoresolu}

\courssub{Method 3 — Comparing Two Characters (the Paired Portrait)}

\begin{methode}[Writing a Comparison]
\begin{enumerate}
\item \textbf{Find the pairing principle} (e.g. Knight/Squire = father and son, ideal vs. youthful version of chivalry; Reeve/Miller = two kinds of rural dishonesty).
\item \textbf{Structure by criteria}, not by character: criterion 1 (values), criterion 2 (appearance), criterion 3 (Chaucer's tone) — discuss BOTH characters under each criterion.
\item \textbf{Use comparative connectives}: whereas, unlike, by contrast, similarly, while.
\item \textbf{Conclude on the meaning of the contrast} (e.g. the decline of chivalric ideals; the varieties of corruption).
\end{enumerate}
\textbf{Trap to avoid:} writing two separate mini-essays glued together. The examiner wants \emph{sustained comparison}.
\end{methode}

\begin{exoresolu}[Comparison: The Knight and the Squire]
\textbf{Statement:} Compare Chaucer's presentation of the Knight and the Squire in the \emph{General Prologue}. What does the contrast between father and son suggest about the chivalric ideal?
\tcblower
\textbf{Introduction.} The Knight and the Squire open the gallery of pilgrims and represent the feudal, military estate in two generations. Chaucer treats both with affection, but not identically.

\textbf{Criterion 1 — Values and experience.} The Knight is the \emph{ideal} of chivalry: he has fought in numerous campaigns across Christendom and beyond, yet he loves "trouthe and honour, freedom and curteisie". His worth is proved by a lifetime of service. The Squire, his son, is a \emph{beginner}: he has fought only briefly, and his service is motivated partly by the hope of winning his lady's favour. Where the father's chivalry is duty, the son's is entangled with love and display.

\textbf{Criterion 2 — Appearance.} The Knight's dress is strikingly modest: a plain tunic stained by his armour, for he has come straight from campaigning to pilgrimage. The Squire, by contrast, wears embroidered clothes "as it were a meede / Al ful of fresshe floures"; he sings, dances, plays the flute and sleeps little, like a nightingale in love. The contrast of dress dramatises the contrast of seriousness.

\textbf{Criterion 3 — Chaucer's tone.} Towards the Knight the tone is one of sincere, unqualified respect — he is "a verray, parfit gentil knyght", and there is no irony. Towards the Squire the tone is amused and tender: his vanity is the charming vanity of youth, not corruption. Unlike the Monk or the Friar, the Squire is not satirised; he is simply young.

\textbf{Conclusion.} The pairing suggests that the chivalric ideal survives, but must be earned by experience. The Knight embodies chivalry achieved; the Squire, chivalry in flower — promising, delightful, not yet tested. Chaucer thus honours the feudal ideal while gently observing how it softens in the next generation.
\end{exoresolu}

\courssub{Method 4 — Analysing a Dramatic Scene (A Raisin in the Sun)}

\begin{methode}[Scene Analysis in Drama]
For a scene such as Act 1 Scene 1 (the Younger family and the \$10,000 insurance cheque):
\begin{enumerate}
\item \textbf{Set the scene}: time, place, stage directions. Hansberry's opening directions (the worn but once-loved furniture, the single window) are themselves a statement about poverty and dignity.
\item \textbf{Identify the central conflict} and who embodies each side (Walter's dream of the liquor store vs. Mama's dream of a house; Ruth's exhaustion; Beneatha's search for identity).
\item \textbf{Trace the money motif}: the \$10,000 from Big Walter's life insurance means something different to each character — freedom (Walter), stability (Mama), medical school (Beneatha), relief (Ruth).
\item \textbf{Analyse dramatic techniques}: dialogue and interruption, stage directions, symbolism (the plant, the window, the eggs Ruth cooks), irony.
\item \textbf{Link to theme}: the deferred American Dream, the cost of racism on family relationships, the dignity of labour.
\end{enumerate}
\textbf{Reflex:} in drama, always comment on \emph{how} meaning is made on stage, not only on what is said.
\end{methode}

\begin{exoresolu}[Scene Analysis: Act 1 Scene 1 of A Raisin in the Sun]
\textbf{Statement:} "In Act 1 Scene 1, Hansberry uses an ordinary morning in the Younger household to present a whole world of frustration and hope." Discuss.
\tcblower
\textbf{Step 1 — Setting.} The stage directions describe a living room whose furniture was once chosen with love but is now worn from too much use; a single small window provides the only natural light. Before any character speaks, the set announces the family's condition: dignity maintained against poverty, and lives cramped in a space too small for their dreams.

\textbf{Step 2 — Conflicts.} The morning routine releases conflict after conflict. Walter Lee asks about the insurance cheque and clashes with Ruth, who refuses to support his plan to invest in a liquor store; he complains that the women of his family do not understand his ambition. Beneatha quarrels with Walter over her medical studies, which he sees as a waste of the family's money. Mama enters as the moral centre, suspicious of liquor money and dreaming instead of a house with a garden.

\textbf{Step 3 — The money motif.} The expected \$10,000 — the price of Big Walter's life — crystallises every dream. Mama's grief that her husband's life should be converted into cash gives the cheque a tragic weight: the family's hope is literally built on a death. Walter's famous question about why Ruth cooks eggs every morning, and his talk of "coloured people" who never get anywhere, show a man measuring his manhood by money.

\textbf{Step 4 — Techniques.} Hansberry uses rapid, overlapping domestic dialogue to create realism; symbolism (Mama's feeble plant struggling towards the little light mirrors the family itself); and irony — the cheque that should unite the family first divides it.

\textbf{Step 5 — Theme.} The scene dramatises the \emph{deferred dream}: each Younger has a dream postponed by race and poverty. The ordinary morning thus opens onto the largest questions of the play — what money can and cannot buy, and what a family owes to its dead and to its future.
\end{exoresolu}

\courssub{Method 5 — Writing a Thematic Essay on Drama}

\begin{methode}[Thematic Essay: Gender, Identity and the American Dream]
For questions on Act 1 Scene 2 and Act 2 Scene 1 (Beneatha, Asagai, George Murchison; Walter and the American Dream):
\begin{enumerate}
\item \textbf{Define the theme precisely} in the introduction (e.g. the American Dream = the belief that hard work brings success; Hansberry tests this belief against racial reality).
\item \textbf{Organise by aspects of the theme}, each illustrated by a scene and a character.
\item \textbf{Use the two suitors as a device}: Asagai (African heritage, pride, "Alaiyo" — one for whom bread is not enough) vs. George Murchison (assimilation, wealth, contempt for African roots). Beneatha's choice between them is a debate about Black identity.
\item \textbf{Include gender}: Beneatha's ambition to be a doctor challenges both racial and gender expectations; Walter's crisis is a crisis of masculinity in a matriarchal household.
\item \textbf{Conclude with the playwright's vision}: dreams must be rooted in dignity and heritage, not merely in money.
\end{enumerate}
\end{methode}

\begin{exoresolu}[Thematic Essay: The American Dream in A Raisin in the Sun]
\textbf{Statement:} "Hansberry's play is both a celebration of dreams and a warning about money." Discuss with reference to Act 1 Scene 2 and Act 2 Scene 1.
\tcblower
\textbf{Introduction.} The play's title, drawn from Langston Hughes's poem about a dream deferred, announces its subject. Hansberry celebrates the Youngers' capacity to dream, while warning that a dream reduced to money alone destroys the dreamer.

\textbf{Paragraph 1 — The dream as dignity.} In Act 1 Scene 2, Mama's dream of a house, Travis's need for space, and Beneatha's medical ambition show dreams tied to family and self-respect. Mama's plant, kept alive in a dark kitchen, symbolises dreams that survive without ideal conditions.

\textbf{Paragraph 2 — The dream corrupted by money.} Walter equates manhood with financial success. In Act 2 Scene 1, his despair and his bitterness towards his family show a man who has allowed money to define his worth. Hansberry warns through him: when the dream is only money, its deferral produces self-contempt.

\textbf{Paragraph 3 — Identity and heritage.} Act 2 Scene 1 stages the debate about cultural identity. Beneatha appears in Nigerian dress, and Asagai's gift of the robes and his nickname for her invite her to root her identity in Africa; George Murchison, wealthy and assimilated, mocks African heritage and asks her to be "quiet" and decorative. The scene argues that money without roots — George's kind of success — is another form of poverty.

\textbf{Paragraph 4 — Gender.} Beneatha refuses to be defined by marriage; her dream of medicine asserts a woman's right to a self. Hansberry thus widens the American Dream to include Black women.

\textbf{Conclusion.} The play celebrates dreaming as an act of resistance, but insists that the true dream is dignity — a house, a heritage, a self — and not merely the cheque.
\end{exoresolu}

\courssub{Method 6 — Reviewing a Novel's Opening Movement (Nineteen Eighty-Four, Part One)}

\begin{methode}[Writing a Review of Part One]
\begin{enumerate}
\item \textbf{Establish the world}: Oceania, Airstrip One (London), the Party and Big Brother, the telescreens, the Thought Police, the Ministries (Truth, Peace, Love, Plenty) — note the ironic names.
\item \textbf{Present the protagonist}: Winston Smith, 39, works at the Ministry of Truth rewriting history; outwardly obedient, inwardly rebellious.
\item \textbf{Trace the rebellion's first steps}: the diary bought in a junk shop, the opening words of April 4th 1984, the writing of "DOWN WITH BIG BROTHER"; his memory of the past; his suspicion of O'Brien and his hatred of the dark-haired girl.
\item \textbf{Identify the key ideas}: control of the past ("Who controls the past controls the future"), doublethink, Newspeak, perpetual war, the Two Minutes Hate, the proles as the only hope.
\item \textbf{Assess tone and method}: Orwell's plain, grey prose mirrors the deadened world; the diary gives us Winston's inner voice.
\end{enumerate}
\textbf{Reflex:} a review is not a summary alone — end with a judgement of Part One's \emph{function} (it builds the prison so that Part Two can stage the escape attempt).
\end{methode}

\begin{exoresolu}[Review of Part One of Nineteen Eighty-Four]
\textbf{Statement:} Write a review of Part One of \emph{Nineteen Eighty-Four}, showing how Orwell establishes his dystopian world and his hero's rebellion.
\tcblower
\textbf{The world.} Part One opens on a cold April day as Winston Smith enters Victory Mansions. Orwell builds Oceania through accumulated detail: posters of Big Brother whose eyes follow you, telescreens that watch and broadcast simultaneously, the Thought Police, the daily Two Minutes Hate against the traitor Goldstein, and the four Ministries whose names invert their functions — the Ministry of Truth falsifies records, the Ministry of Peace wages perpetual war. The Party's slogans (War is Peace; Freedom is Slavery; Ignorance is Strength) define \emph{doublethink}: holding two contradictory beliefs at once.

\textbf{The hero.} Winston, thirty-nine, with a varicose ulcer and a failing marriage behind him, is an ordinary man, not a conventional hero. His job — rewriting old newspapers so that the Party is always right — makes him an accomplice in the very lies he hates. This is Orwell's point: totalitarianism survives by forcing everyone to collaborate.

\textbf{The rebellion.} Winston's first crime is \emph{thought}: he buys a beautiful old diary and begins to write, though he knows the penalty. He writes "DOWN WITH BIG BROTHER" almost without willing it. His rebellion rests on memory — he suspects the past has been altered — and on two hopes: that the \emph{proles}, the unpoliced masses, may one day rise ("If there is hope, it lies in the proles"), and that O'Brien, an Inner Party member, may secretly share his hatred. He also notices a dark-haired girl watching him and fears she is a spy — an irony Part Two will exploit.

\textbf{Function of Part One.} Part One constructs the prison: a world where truth, memory, love and even private thought are crimes. Its closing mood is of a man already half-defeated, writing to a future that may not exist. Orwell's achievement is to make the political personal: we measure the horror of the system by its cost to one ordinary mind. Part One thus prepares the essential question of Part Two — can love exist where everything is watched?
\end{exoresolu}

\courssub{Method 7 — Analysing a Developing Narrative (Part Two, Chapters 1--6: Love as Rebellion)}

\begin{methode}[Analysing the Winston--Julia Episode]
\begin{enumerate}
\item \textbf{Track the plot precisely}: the note "I love you"; the first meeting in the crowd; the countryside rendezvous (the "Golden Country" of Winston's dreams); the room above Mr Charrington's shop; Julia's story; the reading of Goldstein's book begins; the approach to O'Brien.
\item \textbf{Contrast the lovers}: Winston rebels intellectually (he wants truth and the past); Julia rebels instinctively and sensually (she wants pleasure and small freedoms). Neither is complete without the other.
\item \textbf{Read the symbols}: the Golden Country = freedom and natural love; the coral paperweight = the fragile, beautiful past enclosed in glass; the room above the shop = a private world that is, unknown to them, already a trap.
\item \textbf{Analyse the politics of love}: in Oceania, private loyalty is treason; the Party wants all love directed to Big Brother. Hence "their embrace had been a battle, the climax a victory... a political act."
\item \textbf{Note the foreshadowing}: the telescreen behind the picture, Charrington's strange knowledge, the rhyme about the bells — Orwell plants the seeds of the arrest.
\end{enumerate}
\end{methode}

\begin{exoresolu}[Analysis: Part Two, Chapters 1--6 of Nineteen Eighty-Four]
\textbf{Statement:} "In Part Two of \emph{Nineteen Eighty-Four}, love becomes Winston's most dangerous crime." Discuss with close reference to Chapters 1 to 6.
\tcblower
\textbf{Introduction.} Part One ended with Winston's rebellion confined to a diary. In Part Two, rebellion takes flesh: his affair with Julia transforms private thought into shared action — the gravest offence against a Party that demands total, exclusive loyalty.

\textbf{The beginning of the affair (Chapters 1--2).} The dark-haired girl Winston feared as a spy slips him a note: "I love you." Their first meeting must be staged in a crowd, for in Oceania two people alone together are already suspect. In the countryside — the landscape of Winston's dream of the Golden Country — they become lovers. Orwell is explicit about the meaning: their embrace is a political act, a blow struck against the Party's control of the body. Julia's practical cunning (she has done this before; she knows how to evade the patrols) contrasts with Winston's anxious idealism.

\textbf{The room above the shop (Chapters 3--4).} Winston rents the room above Mr Charrington's junk shop, believing it safe. There the lovers create a fragile private world: real coffee, sugar, bread, a bed — ordinary things that have become luxuries and therefore crimes. The coral paperweight Winston has bought becomes the symbol of this world: a tiny, beautiful past preserved inside glass, as the lovers are preserved inside the room. Julia, unlike Winston, has no interest in overthrowing the Party; she rebels for pleasure and personal freedom. Winston loves her for exactly this vitality, yet their difference plants a question: is instinctive rebellion enough?

\textbf{The deepening commitment (Chapters 5--6).} As the affair continues, Winston grows healthier and more hopeful; he and Julia discuss the Brotherhood and agree to approach O'Brien, whose face and manner have convinced Winston that he is a secret enemy of the Party. The reader, however, notices what the lovers ignore: Charrington knows too much, the picture on the wall conceals something, and the old rhyme of the London bells ends with a threat. Orwell's foreshadowing turns every moment of happiness into suspense.

\textbf{Conclusion.} In these chapters Orwell makes his most radical claim: under totalitarianism, love itself is subversion, because it creates a loyalty the Party does not control. But the symbols — the glass around the coral, the hidden eye behind the picture — warn that the lovers' private world is already enclosed. Their love is a victory; the question the rest of the novel will answer is how long a victory can survive inside a trap.
\end{exoresolu}

\begin{attention}[Common Mistakes in the GCE Literature Examination]
\begin{itemize}
\item \textbf{Narrating instead of analysing.} Retelling the plot of the Prologue, the play or the novel earns few marks; every reference must support an argument.
\item \textbf{Quoting inaccurately.} If you are not sure of the exact words, paraphrase closely and name the scene or chapter.
\item \textbf{Ignoring the genre.} Do not treat Hansberry's stage directions as decoration or Chaucer's irony as literal praise.
\item \textbf{Context without purpose.} Biography and history must serve the question, not replace the answer.
\item \textbf{Unbalanced comparisons.} In paired questions (Knight/Squire, Asagai/George), discuss both sides under each point.
\end{itemize}
\end{attention}

\newpage

\coursec{Exercises}

\courssub{I check my knowledge}

\begin{exercice}[Multiple choice: authors and texts]
Choose the correct option for each question.
\begin{enumerate}
\item Geoffrey Chaucer is traditionally known as:
\begin{enumerate}
\item the father of the English novel
\item the father of English poetry
\item the first English playwright
\item the inventor of the sonnet
\end{enumerate}
\item \emph{A Raisin in the Sun} takes its title from:
\begin{enumerate}
\item a sermon by Martin Luther King
\item a poem by Langston Hughes
\item a speech by Malcolm X
\item a Negro spiritual
\end{enumerate}
\item In \emph{Nineteen Eighty-Four}, the official language of Oceania is:
\begin{enumerate}
\item Oldspeak
\item Doublespeak
\item Newspeak
\item Party-speak
\end{enumerate}
\item The pilgrims in \emph{The Canterbury Tales} are travelling to the shrine of:
\begin{enumerate}
\item Saint Augustine at Canterbury
\item Thomas Becket at Canterbury
\item Saint George at Windsor
\item the Virgin Mary at Walsingham
\end{enumerate}
\end{enumerate}
\end{exercice}

\begin{exercice}[True or false — justify your answer]
Say whether each statement is true or false, and justify briefly with reference to the text.
\begin{enumerate}
\item The Knight in the General Prologue is portrayed with gentle irony as a hypocrite.
\item Walter Lee Younger wants to invest the insurance money in a liquor store.
\item Julia, in \emph{Nineteen Eighty-Four}, is a sincere ideological rebel who hates the Party for political reasons.
\item The Friar is presented as a model of poverty and humility.
\item Beneatha Younger is searching for an African cultural identity.
\end{enumerate}
\end{exercice}

\begin{exercice}[Key definitions]
Define each of the following literary terms in two or three sentences, and give one example from any of the three set texts:
\begin{enumerate}
\item \cle{Dramatic irony}
\item \cle{Estates satire}
\item \cle{Dystopia}
\item \cle{Symbolism}
\item \cle{Protagonist}
\end{enumerate}
\end{exercice}

\begin{exercice}[Matching: pilgrims, characters and traits]
Match each name in column A to the correct description in column B.
\begin{center}
\begin{tabularx}{\linewidth}{|l|X|}
\hline
\rowcolor{popblueL} \textbf{Column A} & \textbf{Column B} \\
\hline
1. The Squire & a) A corrupt churchman who sells pardons and loves money \\
\hline
2. The Monk & b) A young lover, poet and musician, fond of finery \\
\hline
3. The Friar & c) A hunter who neglects his cloister and religious rule \\
\hline
4. Mama (Lena Younger) & d) A wanton, merry man who arranges marriages for payment \\
\hline
5. Big Brother & e) The moral anchor of the Younger family, defender of dignity \\
\hline
6. O'Brien & f) The ever-watchful face of the Party in Oceania \\
\hline
& g) The Inner Party member who interrogates Winston \\
\hline
\end{tabularx}
\end{center}
(One description in column B is a distractor and matches nobody.)
\end{exercice}

\courssub{I practise}

\begin{exercice}[Portrait analysis: the Knight]
Re-read the portrait of the Knight in the General Prologue. In about 250 words:
\begin{enumerate}
\item identify three details (dress, behaviour, career) that Chaucer gives about the Knight;
\item explain what each detail reveals about his character and values;
\item state whether the portrait is sincere or ironic, justifying your answer.
\end{enumerate}
\end{exercice}

\begin{exercice}[Irony in the religious portraits]
Choose \textbf{one} of the following pilgrims: the Prioress, the Monk or the Friar.
\begin{enumerate}
\item Quote or paraphrase two details of the portrait that contradict the pilgrim's religious calling.
\item Explain the ironic gap between what the pilgrim \emph{should} be and what he or she \emph{is}.
\item Conclude on the effect of Chaucer's irony on the reader.
\end{enumerate}
Write one well-structured paragraph of about 200 words.
\end{exercice}

\begin{exercice}[The \$10,000 dream]
In Act 1 Scene 1 of \emph{A Raisin in the Sun}, each member of the Younger family has a different plan for the insurance money.
\begin{enumerate}
\item Draw a table with three columns: \emph{Character}, \emph{Dream for the money}, \emph{What the dream reveals about the character}. Complete it for Mama, Walter, Beneatha and Ruth.
\item In ten lines, explain why the money creates conflict in the family.
\end{enumerate}
\end{exercice}

\begin{exercice}[Newspeak and doublethink in practice]
\begin{enumerate}
\item Explain in your own words the Party slogan: \emph{War is Peace. Freedom is Slavery. Ignorance is Strength.}
\item Give one concrete example of doublethink from Part One of \emph{Nineteen Eighty-Four} and show how it controls Winston's mind.
\item Explain the purpose of the telescreen in the daily life of Party members.
\end{enumerate}
Answer in three short paragraphs.
\end{exercice}

\courssub{I prepare for the GCE}

\begin{exercice}[GCE-style essay: Chaucer's mirror of society]
\emph{``Chaucer's General Prologue is a mirror of his society.''} Discuss this statement with close reference to \textbf{at least four pilgrims} studied this term (for example the Knight, the Prioress, the Monk, the Friar, the Merchant, the Clerk, the Franklin or the Miller). In your answer, consider the three estates, the mixture of sincerity and satire, and Chaucer's narrative technique. \textbf{(20 marks)}
\end{exercice}

\begin{exercice}[GCE-style essay: money, dreams and identity]
\emph{``Money is life.''} Examine Walter Lee Younger's philosophy in \emph{A Raisin in the Sun} and show how Hansberry challenges it through Mama and Beneatha. Your essay should refer to Act 1 Scenes 1 and 2 and to Act 2 Scene 1, and should discuss the American Dream, gender and cultural identity. \textbf{(20 marks)}
\end{exercice}

\begin{exercice}[GCE-style guided commentary and comparative question]
\textbf{Part A (10 marks).} Comment on the passage in Part Two, Chapter 2 of \emph{Nineteen Eighty-Four} where Winston and Julia meet in the Golden Country. Bring out the theme of love as rebellion, the symbolism of nature and the thrush, and Orwell's narrative technique.

\textbf{Part B (10 marks).} \emph{``In all three texts, individuals struggle against the pressure of their society.''} Evaluate this statement using \emph{The Canterbury Tales}, \emph{A Raisin in the Sun} and \emph{Nineteen Eighty-Four}. Support your argument with precise references to at least one character or pilgrim from each text. \textbf{(20 marks in total)}
\end{exercice}

\newpage

\coursec{Solutions to the exercises}

\begin{corrige}[Exercise 1 — Multiple choice: authors and texts]
\begin{enumerate}
\item \textbf{Correct answer: b) the father of English poetry.} Chaucer earned this title (later popularised by Dryden) because he was the first major writer to compose ambitious, sophisticated verse in Middle English at a time when Latin and French dominated literature. Option a) is wrong because the English novel develops in the eighteenth century (Defoe, Richardson); c) is wrong because drama flourished later; d) is wrong because the sonnet was imported from Italy by Wyatt and Surrey in the sixteenth century.
\item \textbf{Correct answer: b) a poem by Langston Hughes.} The title comes from Hughes's poem \emph{Harlem} (also called \emph{A Dream Deferred}), which asks what happens to a deferred dream and wonders whether it dries up ``like a raisin in the sun''. The poem is quoted as the epigraph of the play.
\item \textbf{Correct answer: c) Newspeak.} Newspeak is the official engineered language of Oceania, designed to make ``thoughtcrime'' impossible by destroying words. \emph{Oldspeak} is the name given to standard English; \emph{doublespeak} is a related but unofficial term.
\item \textbf{Correct answer: b) Thomas Becket at Canterbury.} The pilgrims travel to the shrine of Saint Thomas Becket, Archbishop of Canterbury, murdered in the cathedral in 1170 and venerated as a martyr. Chaucer calls him ``the holy blissful martyr'' whom they thank for help in sickness.
\end{enumerate}
\end{corrige}

\begin{corrige}[Exercise 2 — True or false]
\begin{enumerate}
\item \textbf{False.} The Knight is one of the very few pilgrims portrayed \emph{without} irony. Chaucer presents him as a genuine ideal: ``a worthy man'', a perfect ``gentil knight'' who loves chivalry, truth, honour and courtesy. The gentle irony is reserved for corrupt pilgrims such as the Monk and the Friar.
\item \textbf{True.} In Act 1 Scene 1, Walter Lee tells Ruth of his plan to invest part of the \$10,000 insurance money in a liquor store with his friends Willy and Bobo, believing this business will lift the family out of poverty. Ruth and Mama are horrified by the idea.
\item \textbf{False.} Julia is a rebel ``from the waist downwards'', as Winston observes: her rebellion is personal and sensual, not ideological. She enjoys outwitting the Party for her own pleasure and falls asleep during Winston's political reflections; she does not share his intellectual hatred of the regime.
\item \textbf{False.} The Friar is the opposite of a model of poverty and humility: he is ``wanton and merry'', knows the taverns and barmaids better than the lepers he should serve, arranges marriages (for a fee) for young women he has made pregnant, and is an excellent beggar who grows rich from his begging district.
\item \textbf{True.} Beneatha rejects assimilation, cuts her chemically straightened hair to wear it natural, studies African culture with Joseph Asagai, who nicknames her ``Alaiyo'', and dreams of reconnecting with her African heritage — a clear search for cultural identity.
\end{enumerate}
\end{corrige}

\begin{corrige}[Exercise 3 — Key definitions]
\begin{enumerate}
\item \textbf{Dramatic irony:} a device by which the audience or reader knows more about a situation than a character does, so that the character's words or actions acquire a meaning he or she does not intend. \emph{Example:} in \emph{Nineteen Eighty-Four}, the reader suspects from early on that O'Brien is loyal to the Party, while Winston trusts him as a fellow rebel — a trust that leads to his arrest.
\item \textbf{Estates satire:} a medieval genre in which the writer surveys the different ``estates'' or social classes (nobility, clergy, commoners) and exposes how far their members fall short of the ideals of their calling. \emph{Example:} Chaucer's General Prologue, where the corrupt Monk and Friar are measured against the ideal of the religious life.
\item \textbf{Dystopia:} an imagined society, usually set in the future, in which life is degraded by oppression, surveillance or dehumanisation; it is the opposite of a utopia and serves as a warning. \emph{Example:} Orwell's Oceania, ruled by Big Brother through the telescreens, Newspeak and the Thought Police.
\item \textbf{Symbolism:} the use of a concrete object, place or person to represent an abstract idea. \emph{Example:} Mama's little plant in \emph{A Raisin in the Sun} symbolises her persistent dream of a garden and a better life for her family despite poor light and cramped conditions.
\item \textbf{Protagonist:} the central character of a literary work, around whose desires, conflicts and development the plot is organised. \emph{Example:} Winston Smith is the protagonist of \emph{Nineteen Eighty-Four}; Walter Lee Younger can be considered the protagonist of \emph{A Raisin in the Sun}.
\end{enumerate}
\end{corrige}

\begin{corrige}[Exercise 4 — Matching: pilgrims, characters and traits]
\begin{center}
\begin{tabularx}{\linewidth}{|l|X|X|}
\hline
\rowcolor{popblueL} \textbf{Column A} & \textbf{Match} & \textbf{Justification} \\
\hline
1. The Squire & b) & A young ``lover and lusty bachelor'' of about twenty, who sings, dances, writes songs and wears embroidered clothes. \\
\hline
2. The Monk & c) & He loves hunting, keeps fine horses and greyhounds, and dismisses the old monastic rule as outdated. \\
\hline
3. The Friar & d) & A ``wanton and merry'' limiter who performs marriages ``at his own cost'' and is intimate with taverners. \\
\hline
4. Mama (Lena Younger) & e) & She holds the family together, defends its dignity and Christian values, and buys the house in Clybourne Park. \\
\hline
5. Big Brother & f) & His moustached face stares from every poster and telescreen with the caption ``Big Brother is watching you''. \\
\hline
6. O'Brien & g) & The Inner Party member who pretends to be a rebel, then supervises Winston's interrogation in the Ministry of Love. \\
\hline
\end{tabularx}
\end{center}
\textbf{Distractor:} description a) — ``a corrupt churchman who sells pardons and loves money'' — describes the \emph{Pardoner}, who is not listed in column A, so it matches nobody.
\end{corrige}

\begin{corrige}[Exercise 5 — Portrait analysis: the Knight]
\textbf{Model answer (about 250 words).}

Chaucer builds the Knight's portrait through three types of detail. First, his \textbf{career}: he has fought in numerous campaigns across the Christian and pagan worlds — Alexandria, Prussia, Lithuania, Granada, North Africa — always ``in his lord's war''. This long list of crusading battles reveals a life entirely devoted to the chivalric and Christian ideal of the warrior; he fights for faith and lord, not for personal gain. Second, his \textbf{behaviour}: though he has always won ``sovereign honour'', he is ``as meek as is a maid'', never speaks rudely to anyone, and loves ``chivalry, truth, honour, freedom and courtesy''. His worthiness is matched by perfect modesty — the true mark of the medieval \emph{gentil} man, whose nobility is moral as much as social. Third, his \textbf{dress}: he wears a simple tunic of coarse fustian, stained by his chain-mail, and comes straight from his voyage to give thanks at Canterbury. His shabby clothing and lack of ostentation contrast deliberately with the fine horses and fur-trimmed sleeves of the Monk: the Knight possesses real worth and does not need to display it.

The portrait is essentially \textbf{sincere}. Unlike the religious pilgrims, the Knight is measured against the ideal of his estate and found equal to it; there is no ironic gap between what he should be and what he is. Placed first in the Prologue, he sets the moral standard by which the corruption of the other pilgrims will be judged. Any irony is at most a light smile at his old-fashioned perfection, not satire.
\end{corrige}

\begin{corrige}[Exercise 6 — Irony in the religious portraits]
\textbf{Model paragraph (the Monk; about 200 words).}

Chaucer's portrait of the Monk is built on a systematic contradiction between the monastic ideal and the man himself. A monk should live in poverty, chastity and obedience, enclosed in his cloister, devoted to prayer and manual labour. This Monk, however, is ``an outrider, that loved venery'' — a passionate hunter who keeps many ``dainty'' horses and greyhounds, whose bridle jingles like a chapel bell, and who wears sleeves trimmed with expensive grey fur and a gold pin. He is fat, bald and greasy-faced, fond of roast swan, and openly scorns the ancient rule of Saint Benedict, calling the text that says hunters are not holy men ``not worth an oyster''. The ironic gap is glaring: everything the Rule forbids — luxury, hunting, fine food, life outside the cloister — is exactly what this Monk loves and practises. Chaucer's technique is to let the Monk condemn himself in his own words while the narrator pretends to agree with him (``And I said his opinion was good''), a mock-approval that sharpens the satire. The effect on the reader is double: we laugh at the Monk's hearty worldliness, but we also recognise the corruption of a Church whose officials have abandoned their spiritual mission for comfort and pleasure. Irony thus turns a comic portrait into social criticism.
\end{corrige}

\begin{corrige}[Exercise 7 — The \$10,000 dream]
\textbf{1. Table of dreams.}
\begin{center}
\begin{tabularx}{\linewidth}{|l|X|X|}
\hline
\rowcolor{popblueL} \textbf{Character} & \textbf{Dream for the money} & \textbf{What it reveals} \\
\hline
Mama (Lena) & Buy a house with a garden (she later puts a down payment on a house in Clybourne Park); keep money for Beneatha's studies. & Her selflessness, her faith, her dream of dignity and rootedness shared with her late husband. \\
\hline
Walter Lee & Invest in a liquor store with Willy and Bobo. & His hunger for status and independence; he equates money with manhood and life itself. \\
\hline
Beneatha & Use part of it for her medical school tuition. & Her ambition, her refusal of traditional female roles, her desire to heal and to be somebody. \\
\hline
Ruth & Escape the cramped, worn-out apartment; a better home for Travis and her unborn child. & Her quiet desperation, her longing for rest, light and family stability. \\
\hline
\end{tabularx}
\end{center}

\textbf{2. Why the money creates conflict (about ten lines).}

The \$10,000 comes from Big Walter's death, so it is charged with emotion: for Mama it is made of her husband's flesh and blood, and spending it on a liquor store feels like a betrayal of his memory and of Christian morality. Each character projects onto the money a different vision of freedom — property and respectability for Mama, business and manhood for Walter, education for Beneatha, comfort for Ruth — and these dreams are incompatible because the sum cannot satisfy them all. Moreover, the money legally belongs to Mama, which humiliates Walter, who feels robbed of his role as head of the family; gender and generational tensions therefore crystallise around it. The cheque, which should be a blessing, becomes a test of the family's unity and of each member's values.
\end{corrige}

\begin{corrige}[Exercise 8 — Newspeak and doublethink in practice]
\textbf{1. The Party slogans.} Each slogan is a paradox that reverses common sense. \emph{War is Peace}: permanent war against Eurasia or Eastasia keeps the population frightened, united and obedient, so external war produces internal ``peace'' for the regime. \emph{Freedom is Slavery}: a person left free to choose becomes the slave of his own desires and errors; only submission to the Party gives true order. \emph{Ignorance is Strength}: the less the masses know, the more powerful the Party is, because no one can compare the present with the past or imagine an alternative. The slogans train citizens to accept contradictions without protest.

\textbf{2. An example of doublethink.} At the Ministry of Truth, Winston rewrites old newspapers so that the Party's predictions always appear correct: he knows he is falsifying the past, yet he must simultaneously believe that the Party has never been wrong. Holding these two contradictory truths at once — knowing the record is forged and believing the forged record — is doublethink. It controls his mind by destroying the very idea of objective truth: if the past is whatever the Party says, resistance becomes literally unthinkable.

\textbf{3. The purpose of the telescreen.} The telescreen is a two-way device that broadcasts Party propaganda continuously while watching and listening to Party members at every moment; it can never be switched off, only dimmed. It creates permanent surveillance: any facial expression, gesture or murmur may betray ``thoughtcrime''. Its deeper purpose is psychological — since citizens never know when they are being observed, they learn to police themselves, internalising the Party's gaze until obedience becomes automatic.
\end{corrige}

\begin{corrige}[Exercise 9 — GCE-style essay: Chaucer's mirror of society]
\textbf{Indicative mark scheme (20 marks).}
\begin{itemize}
\item Introduction and thesis (2 marks): definition of the General Prologue as estates satire; clear position on the statement.
\item Knowledge of the text (6 marks): accurate reference to at least four pilgrims with precise details.
\item Analysis (8 marks): treatment of the three estates, the balance of sincerity and satire, narrative technique (naive narrator, mock-approval, telling detail).
\item Organisation, expression and conclusion (4 marks): coherent paragraphs, accurate English, personal judgement.
\end{itemize}

\textbf{Examiner's expectations and model plan.}

\emph{Introduction.} In the General Prologue, Chaucer assembles at the Tabard Inn a cross-section of fourteenth-century English society, ``of sundry folk by adventure y-fall in fellowship''. The statement is largely true: the Prologue functions as a mirror through estates satire — but it is a selective, artful mirror that judges as much as it reflects.

\emph{Paragraph 1 — The feudal estate faithfully reflected.} The Knight, the Squire and the Yeoman represent the military order. The Knight embodies its ideal (chivalry, truth, honour); the Squire shows its youthful, courtly side. Chaucer reflects this estate with sincerity, setting a standard of measurement.

\emph{Paragraph 2 — The religious estate satirised.} The Prioress (more concerned with table manners, her little dogs and her brooch than with charity), the Monk (hunter, fur-trimmed sleeves, scornful of the Rule) and the Friar (wanton, mercenary, a better beggar than confessor) reveal a Church that has abandoned poverty, chastity and obedience. The mirror here is satirical: the gap between ideal and reality exposes corruption.

\emph{Paragraph 3 — The rising middle class and rural world.} The Merchant with his forked beard and hidden debts, the impoverished Oxford Clerk who would rather have books than riches, the pleasure-loving Franklin and the dishonest, thieving Miller reflect new social energies: money, learning, comfort and rural cunning. Society is shown in motion, no longer fixed in three neat estates.

\emph{Paragraph 4 — Technique: how the mirror is made.} Chaucer uses a naive pilgrim-narrator who admires almost everyone, forcing the reader to detect the irony; mock-approval (``he was a worthy man'' applied dubiously); concrete, telling physical detail; and the hierarchy of the portraits, from the ideal Knight downwards.

\emph{Conclusion.} The Prologue mirrors medieval society in its ranks, professions and values, but the reflection is angled: through irony Chaucer invites us to judge each pilgrim against the ideal of his estate. The mirror is also a measure.

\textbf{Common mistakes to avoid:} retelling the portraits without analysis; treating all pilgrims as satirical (the Knight is sincere); ignoring narrative technique, which the question explicitly demands.
\end{corrige}

\begin{corrige}[Exercise 10 — GCE-style essay: money, dreams and identity]
\textbf{Indicative mark scheme (20 marks).}
\begin{itemize}
\item Introduction and thesis (2 marks): presentation of Walter's philosophy and of the question's stakes.
\item Knowledge of the text (6 marks): precise reference to Act 1 Scenes 1 and 2 and Act 2 Scene 1.
\item Analysis (8 marks): the American Dream, gender, cultural identity; how Mama and Beneatha challenge Walter.
\item Organisation, expression and conclusion (4 marks).
\end{itemize}

\textbf{Examiner's expectations and model plan.}

\emph{Introduction.} Walter Lee Younger, a chauffeur in Chicago's Southside, declares to Mama that ``money is life''. Hansberry stages this philosophy in order to test it: through Mama's dignity and Beneatha's idealism, the play asks whether money can really buy freedom, manhood and identity.

\emph{Paragraph 1 — Walter's philosophy and its roots.} In Act 1 Scene 1, Walter sees the insurance money as his one chance to escape servitude: the liquor store would make him a businessman, able to give Travis the future he never had. His equation of money with life is the American Dream in its materialist form — success measured in dollars — sharpened by the humiliation of serving rich white men.

\emph{Paragraph 2 — Mama's challenge: money against dignity.} Mama refuses to see money as life: she recalls that freedom itself was once the dream, and that Big Walter's life was worth more than the cheque his death produced. Her down payment on the Clybourne Park house (Act 1 Scene 2) redefines the dream as home, family and self-respect — values no price can express.

\emph{Paragraph 3 — Beneatha's challenge: identity beyond money.} Beneatha dreams of medicine, of healing, and of recovering an African identity with Asagai; she rejects George Murchison's assimilationist wealth. In Act 2 Scene 1, her natural hair and African dress dramatise a richness that is cultural, not financial. She tells Walter that money is not everything — though the play honestly shows her tuition also depends on it.

\emph{Paragraph 4 — Hansberry's verdict.} The play does not deny that money matters: poverty crushes Ruth and twists Walter. But Walter's loss of the money and his final refusal of Lindner's buy-out show that dignity must govern money, not the reverse. Manhood, as Mama says, is measured by how a man stands, not by what he owns.

\emph{Conclusion.} Walter's ``money is life'' is exposed as a half-truth born of oppression. Hansberry answers it with a fuller dream — one that joins material need to dignity, family and cultural pride.

\textbf{Common mistakes to avoid:} summarising the plot; ignoring Beneatha; forgetting Act 2 Scene 1, which the question explicitly requires.
\end{corrige}

\begin{corrige}[Exercise 11 — GCE-style guided commentary and comparative question]
\textbf{Part A — Commentary on the Golden Country (10 marks).}

\emph{Indicative mark scheme:} identification of context and themes (3); analysis of symbolism (3); narrative technique (2); organisation and expression (2).

\emph{Model answer.} The passage occurs in Part Two, Chapter 2, when Winston and Julia meet secretly in the countryside, in the clearing Winston has dreamed of as ``the Golden Country''. The theme of \textbf{love as rebellion} is central: in Oceania, the Party has abolished private loyalty and made sex a mere ``duty to the Party'', so a freely chosen love affair is a political act. As Winston understands, their embrace is ``a blow struck against the Party'' — desire itself becomes resistance. The \textbf{symbolism of nature} reinforces this: the open meadow, the bluebells and the stream form a space outside the telescreen's reach, an image of the natural, spontaneous life the Party denies. The \textbf{thrush}, which sings freely and beautifully for no audience and no purpose, symbolises useless, gratuitous joy — everything totalitarianism cannot tolerate; its song suggests that beauty and freedom survive beyond ideology. Orwell's \textbf{narrative technique} is significant: the third-person limited narration stays close to Winston's consciousness, so that the dreamlike, lyrical tone of the scene contrasts violently with the grey, mechanical style of the London chapters. The pastoral interlude is deliberately fragile: the reader, aware of the Thought Police, senses that this Eden cannot last, and the beauty of the passage measures the depth of the coming loss.

\textbf{Part B — Comparative essay (10 marks).}

\emph{Indicative mark scheme:} thesis and comparative structure (2); accurate treatment of each text (2 each = 6); expression and conclusion (2).

\emph{Model plan and answer sketch.} The statement is broadly true, though each text frames the struggle differently.

\emph{The Canterbury Tales:} individuals strain against the expectations of their estate. The Monk refuses the cloister and the Rule, living as a hunter and lord; the Friar turns a spiritual office into a business. Chaucer shows social pressure through the ironic gap between ideal and reality — and, exceptionally, the Knight shows that the ideal can still be lived.

\emph{A Raisin in the Sun:} the Youngers struggle against racial segregation and economic pressure. Walter fights the humiliation of servitude; Beneatha resists both racism and gender expectations by pursuing medicine and her African roots; Mama confronts housing segregation by buying in Clybourne Park. Their final refusal of Lindner's money is a collective victory of dignity over pressure.

\emph{Nineteen Eighty-Four:} Winston and Julia rebel against the most total pressure imaginable — surveillance, Newspeak, doublethink — through memory, writing and love. Their defeat in the Ministry of Love shows the darkest outcome: a society that can crush the individual completely.

\emph{Conclusion:} all three texts dramatise the individual against society, but on a scale of hope: Chaucer's pilgrims negotiate their roles with comic freedom, the Youngers win a fragile dignity, and Winston is annihilated. Together they chart the price of conformity and the courage of resistance.

\textbf{Common mistakes to avoid:} treating the three texts in isolation without comparison; vague references (name precise characters and scenes); forgetting to \emph{evaluate} — the question invites nuance, not simple agreement.
\end{corrige}

\newpage

\coursec{Summary sheet}

\begin{popfigure}
\begin{tikzpicture}[
  mindmap,
  grow cyclic,
  every node/.style={concept, text=white, font=\small},
  root concept/.append style={concept color=popink, font=\small\bfseries, text width=2.6cm, align=center},
  level 1 concept/.append style={sibling angle=51, level distance=3.6cm, font=\footnotesize, text width=2.4cm, align=center},
  level 2 concept/.append style={level distance=2.2cm, font=\scriptsize, text width=1.9cm, align=center},
]
\node[root concept]{Literary Appreciation: Three Set Texts}
  child[concept color=popblue]{ node {Chaucer: \emph{Canterbury Tales}}
    child[concept color=popblue!70]{ node {Background \& structure}}
    child[concept color=popblue!70]{ node {General Prologue: estates satire}}
  }
  child[concept color=poporange]{ node {Feudal \& Religious Orders}
    child[concept color=poporange!70]{ node {Knight, Squire, Yeoman}}
    child[concept color=poporange!70]{ node {Prioress, Monk, Friar}}
  }
  child[concept color=poppurple]{ node {Middle Class \& Rural}
    child[concept color=poppurple!70]{ node {Merchant, Clerk, Sergeant of Law}}
    child[concept color=poppurple!70]{ node {Franklin, Reeve, Miller}}
  }
  child[concept color=popgreen]{ node {Hansberry: \emph{A Raisin in the Sun}}
    child[concept color=popgreen!70]{ node {Act 1: the \$10{,}000 dream}}
    child[concept color=popgreen!70]{ node {Gender \& identity}}
  }
  child[concept color=popgold]{ node {Act 2 Scene 1}
    child[concept color=popgold!70]{ node {American Dream vs money}}
    child[concept color=popgold!70]{ node {Cultural heritage}}
  }
  child[concept color=poppink]{ node {Orwell: \emph{Nineteen Eighty-Four}}
    child[concept color=poppink!70]{ node {Part One review}}
    child[concept color=poppink!70]{ node {Context \& biography}}
  }
  child[concept color=popdark]{ node {Part Two, Ch.\ 1--6}
    child[concept color=popdark!70]{ node {Love as rebellion}}
    child[concept color=popdark!70]{ node {Julia, the room, the Brotherhood}}
  };
\end{tikzpicture}
\legende{Mind map of Chapter 1: Emphasis on Literary Appreciation --- the three set texts and the seven sections of the chapter.}
\end{popfigure}

\begin{bilan}
\textbf{1. Key definitions to master}
\begin{itemize}
  \item \cle{Estates satire}: Chaucer's technique in the \emph{General Prologue} of presenting the three medieval estates (those who fight, those who pray, those who work) to expose the gap between the ideal of each estate and its corrupt reality.
  \item \cle{Frame narrative}: an outer story (the pilgrimage to Canterbury and the tale-telling contest) that contains inner stories (the tales told by the pilgrims).
  \item \cle{Heroic couplet}: pairs of rhyming lines in iambic pentameter, Chaucer's characteristic verse form.
  \item \cle{Verbal irony}: saying the opposite of what is meant, as when Chaucer calls the pleasure-loving Monk a manly and worthy figure.
  \item \cle{Domestic drama}: a play centred on family life within the home; \emph{A Raisin in the Sun} unfolds almost entirely in the Youngers' living room.
  \item \cle{American Dream}: the belief that hard work brings prosperity, freedom and upward mobility; Hansberry tests it against the racial barriers of 1950s America.
  \item \cle{Dystopia}: an imagined nightmarish society, the opposite of a utopia; \emph{Nineteen Eighty-Four} is the classic twentieth-century example.
  \item \cle{Newspeak}, \cle{doublethink}, \cle{Thought Police}: Orwell's terms for the control of language, the holding of two contradictory beliefs at once, and the policing of thought itself.
\end{itemize}

\textbf{2. Facts and content to know}
\begin{itemize}
  \item \textbf{Chaucer} (c.\ 1343--1400): often called the ``Father of English poetry''; wrote in Middle English at a time when Latin and French dominated official writing. The \emph{General Prologue} presents about thirty pilgrims travelling to the shrine of Thomas Becket at Canterbury; the Host proposes that each pilgrim tell two tales on the way out and two on the way back --- a plan the work, left unfinished, never completes.
  \item \textbf{Feudal / military order}: the \cle{Knight} (the ideal: chivalrous, brave, experienced in many campaigns, yet modest and plainly dressed), the \cle{Squire} (his son: young, fashionable, vain, more lover than warrior), the \cle{Yeoman} (a forester and practical servant, well equipped with bow and arrows).
  \item \textbf{Corrupt religious order}: the \cle{Prioress} (false refinement, worldly sentimentality, more tender to her lapdogs than to the poor), the \cle{Monk} (loves hunting, horses and rich food, openly despises the cloister's rules), the \cle{Friar} (begs for profit, hears confessions and grants easy absolution in return for money).
  \item \textbf{Middle-class professionals}: the \cle{Merchant} (pompous and dignified in appearance, yet secretly in debt), the Oxford \cle{Clerk} (thin, poor, devoted to study and books --- one of the few idealised portraits), the \cle{Sergeant of Law} (gives the impression of being busier and wiser than he truly is).
  \item \textbf{Rural / land managers}: the \cle{Franklin} (a wealthy landowner devoted to food, drink and hospitality), the \cle{Reeve} (a shrewd, feared manor manager whom no auditor can catch out), the \cle{Miller} (coarse, strong and dishonest, famous for cheating customers with his thumb on the scale).
  \item \textbf{Hansberry} (1930--1965): the first Black woman to have a play produced on Broadway; the title comes from Langston Hughes's poem ``Harlem'', which asks what happens to ``a dream deferred''.
  \item \textbf{Act 1 Scene 1}: the Younger family, crowded into a small Chicago Southside apartment, awaits a \$10{,}000 life-insurance cheque following Big Walter's death; conflict over its use: Walter Lee's liquor-store project, Beneatha's medical studies, Mama's dream of a house.
  \item \textbf{Act 1 Scene 2 and Act 2 Scene 1}: Beneatha's search for identity (Asagai and African heritage versus George Murchison's assimilation); Ruth's pregnancy; Mama buys a house in the white neighbourhood of Clybourne Park and entrusts the remaining money to Walter.
  \item \textbf{Orwell} (Eric Arthur Blair, 1903--1950): a democratic socialist whose experiences of Stalinism and fascism led him to write the novel as a satire and warning against totalitarianism.
  \item \textbf{Part One}: Winston Smith's life under Big Brother --- telescreens, the Ministries (Truth, Peace, Love, Plenty), the forbidden diary, the Two Minutes Hate, and the proles as the only hope of change.
  \item \textbf{Part Two, Chapters 1--6}: Julia's note (``I love you''); their secret affair in the countryside and in the rented room above Charrington's shop; love experienced as a political act of rebellion; the approach to O'Brien and the supposed Brotherhood.
\end{itemize}

\textbf{3. Methods for the GCE essay}
\begin{itemize}
  \item Use the \cle{PEEL} paragraph: Point, Evidence (a short quotation or a precise reference), Explanation, Link back to the question.
  \item Always connect a \textbf{character portrait} to a \textbf{theme} and to the \textbf{author's technique} (irony, symbolism, stage directions, narrative voice).
  \item For Chaucer, contrast the \emph{ideal} figures (Knight, Clerk) with the \emph{corrupt} ones (Monk, Friar, Miller) to bring out the satirical purpose.
  \item For drama, comment on \cle{stage directions} and \cle{symbols}: Mama's plant, Beneatha's hair, the insurance cheque, the light from the single window.
  \item For Orwell, link plot events to \textbf{themes}: surveillance, truth, the control of language, love as rebellion.
\end{itemize}

\textbf{4. Common mistakes to avoid}
\begin{itemize}
  \item Retelling the plot instead of \emph{analysing} --- the examiner rewards argument, not summary.
  \item Confusing characters: the \textbf{Squire} is the Knight's son; the \textbf{Yeoman} is a servant, not a landowner; the \textbf{Reeve} manages a manor, while the \textbf{Miller} grinds corn.
  \item Treating \emph{Nineteen Eighty-Four} as a mere prediction of the future: it is a \emph{satire and warning} rooted in the totalitarianisms of Orwell's own time.
  \item Ignoring context: medieval feudal society, 1950s segregated Chicago, post-war totalitarian Europe.
  \item Quoting long passages inaccurately --- prefer short, certain quotations or accurate paraphrase.
  \item Mixing up Beneatha's two suitors: \textbf{Asagai} (African pride) versus \textbf{George Murchison} (assimilation).
\end{itemize}

\textbf{5. GCE tips}
\begin{itemize}
  \item Learn two or three short quotations per text (for example Hughes's ``dream deferred''; Orwell's ``Big Brother is watching you'').
  \item Prepare one essay plan per major theme: satire in the \emph{General Prologue}; dreams and money in \emph{A Raisin in the Sun}; rebellion in \emph{Nineteen Eighty-Four}.
  \item Structure every answer: introduction (text, author, context, thesis), three or four PEEL paragraphs, conclusion.
  \item Manage your time: about 45--50 minutes per essay, always keeping three minutes to proofread.
\end{itemize}
\end{bilan}

\findecours

\end{document}
